\documentclass[10pt, letterpaper]{article}

% Inhaltsverzeichnis für Pakettypen (nur für Übersicht im Header, wird nicht im Dokument angezeigt)
% 1. Seitenlayout und Ränder
% 2. Sprache und Zeichensatz
% 3. Mathematik und Theorem-Umgebungen
% 4. Eigene Makros
% 5. Diagramme und Grafiken
% 6. Tabellen und Aufzählungen
% 7. Inhaltsverzeichnis
% 8. Abschnittsüberschriften
% 9. Abstrakt-Umgebung
% 10. Todos/Notizen
% 11. Rahmen/Box-Umgebungen
% 12. Python-Integration
% 13. Literaturverwaltung
% 14. Hyperlinks
% 15. Absatzeinstellungen
% 16. Umgebungen
% 17  Graphik
% 18  Extra
% 00. Titel und Autor

\usepackage{slashed}

% --- 1. Seitenlayout und Ränder ---
\usepackage[margin=3cm]{geometry}

% --- 2. Sprache und Zeichensatz ---
\usepackage[english]{babel}
\usepackage[T1]{fontenc}
\usepackage[utf8]{inputenc}

% --- 3. Mathematik und Theorem-Umgebungen ---
\usepackage{amsmath, amssymb, amsthm}
\usepackage{mathrsfs}
\DeclareMathOperator{\WF}{WF}

% --- 4. Eigene Makros ---
\usepackage{xcolor}
\newcommand{\SKP}{\langle\cdot,\cdot\rangle}
\newcommand{\R}{\mathbb{R}}
\newcommand{\N}{\mathbb{N}}
\newcommand{\Q}{\mathbb{Q}}
\newcommand{\Z}{\mathbb{Z}}
\newcommand{\C}{\mathbb{C}}
\newcommand{\entwurf}[1]{\textcolor{red}{#1}}

% --- 5. Diagramme und Grafiken ---
\usepackage{graphicx}
\graphicspath{{../../Library/Images/}}
\usepackage[export]{adjustbox}
\usepackage{tikz}
\usetikzlibrary{decorations.pathreplacing, arrows.meta, positioning}
\usepackage{tikz-cd}

% --- 6. Tabellen und Aufzählungen ---
\usepackage{enumitem}
\setlist[itemize]{left=0.5cm}

\newenvironment{romanenum}[1][]
  {%
    \ifx&#1&
    \else
      \textbf{#1}\quad
    \fi
    \begin{enumerate}[label=\roman*)]
  }
  {%
    \end{enumerate}%
  }

% --- 7. Inhaltsverzeichnis ---
\usepackage{tocloft}
\renewcommand{\cftsecfont}{\footnotesize}
\renewcommand{\cftsubsecfont}{\footnotesize}
\renewcommand{\cftsubsubsecfont}{\footnotesize}
\renewcommand{\cftsecpagefont}{\footnotesize}
\renewcommand{\cftsubsecpagefont}{\footnotesize}
\renewcommand{\cftsubsubsecpagefont}{\footnotesize}
\usepackage{etoc}

% --- 8. Abschnittsüberschriften ---
\usepackage{titlesec}
\titleformat{\section}{\normalfont\large\bfseries}{\thesection}{1em}{}
\titleformat{\subsection}{\normalfont\normalsize\bfseries}{\thesubsection}{0.5em}{}
\titleformat{\subsubsection}{\normalfont\normalsize\bfseries}{\thesubsubsection}{0.5em}{}
\setcounter{secnumdepth}{4}

% --- 9. Abstrakt-Umgebung ---
\usepackage{changepage}
\renewenvironment{abstract}
  {
    \begin{adjustwidth}{1.5cm}{1.5cm}
    \small
    \textsc{Abstract. –}%
  }
  {
    \end{adjustwidth}
  }

% --- 10. Todos/Notizen ---
\usepackage{todonotes}
% Hellblau definieren
\definecolor{hellblau}{rgb}{0.3,0.6,1}
\newenvironment{note}{\par\begingroup\color{hellblau}\itshape}{\par\endgroup}

% --- 11. Rahmen/Box-Umgebungen ---
\usepackage{mdframed}
\usepackage{tcolorbox}
\colorlet{shadecolor}{gray!25}

\newenvironment{customTheorem}
  {\vspace{10pt}%
   \begin{mdframed}[
     backgroundcolor=gray!20,
     linewidth=0pt,
     innertopmargin=10pt,
     innerbottommargin=10pt,
     skipabove=\dimexpr\topsep+\ht\strutbox\relax,
     skipbelow=\topsep,
   ]}
  {\end{mdframed}
   \vspace{10pt}%
  }

% --- 12. Python-Integration ---
% (Deaktiviert in dieser Version, aktiviere bei Bedarf)
% \usepackage{pythontex}
% \usepackage[makestderr]{pythontex}

% --- 13. Literaturverwaltung ---
\usepackage{csquotes}
\usepackage[backend=biber, style=alphabetic, citestyle=alphabetic]{biblatex}
\addbibresource{bibliography.bib}

% --- 14. Hyperlinks ---
\usepackage{hyperref}
\hypersetup{
  colorlinks   = true,
  urlcolor     = blue,
  linkcolor    = blue,
  citecolor    = blue,
  frenchlinks  = true
}

% --- 15. Absatzeinstellungen ---
\usepackage[parfill]{parskip}
\sloppy

% --- 16. Umgebungen ---
\usepackage{thmtools}

\newcommand{\CustomHeading}[3]{%
  \par\medskip\noindent%
  \textbf{#1 #2} \textnormal{(#3)}.\enskip%
}

\newenvironment{DEF}[2]{\begin{unitbox}\CustomHeading{Definition}{#1}{#2}}{\end{unitbox}}
\newenvironment{PROP}[2]{\begin{unitbox}\CustomHeading{Proposition}{#1}{#2}}{\end{unitbox}}
\newenvironment{THEO}[2]{\begin{unitbox}\CustomHeading{Theorem}{#1}{#2}}{\end{unitbox}}
\newenvironment{LEM}[2]{\begin{unitbox}\CustomHeading{Lemma}{#1}{#2}}{\end{unitbox}}
\newenvironment{KORO}[2]{\begin{unitbox}\CustomHeading{Corollar}{#1}{#2}}{\end{unitbox}}
\newenvironment{REM}[2]{\begin{unitbox}\CustomHeading{Remark}{#1}{#2}}{\end{unitbox}}
\newenvironment{EXA}[2]{\begin{unitbox}\CustomHeading{Example}{#1}{#2}}{\end{unitbox}}
\newenvironment{STUD}[2]{\begin{unitbox}\CustomHeading{Study}{#1}{#2}}{\end{unitbox}}
\newenvironment{CONC}[2]{\begin{unitbox}\CustomHeading{Concept}{#1}{#2}}{\end{unitbox}}
\newenvironment{OTH}[2]{\begin{unitbox}\CustomHeading{Other}{#1}{#2}}{\end{unitbox}}
\newenvironment{EXE}[2]{\begin{unitbox}\CustomHeading{Exercise}{#1}{#2}}{\end{unitbox}}
\newenvironment{MOT}[2]{\begin{unitbox}\CustomHeading{Motivation}{#1}{#2}}{\end{unitbox}}
\newenvironment{PROOF}[2]{\begin{unitbox}\CustomHeading{Proof}{#1}{#2}}{\end{unitbox}}

% --- Unit Umgebung für Source-Inhalte ---
\usepackage{mdframed}
\newmdenv[
  linewidth=1pt,
  topline=false,
  bottomline=false,
  rightline=false,
  leftmargin=0cm,
  rightmargin=0cm,
  skipabove=10pt,
  skipbelow=10pt,
  innertopmargin=0.5\baselineskip,
  innerbottommargin=0.5\baselineskip,
  backgroundcolor=gray!10,
  linecolor=gray
]{unitbox}

\newenvironment{unit}[1]
  {\begin{unitbox}\textbf{Unit #1}\par\smallskip}
  {\end{unitbox}}

% --- 17. ... ---

% --- 18. Extras ---
\usepackage{stmaryrd}
\usepackage{bbold}  % falls du \mathbb{1} nutzen willst

% --- 00. Titel und Autor ---
\title{Mein Titel}
\author{Tim Jaschik}
\date{\today}

\begin{document}

\maketitle
\rule{\textwidth}{0.5pt}
\begin{abstract}
Kurze Beschreibung …
\end{abstract}
\rule{\textwidth}{0.5pt}
\vspace{0.5cm}

\tableofcontents

\pagebreak

\section*{Chapter 1 Matrix Lie Groups}
\subsection*{Definitions}
A Lie group is, roughly speaking, a continuous group, that is, a group described by several real parameters. In this book, we consider matrix Lie groups, which are Lie groups realized as groups of matrices. As an example, consider the set of all $2 \times 2$ real matrices with determinant 1 , customarily denoted $\operatorname{SL}(2 ; \mathbb{R})$. Since the determinant of a product is the product of the determinants, this set forms a group under the operation of matrix multiplication. If we think of the set of all $2 \times 2$ matrices, with entries $a, b, c, d$, as $\mathbb{R}^{4}$, then $\operatorname{SL}(2 ; \mathbb{R})$ is the set of points in $\mathbb{R}^{4}$ for which the smooth function $a d-b c$ has the value 1 .

Suppose $f$ is a smooth function on $\mathbb{R}^{k}$ and we consider the set $E$ where $f(\mathbf{x})$ equals some constant value $c$. If, at each point $\mathbf{x}_{0}$ in $E$, at least one of the partial derivatives of $f$ is nonzero, then the implicit function theorem tells us that we can solve the equation $f(\mathbf{x})=c$ near $\mathbf{x}_{0}$ for one of the variables as a function of the other $k-1$ variables. Thus, $E$ is a smooth "surface" (or embedded submanifold) in $\mathbb{R}^{k}$ of dimension $k-1$. In the case of $\operatorname{SL}(2 ; \mathbb{R})$ inside $\mathbb{R}^{4}$, we note that the partial derivatives of $a d-b c$ with respect to $a, b, c$, and $d$ are $d,-c,-b$, and $a$, respectively. Thus, at each point where $a d-b c=1$, at least one of these partial derivatives is nonzero, and we conclude that $\mathrm{SL}(2 ; \mathbb{R})$ is a smooth surface of dimension 3 . Thus, $\mathrm{SL}(2 ; \mathbb{R})$ is a Lie group of dimension 3.

For other groups of matrices (such as the ones we will encounter later in this section), one could use a similar approach. The analysis is, however, more complicated because most of the groups are defined by setting several different smooth functions equal to constants. One therefore has to check that these functions are "independent" in the sense of the implicit function theorem, which means that their gradient vectors have to be linearly independent at each point in the group.

We will use an alternative approach that makes all such analysis unnecessary. We consider groups $G$ of matrices that are closed in the sense of Definition 1.4. To each such $G$, we will associate in Chapter 3 a "Lie algebra" $\mathfrak{g}$, which is a real vector\\
space. A general result (Corollary 3.45) will then show that $G$ is a smooth manifold whose dimension is equal the dimension of $\mathfrak{g}$ as a vector space.

This chapter makes use of various standard results from linear algebra, which are summarized in Appendix A.

Definition 1.1. The general linear group over the real numbers, denoted $\mathrm{GL}(n ; \mathbb{R})$, is the group of all $n \times n$ invertible matrices with real entries. The general linear group over the complex numbers, denoted $\operatorname{GL}(n ; \mathbb{C})$, is the group of all $n \times n$ invertible matrices with complex entries.

Definition 1.2. Let $M_{n}(\mathbb{C})$ denote the space of all $n \times n$ matrices with complex entries.

We may identify $M_{n}(\mathbb{C})$ with $\mathbb{C}^{n^{2}}$ and use the standard notion of convergence in $\mathbb{C}^{n^{2}}$. Explicitly, this means the following.

Definition 1.3. Let $A_{m}$ be a sequence of complex matrices in $M_{n}(\mathbb{C})$. We say that $A_{m}$ converges to a matrix $A$ if each entry of $A_{m}$ converges (as $m \rightarrow \infty$ ) to the corresponding entry of $A$ (i.e., if $\left(A_{m}\right)_{j k}$ converges to $A_{j k}$ for all $1 \leq j, k \leq n$ ).

We now consider subgroups of $\mathrm{GL}(n ; \mathbb{C})$, that is, subsets $G$ of $\mathrm{GL}(n ; \mathbb{C})$ such that the identity matrix is in $G$ and such that for all $A$ and $B$ in $G$, the matrices $A B$ and $A^{-1}$ are also in $G$.

Definition 1.4. A matrix Lie group is a subgroup $G$ of $\mathrm{GL}(n ; \mathbb{C})$ with the following property: If $A_{m}$ is any sequence of matrices in $G$, and $A_{m}$ converges to some matrix $A$, then either $A$ is in $G$ or $A$ is not invertible.

The condition on $G$ amounts to saying that $G$ is a closed subset of $\operatorname{GL}(n ; \mathbb{C})$. (This does not necessarily mean that $G$ is closed in $M_{n}(\mathbb{C})$.) Thus, Definition 1.4 is equivalent to saying that a matrix Lie group is a closed subgroup of $\operatorname{GL}(n ; \mathbb{C})$.

The condition that $G$ be a closed subgroup, as opposed to merely a subgroup, should be regarded as a technicality, in that most of the interesting subgroups of $\mathrm{GL}(n ; \mathbb{C})$ have this property. Most of the matrix Lie groups $G$ we will consider have the stronger property that if $A_{m}$ is any sequence of matrices in $G$, and $A_{m}$ converges to some matrix $A$, then $A \in G$ (i.e., that $G$ is closed in $M_{n}(\mathbb{C})$ ).

An example of a subgroup of $\mathrm{GL}(n ; \mathbb{C})$ which is not closed (and hence is not a matrix Lie group) is the set of all $n \times n$ invertible matrices with rational entries. This set is, in fact, a subgroup of $\mathrm{GL}(n ; \mathbb{C})$, but not a closed subgroup. That is, one can (easily) have a sequence of invertible matrices with rational entries converging to an invertible matrix with some irrational entries. (In fact, every real invertible matrix is the limit of some sequence of invertible matrices with rational entries.)

Another example of a group of matrices which is not a matrix Lie group is the following subgroup of $\operatorname{GL}(2 ; \mathbb{C})$. Let $a$ be an irrational real number and let

$$
G=\left\{\left.\left(\begin{array}{cc}
e^{i t} & 0 \\
0 & e^{i t a}
\end{array}\right) \right\rvert\, t \in \mathbb{R}\right\} .
$$

\includegraphics[scale=0.2, center]{2025_08_28_fc86fd4ade69f41b5842g-018}

Fig. 1.1 A small portion of the group $G$ inside $\bar{G}$ (left) and a larger portion (right)

Clearly, $G$ is a subgroup of $\mathrm{GL}(2 ; \mathbb{C})$. According to Exercise 10, the closure of $G$ is the group

$$
\bar{G}=\left\{\left.\left(\begin{array}{cc}
e^{i \theta} & 0 \\
0 & e^{i \phi}
\end{array}\right) \right\rvert\, \theta, \phi \in \mathbb{R}\right\} .
$$

The group $G$ inside $\bar{G}$ is known as an "irrational line in a torus"; see Figure 1.1.

\subsection*{Examples}
Mastering the subject of Lie groups involves not only learning the general theory but also familiarizing oneself with examples. In this section, we introduce some of the most important examples of (matrix) Lie groups. Among these are the classical groups, consisting of the general and special linear groups, the unitary and orthogonal groups, and the symplectic groups. The classical groups, and their associated Lie algebras, will be key examples in Parts II and III of the book.

\subsubsection*{General and Special Linear Groups}
The general linear groups (over $\mathbb{R}$ or $\mathbb{C}$ ) are themselves matrix Lie groups. Of course, $\mathrm{GL}(n ; \mathbb{C})$ is a subgroup of itself. Furthermore, if $A_{m}$ is a sequence of matrices in $\mathrm{GL}(n ; \mathbb{C})$ and $A_{m}$ converges to $A$, then by the definition of $\mathrm{GL}(n ; \mathbb{C})$, either $A$ is in $\mathrm{GL}(n ; \mathbb{C})$, or $A$ is not invertible.

Moreover, $\mathrm{GL}(n ; \mathbb{R})$ is a subgroup of $\mathrm{GL}(n ; \mathbb{C})$, and if $A_{m} \in \mathrm{GL}(n ; \mathbb{R})$ and $A_{m}$ converges to $A$, then the entries of $A$ are real. Thus, either $A$ is not invertible or $A \in \mathrm{GL}(n ; \mathbb{R})$.

The special linear group (over $\mathbb{R}$ or $\mathbb{C}$ ) is the group of $n \times n$ invertible matrices (with real or complex entries) having determinant one. Both of these are subgroups of $\mathrm{GL}(n ; \mathbb{C})$. Furthermore, if $A_{n}$ is a sequence of matrices with determinant one and $A_{n}$ converges to $A$, then $A$ also has determinant one, because the determinant is a continuous function. Thus, $\mathrm{SL}(n ; \mathbb{R})$ and $\mathrm{SL}(n ; \mathbb{C})$ are matrix Lie groups.

\subsubsection*{Unitary and Orthogonal Groups}
An $n \times n$ complex matrix $A$ is said to be unitary if the column vectors of $A$ are orthonormal, that is, if

$$
\sum_{l=1}^{n} \overline{A_{l j}} A_{l k}=\delta_{j k}
$$

We may rewrite (1.2) as

$$
\sum_{l=1}^{n}\left(A^{*}\right)_{j l} A_{l k}=\delta_{j k}
$$

where $\delta_{j k}$ is the Kronecker delta, equal to 1 if $j=k$ and equal to zero if $j \neq k$. Here $A^{*}$ is the adjoint of $A$, defined by

$$
\left(A^{*}\right)_{j k}=\overline{A_{k j}}
$$

Equation (1.3) says that $A^{*} A=I$; thus, we see that $A$ is unitary if and only if $A^{*}=A^{-1}$. In particular, every unitary matrix is invertible.

The adjoint operation on matrices satisfies $(A B)^{*}=B^{*} A^{*}$. From this, we can see that if $A$ and $B$ are unitary, then

$$
(A B)^{*}(A B)=B^{*} A^{*} A B=B^{-1} A^{-1} A B=I
$$

showing that $A B$ is also unitary. Furthermore, since $\left(A A^{-1}\right)^{*}=I^{*}=I$, we see that $\left(A^{-1}\right)^{*} A^{*}=I$, which shows that $\left(A^{-1}\right)^{*}=\left(A^{*}\right)^{-1}$. Thus, if $A$ is unitary, we have

$$
\left(A^{-1}\right)^{*} A^{-1}=\left(A^{*}\right)^{-1} A^{-1}=\left(A A^{*}\right)^{-1}=I,
$$

showing that $A^{-1}$ is again unitary.\\
Thus, the collection of unitary matrices is a subgroup of $\operatorname{GL}(n ; \mathbb{C})$. We call this group the unitary group and we denote it by $\mathrm{U}(n)$. We may also define the special\\
unitary group $\mathrm{SU}(n)$, the subgroup of $\mathrm{U}(n)$ consisting of unitary matrices with determinant 1 . It is easy to check that both $\mathrm{U}(n)$ and $\mathrm{SU}(n)$ are closed subgroups of $\mathrm{GL}(n ; \mathbb{C})$ and thus matrix Lie groups.

Meanwhile, let $\langle\cdot, \cdot\rangle$ denote the standard inner product on $\mathbb{C}^{n}$, given by

$$
\langle x, y\rangle=\sum_{j} \overline{x_{j}} y_{j}
$$

(Note that we put the conjugate on the first factor in the inner product.) By Proposition A.8, we have

$$
\langle x, A y\rangle=\left\langle A^{*} x, y\right\rangle
$$

for all $x, y \in \mathbb{C}^{n}$. Thus,

$$
\langle A x, A y\rangle=\left\langle A^{*} A x, y\right\rangle,
$$

from which we can see that if $A$ is unitary, then $A$ preserves the inner product on $\mathbb{C}^{n}$, that is,

$$
\langle A x, A y\rangle=\langle x, y\rangle
$$

for all $x$ and $y$. Conversely, if $A$ preserves the inner product, we must have $\left\langle A^{*} A x, y\right\rangle=\langle x, y\rangle$ for all $x, y$. It is not hard to see that this condition holds only if $A^{*} A=I$. Thus, an equivalent characterization of unitarity is that $A$ is unitary if and only if $A$ preserves the standard inner product on $\mathbb{C}^{n}$.

Finally, for any matrix $A$, we have that $\operatorname{det} A^{*}=\overline{\operatorname{det} A}$. Thus, if $A$ is unitary, we have

$$
\operatorname{det}\left(A^{*} A\right)=|\operatorname{det} A|^{2}=\operatorname{det} I=1
$$

Hence, for all unitary matrices $A$, we have $|\operatorname{det} A|=1$.\\
In a similar fashion, an $n \times n$ real matrix $A$ is said to be orthogonal if the column vectors of $A$ are orthonormal. As in the unitary case, we may give equivalent versions of this condition. The only difference is that if $A$ is real, $A^{*}$ is the same as the transpose $A^{t r}$ of $A$, given by

$$
\left(A^{t r}\right)_{j k}=A_{k j} .
$$

Thus, $A$ is orthogonal if and only if $A^{t r}=A^{-1}$, and this holds if and only if $A$ preserves the inner product on $\mathbb{R}^{n}$. Since $\operatorname{det}\left(A^{t r}\right)=\operatorname{det} A$, if $A$ is orthogonal, we have

$$
\operatorname{det}\left(A^{t r} A\right)=\operatorname{det}(A)^{2}=\operatorname{det}(I)=1,
$$

so that $\operatorname{det}(A)= \pm 1$. The collection of all orthogonal matrices forms a closed subgroup of $\mathrm{GL}(n ; \mathbb{C})$, which we call the orthogonal group and denote by $\mathrm{O}(n)$. The set of $n \times n$ orthogonal matrices with determinant one is the special orthogonal group, denoted $\mathrm{SO}(n)$. Geometrically, elements of $\mathrm{SO}(n)$ are rotations, while the elements of $\mathrm{O}(n)$ are either rotations or combinations of rotations and reflections.

Consider now the bilinear form ( $\cdot, \cdot$ ) on $\mathbb{C}^{n}$ defined by

$$
(x, y)=\sum_{j} x_{j} y_{j} .
$$

This form is not an inner product (Sect. A.6) because, for example, it is symmetric rather than conjugate-symmetric. The set of all $n \times n$ complex matrices $A$ which preserve this form (i.e., such that $(A x, A y)=(x, y)$ for all $x, y \in \mathbb{C}^{n}$ ) is the complex orthogonal group $\mathrm{O}(n ; \mathbb{C})$, and it is a subgroup of $\mathrm{GL}(n ; \mathbb{C})$. Since there are no conjugates in the definition of the form ( $\cdot, \cdot$ ), we have

$$
(x, A y)=\left(A^{t r} x, y\right),
$$

for all $x, y \in \mathbb{C}^{n}$, where on the right-hand side of the above relation, we have $A^{t r}$ rather than $A^{*}$. Repeating the arguments for the case of $\mathrm{O}(n)$, but now allowing complex entries in our matrices, we find that an $n \times n$ complex matrix $A$ is in $\mathrm{O}(n ; \mathbb{C})$ if and only if $A^{t r} A=I$, that $\mathrm{O}(n ; \mathbb{C})$ is a matrix Lie group, and that $\operatorname{det} A= \pm 1$ for all $A$ in $\mathrm{O}(n ; \mathbb{C})$. Note that $\mathrm{O}(n ; \mathbb{C})$ is not the same as the unitary group $\mathrm{U}(n)$. The group $\mathrm{SO}(n ; \mathbb{C})$ is defined to be the set of all $A$ in $\mathrm{O}(n ; \mathbb{C})$ with $\operatorname{det} A=1$ and it is also a matrix Lie group.

\subsubsection*{Generalized Orthogonal and Lorentz Groups}
Let $n$ and $k$ be positive integers, and consider $\mathbb{R}^{n+k}$. Define a symmetric bilinear form $[\cdot, \cdot]_{n, k}$ on $\mathbb{R}^{n+k}$ by the formula

$$
[x, y]_{n, k}=x_{1} y_{1}+\cdots+x_{n} y_{n}-x_{n+1} y_{n+1}-\cdots-x_{n+k} y_{n+k}
$$

The set of $(n+k) \times(n+k)$ real matrices $A$ which preserve this form (i.e., such that $[A x, A y]_{n, k}=[x, y]_{n, k}$ for all $x, y \in \mathbb{R}^{n+k}$ ) is the generalized orthogonal group $\mathrm{O}(n ; k)$. It is a subgroup of $\mathrm{GL}(n+k ; \mathbb{R})$ and a matrix Lie group (Exercise 1). Of particular interest in physics is the Lorentz group $\mathrm{O}(3 ; 1)$. We also define $\mathrm{SO}(n ; k)$ to be the subgroup of $\mathrm{O}(n ; k)$ consisting of elements with determinant 1 .

If $A$ is an $(n+k) \times(n+k)$ real matrix, let $A^{(j)}$ denote the $j$ th column vector of $A$, that is,

$$
A^{(j)}=\left(\begin{array}{c}
A_{1, j} \\
\vdots \\
A_{n+k, j}
\end{array}\right) .
$$

Note that $A^{(j)}$ is equal to $A e_{j}$, that is, the result of applying $A$ to the $j$ th standard basis element $e_{j}$. Then $A$ will belong to $\mathrm{O}(n ; k)$ if and only $\left[A e_{j}, A e_{l}\right]$ for all $1 \leq j, l \leq n+k$. Explicitly, this means that $A \in \mathrm{O}(n ; k)$ if and only if the following conditions are satisfied:

$$
\begin{array}{ll}
{\left[A^{(j)}, A^{(l)}\right]_{n, k}=0} & j \neq l, \\
{\left[A^{(j)}, A^{(j)}\right]_{n, k}=1} & 1 \leq j \leq n, \\
{\left[A^{(j)}, A^{(j)}\right]_{n, k}=-1} & n+1 \leq j \leq n+k .
\end{array}
$$

Let $g$ denote the $(n+k) \times(n+k)$ diagonal matrix with ones in the first $n$ diagonal entries and minus ones in the last $k$ diagonal entries:

$$
g=\left(\begin{array}{llllll}
1 & & & & & \\
& \ddots & & & & \\
& & 1 & & & \\
& & & -1 & & \\
& & & & \ddots & \\
& & & & & -1
\end{array}\right)
$$

Then $A$ is in $\mathrm{O}(n ; k)$ if and only if $A^{t r} g A=g$ (Exercise 1). Taking the determinant of this equation gives $(\operatorname{det} A)^{2} \operatorname{det} g=\operatorname{det} g$, or $(\operatorname{det} A)^{2}=1$. Thus, for any $A$ in $\mathrm{O}(n ; k), \operatorname{det} A= \pm 1$.

\subsubsection*{Symplectic Groups}
Consider the skew-symmetric bilinear form $B$ on $\mathbb{R}^{2 n}$ defined as follows:

$$
\omega(x, y)=\sum_{j=1}^{n}\left(x_{j} y_{n+j}-x_{n+j} y_{j}\right) .
$$

The set of all $2 n \times 2 n$ matrices $A$ which preserve $\omega$ (i.e., such that $\omega(A x, A y)= \omega(x, y)$ for all $x, y \in \mathbb{R}^{2 n}$ ) is the real symplectic group $\operatorname{Sp}(n ; \mathbb{R})$, and it is a closed subgroup of $\operatorname{GL}(2 n ; \mathbb{R})$. (Some authors refer to the group we have just defined as $\operatorname{Sp}(2 n ; \mathbb{R})$ rather than $\operatorname{Sp}(n ; \mathbb{R})$.) If $\Omega$ is the $2 n \times 2 n$ matrix

$$
\Omega=\left(\begin{array}{cc}
0 & I \\
-I & 0
\end{array}\right),
$$

then

$$
\omega(x, y)=\langle x, \Omega y\rangle .
$$

From this, it is not hard to show that a $2 n \times 2 n$ real matrix $A$ belongs to $\operatorname{Sp}(n ; \mathbb{R})$ if and only if

$$
-\Omega A^{t r} \Omega=A^{-1} .
$$

(See Exercise 2.) Taking the determinant of this identity gives $\operatorname{det} A=(\operatorname{det} A)^{-1}$, i.e., $(\operatorname{det} A)^{2}=1$. This shows that $\operatorname{det} A= \pm 1$, for all $A \in \operatorname{Sp}(n ; \mathbb{R})$. In fact, $\operatorname{det} A=1$ for all $A \in \operatorname{Sp}(n ; \mathbb{R})$, although this is not obvious.

One can define a bilinear form $\omega$ on $\mathbb{C}^{2 n}$ by the same formula as in (1.7) (with no conjugates). Over $\mathbb{C}$, we have the relation

$$
\omega(z, w)=(z, \Omega w),
$$

where ( $\cdot, \cdot$ ) is the complex bilinear form in (1.4). The set of $2 n \times 2 n$ complex matrices which preserve this form is the complex symplectic group $\operatorname{Sp}(n ; \mathbb{C})$. A $2 n \times 2 n$ complex matrix $A$ is in $\operatorname{Sp}(n ; \mathbb{C})$ if and only if (1.9) holds. (Note: This condition involves $A^{t r}$, not $A^{*}$.) Again, we can easily show that each $A \in \operatorname{Sp}(n ; \mathbb{C})$ satisfies $\operatorname{det} A= \pm 1$ and, again, it is actually the case that $\operatorname{det} A=1$. Finally, we have the compact symplectic group $\mathrm{Sp}(n)$ defined as

$$
\operatorname{Sp}(n)=\operatorname{Sp}(n ; \mathbb{C}) \cap \mathrm{U}(2 n) .
$$

That is to say, $\operatorname{Sp}(n)$ is the group of $2 n \times 2 n$ matrices that preserve both the inner product and the bilinear form $\omega$. For more information about $\operatorname{Sp}(n)$, see Sect. 1.2.8.

\subsubsection*{The Euclidean and Poincaré Groups}
The Euclidean group $\mathrm{E}(n)$ is the group of all transformations of $\mathbb{R}^{n}$ that can be expressed as a composition of a translation and an orthogonal linear transformation. We write elements of $\mathrm{E}(n)$ as pairs $\{x, R\}$ with $x \in \mathbb{R}^{n}$ and $R \in \mathrm{O}(n)$, and we let $\{x, R\}$ act on $\mathbb{R}^{n}$ by the formula

$$
\{x, R\} y=R y+x .
$$

Since

$$
\left\{x_{1}, R_{1}\right\}\left\{x_{2}, R_{2}\right\} y=R_{1}\left(R_{2} y+x_{2}\right)+x_{1}=R_{1} R_{2} y+\left(x_{1}+R_{1} x_{2}\right),
$$

the product operation for $\mathrm{E}(n)$ is the following:

$$
\left\{x_{1}, R_{1}\right\}\left\{x_{2}, R_{2}\right\}=\left\{x_{1}+R_{1} x_{2}, R_{1} R_{2}\right\} .
$$

The inverse of an element of $\mathrm{E}(n)$ is given by

$$
\{x, R\}^{-1}=\left\{-R^{-1} x, R^{-1}\right\}
$$

The group $\mathrm{E}(n)$ is not a subgroup of $\mathrm{GL}(n ; \mathbb{R})$, since translations are not linear maps. However, $\mathrm{E}(n)$ is isomorphic to the (closed) subgroup of $\mathrm{GL}(n+1 ; \mathbb{R})$ consisting of matrices of the form

$$
\left(\begin{array}{ccc} 
& & x_{1} \\
& R & \vdots \\
& & x_{n} \\
0 & \cdots & 0
\end{array}\right)
$$

with $R \in \mathrm{O}(n)$. (The reader may easily verify that matrices of the form (1.11) multiply according to the formula in (1.10).)

We similarly define the Poincaré group $\mathrm{P}(n ; 1)$ (also known as the inhomogeneous Lorentz group) to be the group of all transformations of $\mathbb{R}^{n+1}$ of the form

$$
T=T_{x} A
$$

with $x \in \mathbb{R}^{n+1}$ and $A \in \mathrm{O}(n ; 1)$. This group is isomorphic to the group of $(n+2) \times$ ( $n+2$ ) matrices of the form

$$
\left(\begin{array}{ccc} 
& & x_{1} \\
& A & \vdots \\
& & x_{n+1} \\
0 & \cdots & 0
\end{array}\right)
$$

with $A \in \mathrm{O}(n ; 1)$.

\subsubsection*{The Heisenberg Group}
The set of all $3 \times 3$ real matrices $A$ of the form

$$
A=\left(\begin{array}{lll}
1 & a & b \\
0 & 1 & c \\
0 & 0 & 1
\end{array}\right)
$$

where $a, b$, and $c$ are arbitrary real numbers, is the Heisenberg group. It is easy to check that the product of two matrices of the form (1.13) is again of that form, and,\\
clearly, the identity matrix is of the form (1.13). Furthermore, direct computation shows that if $A$ is as in (1.13), then

$$
A^{-1}=\left(\begin{array}{ccc}
1 & -a & a c-b \\
0 & 1 & -c \\
0 & 0 & 1
\end{array}\right)
$$

Thus, $H$ is a subgroup of $\operatorname{GL}(3 ; \mathbb{R})$. The Heisenberg group is a model for the Heisenberg-Weyl commutation relations in physics and also serves as a illuminating example for the Baker-Campbell-Hausdorff formula (Sect. 5.2). See also Exercise 8.

\subsubsection*{The Groups $\mathbb{R}^{*}, \mathbb{C}^{*}, S^{1}, \mathbb{R}$, and $\mathbb{R}^{n}$}
Several important groups which are not defined as groups of matrices can be thought of as such. The group $\mathbb{R}^{*}$ of non-zero real numbers under multiplication is isomorphic to $\mathrm{GL}(1 ; \mathbb{R})$. Similarly, the group $\mathbb{C}^{*}$ of nonzero complex numbers under multiplication is isomorphic to $\mathrm{GL}(1 ; \mathbb{C})$ and the group $S^{1}$ of complex numbers with absolute value one is isomorphic to $\mathrm{U}(1)$.

The group $\mathbb{R}$ under addition is isomorphic to $\mathrm{GL}(1 ; \mathbb{R})^{+}$( $1 \times 1$ real matrices with positive determinant) via the map $x \rightarrow\left[e^{x}\right]$. The group $\mathbb{R}^{n}$ (with vector addition) is isomorphic to the group of diagonal real matrices with positive diagonal entries, via the map

$$
\left(x_{1}, \ldots, x_{n}\right) \rightarrow\left(\begin{array}{ccc}
e^{x_{1}} & & 0 \\
& \ddots & \\
0 & & e^{x_{n}}
\end{array}\right)
$$

\subsubsection*{The Compact Symplectic Group}
Of the groups introduced in the preceding subsections, the compact symplectic group $\operatorname{Sp}(n):=\operatorname{Sp}(n ; \mathbb{C}) \cap \mathrm{U}(2 n)$ is the most mysterious. In this section, we attempt to understand the structure of $\operatorname{Sp}(n)$ and to show that it can be understood as being the "unitary group over the quaternions."

Since the definition of $\operatorname{Sp}(n)$ involves unitarity, it is convenient to express the bilinear form $\omega$ on $\mathbb{C}^{2 n}$ in terms of the inner product $\langle\cdot, \cdot\rangle$, rather than in terms of the bilinear form ( $\cdot, \cdot$ ), as we did in Sect. 1.2.4. To this end, define a conjugate-linear map $J: \mathbb{C}^{2 n} \rightarrow \mathbb{C}^{2 n}$ by

$$
J(\alpha, \beta)=(-\bar{\beta}, \bar{\alpha})
$$

where $\alpha$ and $\beta$ are in $\mathbb{C}^{n}$ and $(\alpha, \beta)$ is in $\mathbb{C}^{2 n}$. We can easily check that for all $z, w \in \mathbb{C}^{2 n}$, we have

$$
\omega(z, w)=\langle J z, w\rangle
$$

Recall that we take our inner product to be conjugate linear in the first factor; since $J$ is also conjugate linear, $\langle J z, w\rangle$ is actually linear in $z$. We may easily check that

$$
\langle J z, w\rangle=-\overline{\langle z, J w\rangle}=-\langle J w, z\rangle
$$

for all $z, w \in \mathbb{C}^{2 n}$ and that

$$
J^{2}=-I .
$$

Proposition 1.5. If $U$ belongs to $\mathrm{U}(2 n)$ then $U$ belongs to $\operatorname{Sp}(n)$ if and only if $U$ commutes with $J$.

Proof. Fix some $U$ in $\mathrm{U}(2 n)$. Then for $z$ and $w$ in $\mathbb{C}^{2 n}$, we have, on the one hand,

$$
\omega(U z, U w)=\langle J U z, U w\rangle=\left\langle U^{*} J U z, w\right\rangle=\left\langle U^{-1} J U z, w\right\rangle,
$$

and, on the other hand,

$$
\omega(z, w)=\langle J z, w\rangle
$$

From this it is each to check that $U$ preserves $\omega$ if and only if

$$
U^{-1} J U=J
$$

which is equivalent to $J U=U J$.\\
The preceding result can be used to give a different perspective on the definition of $\mathrm{Sp}(n)$, as follows. The quaternion algebra $\mathbb{H}$ is the four-dimensional associative algebra over $\mathbb{R}$ spanned by elements 1 (the identity), $\mathbf{i}, \mathbf{j}$, and $\mathbf{k}$ satisfying

$$
\mathbf{i}^{2}=\mathbf{j}^{2}=\mathbf{k}^{2}=-1
$$

and

$$
\begin{aligned}
\mathbf{i j}=\mathbf{k} ; & \mathbf{j i}=-\mathbf{k} ; \\
\mathbf{j k}=\mathbf{i} ; & \mathbf{k j}=-\mathbf{i} ; \\
\mathbf{k i}=\mathbf{j} ; & \mathbf{i k}=-\mathbf{j} .
\end{aligned}
$$

We may realize the quaternion algebra inside $M_{2}(\mathbb{C})$ by taking identifying 1 with the identity matrix and setting

$$
\mathbf{i}=\left(\begin{array}{rr}
i & 0 \\
0 & -i
\end{array}\right) ; \quad \mathbf{j}=\left(\begin{array}{rr}
0 & 1 \\
-1 & 0
\end{array}\right) ; \quad \mathbf{k}=\left(\begin{array}{cc}
0 & i \\
i & 0
\end{array}\right) .
$$

The algebra $\mathbb{H}$ is then the space of real linear combinations of $I, \mathbf{i}, \mathbf{j}$, and $\mathbf{k}$.\\
Now, since $J$ is conjugate linear, we have

$$
J(i z)=-i J z
$$

for all $z \in \mathbb{C}^{2 n}$; that is, $i J=-J i$. Thus, if we define $K$ to be $i J$, we have

$$
K^{2}=i J i J=-J(i)^{2} J=J^{2}=-I
$$

and one can easily check that $i I, J$, and $K$ satisfy the same commutation relations as $\mathbf{i}$, $\mathbf{j}$, and $\mathbf{k}$. We can therefore make $\mathbb{C}^{2 n}$ into a "vector space" over the noncommutative algebra $\mathbb{H}$ by setting

$$
\begin{aligned}
\mathbf{i} \cdot z & =i z \\
\mathbf{j} \cdot z & =J z \\
\mathbf{k} \cdot z & =i J z .
\end{aligned}
$$

Now, if $U$ belongs to $\operatorname{Sp}(n)$, then $U$ commutes with multiplication by $i$ and with $J$ (Proposition 1.5) and thus, also, with $K:=i J$. Thus, $U$ is actually "quaternion linear." A $2 n \times 2 n$ matrix $U$ therefore belongs to $\mathrm{Sp}(n)$ if and only if $U$ is quaternion linear and preserves the norm. Thus, we may think of $\operatorname{Sp}(n)$ as the "unitary group over the quaternions." The compact symplectic group then fits naturally with the orthogonal groups (norm-preserving maps over $\mathbb{R}$ ) and the unitary groups (normpreserving maps over $\mathbb{C}$ ).

Every $U \in \mathrm{U}(2 n)$ has an orthonormal basis of eigenvectors, with eigenvalues having absolute value 1 . We now determine the additional properties the eigenvectors and eigenvalues must satisfy in order for $U$ to be in $\operatorname{Sp}(n)=\mathrm{U}(2 n) \cap \operatorname{Sp}(n ; \mathbb{C})$.

Theorem 1.6. If $U \in \operatorname{Sp}(n)$, then there exists an orthonormal basis $u_{1}, \ldots, u_{n}$, $v_{1}, \ldots v_{n}$ for $\mathbb{C}^{n}$ such that the following properties hold: First, $J u_{j}=v_{j}$; second, for some real numbers $\theta_{1}, \ldots, \theta_{n}$, we have

$$
\begin{aligned}
U u_{j} & =e^{i \theta_{j}} u_{j} \\
U v_{j} & =e^{-i \theta_{j}} v_{j} ;
\end{aligned}
$$

and third,

$$
\begin{aligned}
& \omega\left(u_{j}, u_{k}\right)=\omega\left(v_{j}, v_{k}\right)=0 \\
& \omega\left(u_{j}, v_{k}\right)=\delta_{j k}
\end{aligned}
$$

Conversely, if there exists an orthonormal basis with these properties, $U$ belongs to $\mathrm{Sp}(n)$.

Lemma 1.7. Suppose $V$ is a complex subspace of $\mathbb{C}^{2 n}$ that is invariant under the conjugate-linear map $J$. Then the orthogonal complement $V^{\perp}$ of $V$ (with respect to the inner product $\langle\cdot, \cdot\rangle$ ) is also invariant under $J$. Furthermore, $V$ and $V^{\perp}$ are orthogonal with respect to $\omega$; that is,

$$
\omega(z, w)=0
$$

for all $z \in V$ and $w \in V^{\perp}$.\\
Proof. If $w \in V^{\perp}$ then for all $z \in V$, we have

$$
\langle J w, z\rangle=-\langle J z, w\rangle=0
$$

because $J z$ is again in $V$. Thus, $V^{\perp}$ is invariant under $J$. Then if $z \in V$ and $w \in V^{\perp}$, we have

$$
\omega(z, w)=\langle J z, w\rangle=0
$$

because $J z$ is again in $V$.\\
Proof of Theorem 1.6. Consider $U$ in $\operatorname{Sp}(n ; \mathbb{C}) \cap \mathrm{U}(2 n)$, choose an eigenvector for $U$, normalized to be a unit vector, and call it $u_{1}$. Since $U$ preserves the norms of vectors, the eigenvalue $\lambda_{1}$ for $u_{1}$ must be of the form $e^{i \theta_{1}}$, for some $\theta_{1} \in \mathbb{R}$. If we set $v_{1}=J u_{1}$, then since $J$ is conjugate linear and commutes with $U$ (Proposition 1.5), we have

$$
U v_{1}=J\left(U u_{1}\right)=J\left(e^{i \theta_{1}} u_{1}\right)=e^{-i \theta_{1}} v_{1} .
$$

That is to say, $v_{1}$ is an eigenvector for $U$ with eigenvalue $e^{-i \theta_{1}}$. Furthermore,

$$
\left\langle v_{1}, u_{1}\right\rangle=\left\langle J u_{1}, u_{1}\right\rangle=\omega\left(u_{1}, u_{1}\right)=0
$$

since $\omega$ is a skew-symmetric form. On the other hand,

$$
\omega\left(u_{1}, v_{1}\right)=\left\langle J u_{1}, v_{1}\right\rangle=\left\langle J u_{1}, J u_{1}\right\rangle=1,
$$

since $J$ preserves the magnitude of vectors.\\
Now, since $J^{2}=-I$, we can easily check that the span $V$ of $u_{1}$ and $v_{1}=J u_{1}$ is invariant under $J$. Thus, by Lemma 1.7, $V^{\perp}$ is also invariant under $J$ and is\\
$\omega$-orthogonal to $V$. Meanwhile, $V$ is invariant under both $U$ and $U^{*}=U^{-1}$. Thus, by Proposition A.10, $V^{\perp}$ is invariant under both $U^{* *}=U$ and $U^{*}$. Since $U$ preserves $V^{\perp}$, the restriction of $U$ to $V^{\perp}$ will have an eigenvector, which we can normalize to be a unit vector and call $u_{2}$. If we let $v_{2}=J u_{2}$, then we have all the same properties for $u_{2}$ and $v_{2}$ as for $u_{1}$ and $v_{1}$. Furthermore, $u_{2}$ and $v_{2}$ are orthogonal - with respect to both $\langle\cdot, \cdot\rangle$ and $\omega(\cdot, \cdot)$ - to $u_{1}$ and $v_{1}$. We can then proceed on in a similar fashion to obtain the full set of vectors $u_{1}, \ldots, u_{n}$ and $v_{1}, \ldots, v_{n}$. (If $u_{1}, \ldots, u_{k}$ and $v_{1}, \ldots, v_{k}$ have been chosen, we take $u_{k+1}$ and $v_{k+1}:=J u_{k+1}$ in the orthogonal complement of the span of $u_{1}, \ldots, u_{k}$ and $v_{1}, \ldots, v_{k}$.)

The other direction of the theorem is left to the reader (Exercise 6).

\subsection*{Topological Properties}
In this section, we investigate three important topological properties of matrix Lie groups, each of which is satisfied by some groups but not others.

\subsubsection*{Compactness}
The first property we consider is compactness.\\
Definition 1.8. A matrix Lie group $G \subset \mathrm{GL}(n ; \mathbb{C})$ is said to be compact if it is compact in the usual topological sense as a subset of $M_{n}(\mathbb{C}) \cong \mathbb{R}^{2 n^{2}}$.

In light of the Heine-Borel theorem (Theorem 2.41 in [Rud1]), a matrix Lie group $G$ is compact if and only if it is closed (as a subset of $M_{n}(\mathbb{C})$, not just as a subset of $\mathrm{GL}(n ; \mathbb{C})$ ) and bounded. Explicitly, this means that $G$ is compact if and only if (1) whenever $A_{m} \in G$ and $A_{m} \rightarrow A$, then $A$ is in $G$, and (2) there exists a constant $C$ such that for all $A \in G$, we have $\left|A_{j k}\right| \leq C$ for all $1 \leq j, k \leq n$.

The following groups are compact: $\mathrm{O}(n)$ and $\mathrm{SO}(n), \mathrm{U}(n)$ and $\mathrm{SU}(n)$, and $\mathrm{Sp}(n)$. Each of these groups is easily seen to be closed in $M_{n}(\mathbb{C})$ and each satisfies the bound $\left|A_{j k}\right| \leq 1$, since in each case, the columns of $A \in G$ are required to be unit vectors. Most of the other groups we have considered are noncompact. The special linear group $\operatorname{SL}(n ; \mathbb{R})$, for example, is unbounded (except in the trivial case $n=1$ ), because for all $m$, the matrix

$$
A_{m}=\left(\begin{array}{ccccc}
m & & & & \\
& \frac{1}{m} & & & \\
& & 1 & & \\
& & & \ddots & \\
& & & & 1
\end{array}\right)
$$

has determinant one.

\subsubsection*{Connectedness}
The second property we consider is connectedness.\\
Definition 1.9. A matrix Lie group $G$ is said to be connected if for all $A$ and $B$ in $G$, there exists a continuous path $A(t), a \leq t \leq b$, lying in $G$ with $A(a)=A$ and $A(b)=B$. For any matrix Lie group $G$, the identity component of $G$, denoted $G_{0}$, is the set of $A \in G$ for which there exists a continuous path $A(t), a \leq t \leq b$, lying in $G$ with $A(a)=I$ and $A(b)=A$.

The property we have called "connected" in Definition 1.9 what is called path connected in topology, which is not (in general) the same as connected. However, we will eventually prove that a matrix Lie group is connected if and only if it is path-connected. Thus, in a slight abuse of terminology, we shall continue to refer to the above property as connectedness. (See the remarks following Corollary 3.45.)

To show that a matrix Lie group $G$ is connected, it suffices to show that each $A \in G$ can be connected to the identity by a continuous path lying in $G$.

Proposition 1.10. If $G$ is a matrix Lie group, the identity component $G_{0}$ of $G$ is a normal subgroup of $G$.

We will see in Sect. 3.7 that $G_{0}$ is closed and hence a matrix Lie group.\\
Proof. If $A$ and $B$ are any two elements of $G_{0}$, then there are continuous paths $A(t)$ and $B(t)$ connecting $I$ to $A$ and to $B$ in $G$. Then the path $A(t) B(t)$ is a continuous path connecting $I$ to $A B$ in $G$, and $(A(t))^{-1}$ is a continuous path connecting $I$ to $A^{-1}$ in $G$. Thus, both $A B$ and $A^{-1}$ belong to $G_{0}$, showing that $G_{0}$ is a subgroup of $G$. Now suppose $A$ is in $G_{0}$ and $B$ is any element of $G$. Then there is a continuous path $A(t)$ connecting $I$ to $A$ in $G$, and the path $B A(t) B^{-1}$ connects $I$ to $B A B^{-1}$ in $G$. Thus, $B A B^{-1} \in G_{0}$, showing that $G_{0}$ is normal.

Note that because matrix multiplication and matrix inversion are continuous on $\mathrm{GL}(n ; \mathbb{C})$, it follows that if $A(t)$ and $B(t)$ are continuous, then so are $A(t) B(t)$ and $A(t)^{-1}$. The continuity of the matrix product is obvious. The continuity of the inverse follows from the formula for the inverse in terms of cofactors; this formula is continuous as long as we remain in the set of invertible matrices where the determinant in the denominator is nonzero.

Proposition 1.11. The group $\mathrm{GL}(n ; \mathbb{C})$ is connected for all $n \geq 1$.\\
Proof. We make use of the result that every matrix is similar to an upper triangular matrix (Theorem A.4). That is to say, we can express any $A \in M_{n}(\mathbb{C})$ in the form $A=C B C^{-1}$, where

$$
B=\left(\begin{array}{lll}
\lambda_{1} & & * \\
& \ddots & \\
& & \lambda_{n}
\end{array}\right) .
$$

If $A$ is invertible, each $\lambda_{j}$ must be nonzero. Let $B(t)$ be obtained by multiplying the part of $B$ above the diagonal by ( $1-t$ ), for $0 \leq t \leq 1$, and let $A(t)=C B(t) C^{-1}$. Then $A(t)$ is a continuous path lying in $\mathrm{GL}(n ; \mathbb{C})$ which starts at $A$ and ends at $C D C^{-1}$, where $D$ is the diagonal matrix with diagonal entries $\lambda_{1}, \ldots, \lambda_{n}$. We can now define paths $\lambda_{j}(t)$ connecting $\lambda_{j}$ to 1 in $\mathbb{C}^{*}$ as $t$ goes from 1 to 2 , and we can define $A(t)$ on the interval $1 \leq t \leq 2$ by

$$
A(t)=C\left(\begin{array}{ccc}
\lambda_{1}(t) & & 0 \\
& \ddots & \\
0 & & \lambda_{n}(t)
\end{array}\right) C^{-1} .
$$

Then $A(t), 0 \leq t \leq 2$, is a continuous path in $\mathrm{GL}(n ; \mathbb{C})$ connecting $A$ to $I$.\\
An alternative proof of this result is given in Exercise 12.\\
Proposition 1.12. The group $\operatorname{SL}(n ; \mathbb{C})$ is connected for all $n \geq 1$.\\
Proof. The proof is almost the same as for $\operatorname{GL}(n ; \mathbb{C})$, except that we must make sure our path connecting $A \in \mathrm{SL}(n ; \mathbb{C})$ to $I$ lies entirely in $\operatorname{SL}(n ; \mathbb{C})$. We can ensure this by choosing $\lambda_{n}(t)$, in the second part of the preceding proof, to be equal to $\left(\lambda_{1}(t) \cdots \lambda_{n-1}(t)\right)^{-1}$.

Proposition 1.13. The groups $\mathrm{U}(n)$ and $\mathrm{SU}(n)$ are connected, for all $n \geq 1$.\\
Proof. By Theorem A.3, every unitary matrix has an orthonormal basis of eigenvectors, with eigenvalues having absolute value 1 . Thus, each $U \in \mathrm{U}(n)$ can be written as $U_{1} D U_{1}^{-1}$, where $U_{1} \in \mathrm{U}(n)$ and $D$ is diagonal with diagonal entries $e^{i \theta_{1}}, \ldots, e^{i \theta_{n}}$. We may then define

$$
U(t)=U_{1}\left(\begin{array}{ccc}
e^{i(1-t) \theta_{1}} & & 0 \\
& \ddots & \\
& & e^{i(1-t) \theta_{n}}
\end{array}\right) U_{1}^{-1}, \quad 0 \leq t \leq 1 .
$$

It is easy to see that $U(t)$ is in $\mathrm{U}(n)$ for all $t$, and $U(t)$ connects $U$ to $I$. A slight modification of this argument, as in the proof of Proposition 1.12, shows that $\mathrm{SU}(n)$ is connected.

The group $\mathrm{SO}(n)$ is also connected; see Exercise 13.

\subsubsection*{Simple Connectedness}
The last topological property we consider is simple connectedness.

Definition 1.14. A matrix Lie group $G$ is said to be simply connected if it is connected and, in addition, every loop in $G$ can be shrunk continuously to a point in $G$.

More precisely, assume that $G$ is connected. Then $G$ is simply connected if for every continuous path $A(t), 0 \leq t \leq 1$, lying in $G$ and with $A(0)=A(1)$, there exists a continuous function $A(s, t), 0 \leq s, t \leq 1$, taking values in $G$ and having the following properties: (1) $A(s, 0)=A(s, 1)$ for all $s$, (2) $A(0, t)=A(t)$, and (3) $A(1, t)=A(1,0)$ for all $t$.

One should think of $A(t)$ as a loop and $A(s, t)$ as a family of loops, parameterized by the variable $s$ which shrinks $A(t)$ to a point. Condition 1 says that for each value of the parameter $s$, we have a loop; Condition 2 says that when $s=0$ the loop is the specified loop $A(t)$; and Condition 3 says that when $s=1$ our loop is a point. The condition of simple connectedness is important because for simply connected groups, there is a particularly close relationship between the group and the Lie algebra. (See Sect. 5.7.)

Proposition 1.15. The group $\mathrm{SU}(2)$ is simply connected.\\
Proof. Exercise 5 shows that $\mathrm{SU}(2)$ may be thought of (topologically) as the three-dimensional sphere $S^{3}$ sitting inside $\mathbb{R}^{4}$. It is well known that $S^{3}$ is simply connected; see, for example, Proposition 1.14 in [Hat].

If a matrix Lie group $G$ is not simply connected, the degree to which it fails to be simply connected is encoded in the fundamental group of $G$. (See Sect. 13.1.) Sections 13.2 and 13.3 analyze several additional examples. It is shown there, for example, that $\mathrm{SU}(n)$ is simply connected for all $n$.

\subsubsection*{The Topology of SO(3)}
We conclude this section with an analysis of the topological structure of the group $\mathrm{SO}(3)$. We begin by describing real projective spaces.

Definition 1.16. The real projective space of dimension $n$, denoted $\mathbb{R} P^{n}$, is the set of lines through the origin in $\mathbb{R}^{n+1}$. Since each line through the origin intersects the unit sphere exactly twice, we may think of $\mathbb{R} P^{n}$ as the unit sphere $S^{n}$ with "antipodal" points $u$ and $-u$ identified.

Using the second description, we think of points in $\mathbb{R} P^{n}$ as pairs $\{u,-u\}$, with $u \in S^{n}$. There is a natural map $\pi: S^{n} \rightarrow \mathbb{R} P^{n}$, given by

$$
\pi(u)=\{u,-u\} .
$$

We may define a distance function on $\mathbb{R} P^{n}$ by defining

$$
\begin{aligned}
d(\{u,-u\},\{v,-v\}) & =\min (d(u, v), d(u,-v), d(-u, v), d(-u,-v)) \\
& =\min (d(u, v), d(u,-v))
\end{aligned}
$$

(The second equality holds because $d(x, y)=d(-x,-y)$.) With this metric, $\mathbb{R} P^{n}$ is locally isometric to $S^{n}$, since if $u$ and $v$ are nearby points in $S^{n}$, we have $d(\{u,-u\},\{v,-v\})=d(u, v)$.

It is known that $\mathbb{R} P^{n}$ is not simply connected. (See, for example, Example 1.43 in [Hat].) Indeed, suppose $u$ is any unit vector in $\mathbb{R}^{n+1}$ and $B(t)$ is any path in $S^{n}$ connecting $u$ to $-u$. Then

$$
A(t):=\pi(B(t))
$$

is a loop in $\mathbb{R} P^{n}$, and this loop cannot be shrunk continuously to a point in $\mathbb{R} P^{n}$. To prove this claim, suppose that a map $A(s, t)$ as in Definition 1.14. Then $A(s, t)$ can be "lifted" to a continuous map $B(s, t)$ into $S^{n}$ such that $B(0, t)=B(t)$ and such that $A(s, t)=\pi(B(s, t))$. (See Proposition 1.30 in [Hat].) Since $A(s, 0)= A(s, 1)$ for all $s$, we must have $B(s, 0)= \pm B(s, 1)$. But by construction, $B(0,0)= -B(0,1)$. If order for $B(s, t)$ to be continuous in $s$, we must then have $B(s, 0)= -B(s, 1)$ for all $s$. It follows that $B(1, t)$ is a nonconstant path in $S^{n}$. It is then easily verified that $A(1, t)=\pi(B(1, t))$ cannot be constant, contradicting our assumption about $A(s, t)$.

Let $D^{n}$ denote the closed upper hemisphere in $S^{n}$, that is, the set of points $u \in S^{n}$ with $u_{n+1} \geq 0$. Then $\pi$ maps $D^{n}$ onto $\mathbb{R} P^{n}$, since at least one of $u$ and $-u$ is in $D^{n}$. The restriction of $\pi$ to $D^{n}$ is injective except on the equator, that is, the set of $u \in S^{n}$ with $u_{n+1}=0$. If $u$ is in the equator, then $-u$ is also in the equator, and $\pi(-u)=\pi(u)$. Thus, we may also think of $\mathbb{R} P^{n}$ as the upper hemisphere $D^{n}$, with antipodal points on the equator identified (Figure 1.2).

We may now make one last identification using the projection $P$ of $\mathbb{R}^{n+1}$ onto $\mathbb{R}^{n}$. (That is to say, $P$ is the map sending $\left(x_{1}, \ldots, x_{n}, x_{n+1}\right)$ to $\left(x_{1}, \ldots, x_{n}\right)$.) The restriction of $P$ to $D^{n}$ is a continuous bijection between $D^{n}$ and the closed unit ball $B^{n}$ in $\mathbb{R}^{n}$, with the equator in $D^{n}$ mapping to the boundary of the ball. Thus, our


\includegraphics[scale=0.2, center]{2025_08_28_fc86fd4ade69f41b5842g-033}

Fig. 1.2 The space $\mathbb{R} P^{n}$ is the upper hemisphere with antipodal points on the equator identified. The indicated path from $u$ to $-u$ corresponds to a loop in $\mathbb{R} P^{n}$ that cannot be shrunk to a point.


last model of $\mathbb{R} P^{n}$ is the closed unit ball $B^{n} \subset \mathbb{R}^{n}$, with antipodal points on the boundary of $B^{n}$ identified.

We now turn to a topological analysis of $\mathrm{SO}(3)$.\\
Proposition 1.17. There is a continuous bijection between $\mathrm{SO}(3)$ and $\mathbb{R} P^{3}$.\\
Since $\mathbb{R} P^{3}$ is not simply connected, it follows that $\mathrm{SO}(3)$ is not simply connected, either.

Proof. If $v$ is a unit vector in $\mathbb{R}^{3}$, let $R_{v, \theta}$ be the element of $\mathrm{SO}(3)$ consisting of a "right-handed" rotation by angle $\theta$ in the plane orthogonal to $v$. That is to say, let $v^{\perp}$ denote the plane orthogonal to $v$ and choose an orthonormal basis ( $u_{1}, u_{2}$ ) for $v^{\perp}$ in such a way that the linear map taking the orthonormal basis $\left(u_{1}, u_{2}, v\right)$ to the standard basis $\left(e_{1}, e_{2}, e_{3}\right)$ has positive determinant. We use the basis $\left(u_{1}, u_{2}\right)$ to identify $v^{\perp}$ with $\mathbb{R}^{2}$, and the rotation is then in the counterclockwise direction in $\mathbb{R}^{2}$. It is easily seen that $R_{-v, \theta}$ is the same as $R_{v,-\theta}$. It is also not hard to show (Exercise 14) that every element of $\mathrm{SO}(3)$ can be expressed as $R_{v, \theta}$, for some $v$ and $\theta$ with $-\pi \leq \theta \leq \pi$. Furthermore, we can arrange that $0 \leq \theta \leq \pi$ by replacing $v$ with $-v$ if necessary.

If $R=I$, then $R=R_{v, 0}$ for any unit vector $v$. If $R$ is a rotation by angle $\pi$ about some axis $v$, then $R$ can be expressed both as $R_{v, \pi}$ and as $R_{-v, \pi}$. It is not hard to see that if $R \neq I$ and $R$ is not a rotation by angle $\pi$, then $R$ has a unique representation as $R_{v, \theta}$ with $0<\theta<\pi$.

Now let $B^{3}$ denote the closed ball of radius $\pi$ in $\mathbb{R}^{3}$ and consider the map $\Phi$ : $B^{3} \rightarrow \mathrm{SO}(3)$ given by

$$
\begin{aligned}
& \Phi(u)=R_{\hat{u},\|u\|}, \quad u \neq 0 \\
& \Phi(0)=I .
\end{aligned}
$$

Here, $\hat{u}=u /\|u\|$ is the unit vector in the $u$-direction. The map $\Phi$ is continuous, even at $I$, since $R_{v, \theta}$ approaches the identity as $\theta$ approaches zero, regardless of how $v$ is behaving. The discussion in the preceding paragraph shows that $\Phi$ maps $B^{3}$ onto SO(3). The map $\Phi$ is injective except that "antipodal" points on the boundary of $B^{3}$ have the same image: $R_{v, \pi}=R_{-v, \pi}$. Thus, $\Phi$ descends to a continuous, injective map of $\mathbb{R} P^{3}$ onto SO (3). Since both $\mathbb{R} P^{3}$ and SO (3) are compact, Theorem 4.17 in [Rud1] tells us that the inverse map is also continuous, meaning that SO (3) is homeomorphic to $\mathbb{R} P^{3}$.

For a different approach to proving Proposition 1.17, see the discussion following Proposition 1.19.

\subsection*{Homomorphisms}
We now look at the notion of homomorphisms for matrix Lie groups.

Definition 1.18. Let $G$ and $H$ be matrix Lie groups. A map $\Phi$ from $G$ to $H$ is called a Lie group homomorphism if (1) $\Phi$ is a group homomorphism and (2) $\Phi$ is continuous. If, in addition, $\Phi$ is one-to-one and onto and the inverse map $\Phi^{-1}$ is continuous, then $\Phi$ is called a Lie group isomorphism.

The condition that $\Phi$ be continuous should be regarded as a technicality, in that it is very difficult to give an example of a group homomorphism between two matrix Lie groups which is not continuous. In fact, if $G=\mathbb{R}$ and $H=\mathbb{C}^{*}$, then any group homomorphism from $G$ to $H$ which is even measurable (a very weak condition) must be continuous. (See Exercise 17 in Chapter 9 of [Rud2].)

Note that the inverse of a Lie group isomorphism is continuous (by definition) and a group homomorphism (by elementary group theory), and thus a Lie group isomorphism. If $G$ and $H$ are matrix Lie groups and there exists a Lie group isomorphism from $G$ to $H$, then $G$ and $H$ are said to be isomorphic, and we write $G \cong H$.

The simplest interesting example of a Lie group homomorphism is the determinant, which is a homomorphism of $\operatorname{GL}(n ; \mathbb{C})$ into $\mathbb{C}^{*}$. Another simple example is the map $\Phi: \mathbb{R} \rightarrow \mathrm{SO}(2)$ given by

$$
\Phi(\theta)=\left(\begin{array}{rr}
\cos \theta & -\sin \theta \\
\sin \theta & \cos \theta
\end{array}\right) .
$$

This map is clearly continuous, and calculation (using standard trigonometric identities) shows that it is a homomorphism.

An important topic for us will be the relationship between the groups $\operatorname{SU}(2)$ and SO(3), which are almost, but not quite, isomorphic. Specifically, we now show that there exists a Lie group homomorphism $\Phi: \mathrm{SU}(2) \rightarrow \mathrm{SO}(3)$ that is two-to-one and onto. Consider the space $V$ of all $2 \times 2$ complex matrices $X$ which are self-adjoint (i.e., $X^{*}=X$ ) and have trace zero. Elements of $V$ are precisely the matrices of the form

$$
X=\left(\begin{array}{cc}
x_{1} & x_{2}+i x_{3} \\
x_{2}-i x_{3} & -x_{1}
\end{array}\right),
$$

with $x_{1}, x_{2}, x_{3} \in \mathbb{R}$. If we identify $V$ with $\mathbb{R}^{3}$ by means of the coordinates $x_{1}, x_{2}$, and $x_{3}$ in (1.14), then the standard inner product on $\mathbb{R}^{3}$ can be computed as

$$
\left\langle X_{1}, X_{2}\right\rangle=\frac{1}{2} \operatorname{trace}\left(X_{1} X_{2}\right) .
$$

That is to say,

$$
\begin{aligned}
& \frac{1}{2} \operatorname{trace}\left(\left(\begin{array}{cc}
x_{1} & x_{2}+i x_{3} \\
x_{2}-i x_{3} & -x_{1}
\end{array}\right)\left(\begin{array}{cc}
x_{1}^{\prime} & x_{2}^{\prime}+i x_{3}^{\prime} \\
x_{2}^{\prime}-i x_{3}^{\prime} & -x_{1}^{\prime}
\end{array}\right)\right) \\
& \quad=x_{1} x_{1}^{\prime}+x_{2} x_{2}^{\prime}+x_{3} x_{3}^{\prime}
\end{aligned}
$$

as one may easily check by direct calculation.\\
For each $U \in \mathrm{SU}(2)$, define a linear map $\Phi_{U}: V \rightarrow V$ by

$$
\Phi_{U}(X)=U X U^{-1} .
$$

Since $U$ is unitary,

$$
\left(U A U^{-1}\right)^{*}=\left(U^{-1}\right)^{*} A U^{*}=U A U^{-1},
$$

showing that $U A U^{-1}$ is again in $V$.\\
It is easy to see that $\Phi_{U_{1} U_{2}}=\Phi_{U_{1}} \Phi_{U_{2}}$. Furthermore,

$$
\begin{aligned}
\frac{1}{2} \operatorname{trace}\left(\left(U X_{1} U^{-1}\right)\left(U X_{2} U^{-1}\right)\right) & =\frac{1}{2} \operatorname{trace}\left(U X_{1} X_{2} U^{-1}\right) \\
& =\frac{1}{2} \operatorname{trace}\left(X_{1} X_{2}\right),
\end{aligned}
$$

since the trace is invariant under conjugation. Thus, each $\Phi_{U}$ preserves the inner product trace $\left(X_{1} X_{2}\right) / 2$ on $V$. It follows that the map $U \mapsto \Phi_{U}$ is a homomorphism of SU (2) into the group of orthogonal linear transformations of $V \cong \mathbb{R}^{3}$, that is, into $\mathrm{O}(3)$. Since $\mathrm{SU}(2)$ is connected (Proposition 1.13), $\Phi_{U}$ must actually lie in $\mathrm{SO}(3)$ for all $U \in \mathrm{SU}(2)$. Thus, $\Phi$ (i.e., the map $U \mapsto \Phi_{U}$ ) is a homomorphism of $\mathrm{SU}(2)$ into $\mathrm{SO}(3)$, which is easily seen to be continuous. Since $(-I) X(-I)^{-1}=X$, we see that $\Phi_{-I}$ is the identity element of $\operatorname{SO}(3)$.

Suppose, for example, that $U$ is the matrix

$$
U=\left(\begin{array}{cc}
e^{i \theta / 2} & 0 \\
0 & e^{-i \theta / 2}
\end{array}\right) .
$$

Then by direct calculation, we obtain

$$
U\left(\begin{array}{cc}
x_{1} & x_{2}+i x_{3} \\
x_{2}-i x_{3} & -x_{1}
\end{array}\right) U^{-1}=\left(\begin{array}{cc}
x_{1}^{\prime} & x_{2}^{\prime}+i x_{3}^{\prime} \\
x_{2}^{\prime}-i x_{3}^{\prime} & -x_{1}^{\prime}
\end{array}\right),
$$

where $x_{1}^{\prime}=x_{1}$ and

$$
\begin{aligned}
x_{2}^{\prime}+i x_{3}^{\prime} & =e^{i \theta}\left(x_{2}+i x_{3}\right) \\
& =\left(x_{2} \cos \theta-x_{3} \sin \theta\right)+i\left(x_{2} \sin \theta+x_{3} \cos \theta\right) .
\end{aligned}
$$

In this case, then, $\Phi_{U}$ is a rotation by angle $\theta$ in the ( $x_{2}, x_{3}$ )-plane. Note that even though the diagonal entries of $U$ are $e^{ \pm i \theta / 2}$, the map $\Phi_{U}$ is a rotation by angle $\theta$, not $\theta / 2$.

Proposition 1.19. The map $U \mapsto \Phi_{U}$ is a 2-1 and onto map of $\mathrm{SU}(2)$ to $\mathrm{SO}(3)$, with kernel equal to $\{I,-I\}$.

Since $\mathrm{SU}(2)$ is homeomorphic to $S^{3}$, the proposition gives another way of seeing that SO (3) is homeomorphic to $\mathbb{R} P^{3}$, that is, $S^{3}$ with antipodal points identified. This result was obtained in a different way in Proposition 1.17.

It is not hard to show that $\Phi$ is a "covering map" in the topological sense (Section 1.3 of [Hat]). Since SU(2) is simply connected (Proposition 1.15) and the map is 2-1, it follows by the theory of covering maps (e.g., Theorem 1.38 in [Hat]) that SO(3) cannot be simply connected and, indeed, it must have fundamental group $\mathbb{Z} / 2$. See Chapter 13 for general information about fundamental groups and for a computation of the fundamental group of $\mathrm{SO}(n), n>2$.

Proof. Exercise 16 shows that the kernel of $\Phi$ is precisely the set $\{I,-I\}$. To see that $\Phi$ maps onto $\mathrm{SO}(3)$, let $R$ be a rotation of $V \cong \mathbb{R}^{3}$. By Exercise 14, there exists an "axis" $X \in V$ such that $R$ is a rotation by some angle $\theta$ in the plane orthogonal to $X$. If we express $X$ in the form

$$
X=U_{0}\left(\begin{array}{cc}
x_{1} & 0 \\
0 & -x_{1}
\end{array}\right) U_{0}^{-1}
$$

with $U_{0} \in \mathrm{U}(2)$, then the plane orthogonal to $X$ in $V$ is the space of matrices of the form

$$
X^{\prime}=U_{0}\left(\begin{array}{cc}
0 & x_{2}+i x_{3} \\
x_{2}-i x_{3} & 0
\end{array}\right) U_{0}^{-1} .
$$

If we now take

$$
U=U_{0}\left(\begin{array}{cc}
e^{i \theta / 2} & 0 \\
0 & e^{-i \theta / 2}
\end{array}\right) U_{0}^{-1}
$$

we can easily see that $U X U^{-1}=X$. On the other hand, the calculations in (1.15) and (1.16) show that $U X^{\prime} U^{-1}$ is of the same form as in (1.17), but with ( $x_{2}, x_{3}$ ) rotated by angle $\theta$. Thus, $\Phi_{U}$ is a rotation by angle $\theta$ in the plane perpendicular to $X$, showing that $\Phi_{U}$ coincides with $R$.

It is possible, if not terribly useful, to calculate $\Phi$ explicitly. If you write an element of $\operatorname{SU}(2)$ as in Exercise 5, you (or your computer) may calculate

$$
\begin{aligned}
& \left(\begin{array}{cc}
\alpha & -\bar{\beta} \\
\beta & \bar{\alpha}
\end{array}\right)\left(\begin{array}{cc}
x_{1} & x_{2}+i x_{3} \\
x_{2}-i x_{3} & -x_{1}
\end{array}\right)\left(\begin{array}{rc}
\bar{\alpha} & \bar{\beta} \\
-\beta & \alpha
\end{array}\right) \\
& =\left(\begin{array}{cc}
x_{1}^{\prime} & x_{2}^{\prime}+i x_{3}^{\prime} \\
x_{2}^{\prime}-i x_{3}^{\prime} & -x_{3}^{\prime}
\end{array}\right)
\end{aligned}
$$

explicitly. Then ( $x_{1}^{\prime}, x_{2}^{\prime}, x_{3}^{\prime}$ ) will depend linearly on ( $x_{1}, x_{2}, x_{3}$ ) and you can express $\left(x_{1}^{\prime}, x_{2}^{\prime}, x_{3}^{\prime}\right)$ as a matrix applied to $\left(x_{1}, x_{2}, x_{3}\right)$, with the result that

$$
\Phi_{U}=\left(\begin{array}{ccc}
|\alpha|^{2}-|\beta|^{2} & -2 \operatorname{Re}(\alpha \beta) & 2 \operatorname{Im}(\alpha \beta) \\
2 \operatorname{Re}(\alpha \bar{\beta}) & \operatorname{Re}\left(\alpha^{2}-\beta^{2}\right) & \operatorname{Im}\left(\beta^{2}-\alpha^{2}\right) \\
2 \operatorname{Im}(\alpha \bar{\beta}) & \operatorname{Im}\left(\alpha^{2}+\beta^{2}\right) & \operatorname{Re}\left(\alpha^{2}+\beta^{2}\right)
\end{array}\right) .
$$

If we take $\alpha=e^{i \theta / 2}$ and $\beta=0$, we may see directly that $\Phi_{U}$ is a rotation by angle $\theta$ in the ( $x_{2}, x_{3}$ )-plane, as we saw already in (1.15) and (1.16).

\subsection*{Lie Groups}
A Lie group is a smooth manifold equipped with a group structure such that the operations of group multiplication and inversion are smooth. As the terminology suggests, every matrix Lie group is a Lie group. (See Corollary 3.45 in Chapter 3.) The reverse is not true: Not every Lie group is isomorphic to a matrix Lie group. Nevertheless, we have restricted our attention in this book to matrix Lie groups, in order to minimize prerequisites and keep the discussion as concrete as possible. Most of the interesting examples of Lie groups are, in any case, matrix Lie groups.

A manifold is an object $M$ that looks locally like a piece of $\mathbb{R}^{n}$. More precisely, an $n$-dimensional manifold is a second-countable, Hausdorff topological space with the property that each $m \in M$ has a neighborhood that is homeomorphic to an open subset of $\mathbb{R}^{n}$. A two-dimensional torus, for example, looks locally but not globally like $\mathbb{R}^{2}$ and is, thus, a two-dimensional manifold. A smooth manifold is a manifold $M$ together with a collection of local coordinates covering $M$ such that the change-of-coordinates map between two overlapping coordinate systems is smooth.

Definition 1.20. A Lie group is a smooth manifold $G$ which is also a group and such that the group product

$$
G \times G \rightarrow G
$$

and the inverse map $G \rightarrow G$ are smooth.\\
Example 1.21. Let

$$
G=\mathbb{R} \times \mathbb{R} \times S^{1}=\left\{(x, y, u) \mid x \in \mathbb{R}, y \in \mathbb{R}, u \in S^{1} \subset \mathbb{C}\right\}
$$

equipped with the group product given by

$$
\left(x_{1}, y_{1}, u_{1}\right) \cdot\left(x_{2}, y_{2}, u_{2}\right)=\left(x_{1}+x_{2}, y_{1}+y_{2}, e^{i x_{1} y_{2}} u_{1} u_{2}\right) .
$$

Then $G$ is a Lie group.

Proof. It is easily checked that this operation is associative; the product of three elements with either grouping is

$$
\left(x_{1}+x_{2}+x_{3}, y_{1}+y_{2}+y_{3}, e^{i\left(x_{1} y_{2}+x_{1} y_{3}+x_{2} y_{3}\right)} u_{1} u_{2} u_{3}\right) .
$$

There is an identity element in $G$, namely $e=(0,0,1)$ and each element ( $x, y, u$ ) has an inverse given by $\left(-x,-y, e^{i x y} u^{-1}\right)$. Thus, $G$ is, in fact, a group. Furthermore, both the group product and the map that sends each element to its inverse are clearly smooth, showing that $G$ is a Lie group.

Although there is nothing about matrices in the definition of $G$, we may still ask whether $G$ is isomorphic to some matrix Lie group. This turns out to be false. As shown in Sect. 4.8, there is no continuous, injective homomorphism of $G$ into any $\mathrm{GL}(n ; \mathbb{C})$. We conclude, then, that not every Lie group is isomorphic to a matrix Lie group. Nevertheless, most of the interesting examples of Lie groups are matrix Lie groups.

Let us now think briefly about how we might show that every matrix Lie group is a Lie group. We will prove in Sect. 3.7 that every matrix Lie group is an "embedded submanifold" of $M_{n}(\mathbb{R}) \cong \mathbb{R}^{2 n^{2}}$. The operations of matrix multiplication and inversion are smooth on $M_{n}(\mathbb{C})$ (after restricting to the open subset of invertible matrices in the case of inversion). Thus, the restriction of these operations to a matrix Lie group $G \subset M_{n}(\mathbb{C})$ is also smooth, making $G$ into a Lie group.

It is customary to call a map $\Phi$ between two Lie groups a Lie group homomorphism if $\Phi$ is a group homomorphism and $\Phi$ is smooth, whereas we have (in Definition 1.18) required only that $\Phi$ be continuous. We will show, however, that every continuous homomorphism between matrix Lie groups is automatically smooth, so that there is no conflict of terminology. See Corollary 3.50 to Theorem 3.42. Finally, we note that since every matrix Lie group $G$ is a manifold, $G$ must be locally path connected. It then follows by a standard topological argument that $G$ is connected if and only if it is path connected.

\subsection*{Exercises}
\begin{enumerate}
  \item Let $[\cdot, \cdot]_{n, k}$ be the symmetric bilinear form on $\mathbb{R}^{n+k}$ defined in (1.5). Let $g$ be the $(n+k) \times(n+k)$ diagonal matrix with first $n$ diagonal entries equal to one and last $k$ diagonal entries equal to minus one:
\end{enumerate}

$$
g=\left(\begin{array}{rr}
I_{n} & 0 \\
0 & -I_{k}
\end{array}\right) .
$$

Show that for all $x, y \in \mathbb{R}^{n+k}$,

$$
[x, y]_{n, k}=\langle x, g y\rangle .
$$

Show that a $(n+k) \times(n+k)$ real matrix $A$ belongs to $\mathrm{O}(n ; k)$ if and only if $g A^{t r} g=A^{-1}$.\\
2. Let $\omega$ be the skew-symmetric bilinear form on $\mathbb{R}^{2 n}$ given by (1.7). Let $\Omega$ be the $2 n \times 2 n$ matrix

$$
\Omega=\left(\begin{array}{rr}
0 & I \\
-I & 0
\end{array}\right) .
$$

Show that for all $x, y \in \mathbb{R}^{2 n}$, we have

$$
\omega(x, y)=\langle x, \Omega y\rangle .
$$

Show that a $2 n \times 2 n$ matrix $A$ belongs to $\operatorname{Sp}(n ; \mathbb{R})$ if and only if $-\Omega A^{t r} \Omega= A^{-1}$.\\
Note: A similar analysis applies to $\operatorname{Sp}(n ; \mathbb{C})$.\\
3. Show that the symplectic group $\operatorname{Sp}(1 ; \mathbb{R}) \subset \operatorname{GL}(2 ; \mathbb{R})$ is equal to $\operatorname{SL}(2 ; \mathbb{R})$. Show that $\operatorname{Sp}(1 ; \mathbb{C})=\operatorname{SL}(2 ; \mathbb{C})$ and that $\operatorname{Sp}(1)=\operatorname{SU}(2)$.\\
4. Show that a matrix $R$ belongs to $\mathrm{SO}(2)$ if and only if it can be expressed in the form

$$
\left(\begin{array}{rr}
\cos \theta & -\sin \theta \\
\sin \theta & \cos \theta
\end{array}\right)
$$

for some $\theta \in \mathbb{R}$. Show that a matrix $R$ belongs to $\mathrm{O}(2)$ if and only if it is of one of the two forms:

$$
A=\left(\begin{array}{rr}
\cos \theta & -\sin \theta \\
\sin \theta & \cos \theta
\end{array}\right) \quad \text { or } \quad A=\left(\begin{array}{rr}
\cos \theta & \sin \theta \\
\sin \theta & -\cos \theta
\end{array}\right) .
$$

Hint: Recall that for $A$ to be in $\mathrm{O}(2)$, the columns of $A$ must be orthonormal.\\
5. Show that if $\alpha$ and $\beta$ are arbitrary complex numbers satisfying $|\alpha|^{2}+|\beta|^{2}=1$, then the matrix

$$
A=\left(\begin{array}{rr}
\alpha & -\bar{\beta} \\
\beta & \bar{\alpha}
\end{array}\right)
$$

is in $\mathrm{SU}(2)$. Show that every $A \in \mathrm{SU}(2)$ can be expressed in this form for a unique pair $(\alpha, \beta)$ satisfying $|\alpha|^{2}+|\beta|^{2}=1$.\\
6. Suppose $U$ belongs to $M_{2 n}(\mathbb{C})$ and $U$ has an orthonormal basis of eigenvectors satisfying the conditions in Theorem 1.6. Show that $U$ belongs to $\operatorname{Sp}(n)$.\\
Hint: Start by showing that $U$ is unitary. Then show that $\omega(U z, U w)=\omega(z, w)$ if $z$ and $w$ belong to the basis $u_{1}, \ldots, u_{n}, v_{1}, \ldots, v_{n}$.\\
7. Using Theorem 1.6, show that $\operatorname{Sp}(n)$ is connected and that every element of $\operatorname{Sp}(n)$ has determinant 1.\\
8. Determine the center $Z(H)$ of the Heisenberg group $H$. Show that the quotient group $H / Z(H)$ is commutative.\\
9. Suppose $a$ is an irrational real number. Show that the set $E_{a}$ of numbers of the form $e^{2 \pi i n a}, n \in \mathbb{Z}$, is dense in the unit circle $S^{1}$.\\
Hint: Show that if we divide $S^{1}$ into $N$ equally sized "bins" of length $2 \pi / N$, there is at least one bin that contains infinitely many elements of $E_{a}$. Then use the fact that $E_{a}$ is a subgroup of $S^{1}$.\\
10. Let $a$ be an irrational real number and let $G$ be the following subgroup of GL(2; $\mathbb{C}$ ):

$$
G=\left\{\left.\left(\begin{array}{cc}
e^{i t} & 0 \\
0 & e^{i t a}
\end{array}\right) \right\rvert\, t \in \mathbb{R}\right\} .
$$

Show that

$$
\bar{G}=\left\{\left.\left(\begin{array}{cc}
e^{i \theta} & 0 \\
0 & e^{i \phi}
\end{array}\right) \right\rvert\, \theta, \phi \in \mathbb{R}\right\},
$$

where $\bar{G}$ denotes the closure of the set $G$ inside the space of $2 \times 2$ matrices.\\
Hint: Use Exercise 9.\\
11. A subset $E$ of a matrix Lie group $G$ is called discrete if for each $A$ in $E$ there is a neighborhood $U$ of $A$ in $G$ such that $U$ contains no point in $E$ except for $A$. Suppose that $G$ is a connected matrix Lie group and $N$ is a discrete normal subgroup of $G$. Show that $N$ is contained in the center of $G$.\\
12. This problem gives an alternative proof of Proposition 1.11, namely that $\mathrm{GL}(n ; \mathbb{C})$ is connected. Suppose $A$ and $B$ are invertible $n \times n$ matrices. Show that there are only finitely many complex numbers $\lambda$ for which $\operatorname{det}(\lambda A+(1-\lambda) B)=0$. Show that there exists a continuous path $A(t)$ of the form $A(t)=\lambda(t) A+(1-\lambda(t)) B$ connecting $A$ to $B$ and such that $A(t)$ lies in $\mathrm{GL}(n ; \mathbb{C})$. Here, $\lambda(t)$ is a continuous path in the plane with $\lambda(0)=0$ and $\lambda(1)=1$.\\
13. Show that $\mathrm{SO}(n)$ is connected, using the following outline.

For the case $n=1$, there is nothing to show, since a $1 \times 1$ matrix with determinant one must be [1]. Assume, then, that $n \geq 2$. Let $e_{1}$ denote the unit vector with entries $1,0, \ldots, 0$ in $\mathbb{R}^{n}$. For every unit vector $v \in \mathbb{R}^{n}$, show that there exists a continuous path $R(t)$ in $\mathrm{SO}(n)$ with $R(0)=I$ and $R(1) v=e_{1}$. (Thus, any unit vector can be "continuously rotated" to $e_{1}$.)\\
Now, show that any element $R$ of $\mathrm{SO}(n)$ can be connected to a block-diagonal matrix of the form

$$
\left(\begin{array}{ll}
1 & \\
& R_{1}
\end{array}\right)
$$

with $R_{1} \in \mathrm{SO}(n-1)$ and proceed by induction.\\
14. If $R$ is an element of $\mathrm{SO}(3)$, show that $R$ must have an eigenvector $v$ with eigenvalue 1 . Show that $R$ maps the plane orthogonal to $v$ into itself. Conclude that $R$ is a rotation by some angle $\theta$ around the "axis" $v$.\\
Hint: Since $\mathrm{SO}(3) \subset \mathrm{SU}(3)$, every (real or complex) eigenvalue of $R$ must have absolute value 1 . Since, also, $R$ is real, any nonreal eigenvalues of $R$ come in conjugate pairs.\\
15. Let $R$ be an element of $\mathrm{SO}(n)$.\\
(a) Suppose $v \in \mathbb{C}^{n}$ is an eigenvector for $R$ with eigenvalue $\lambda \in \mathbb{C}$, and suppose $\lambda$ is not real. Let $V \subset \mathbb{R}^{n}$ be the two-dimensional span of $(v+\bar{v}) / 2$ and $(v-\bar{v}) /(2 i)$. Show that $V$ is invariant under $R$ and that the restriction of $R$ to $V$ has determinant 1 .\\
(b) Suppose that a subspace $V \subset \mathbb{R}^{n}$ is invariant under both $R$ and $R^{-1}$. Show that the orthogonal complement $V^{\perp}$ of $V$ is also invariant under both $R$ and $R^{-1}$.\\
(c) Show that if $n=2 k$, there exists $S \in \mathrm{SO}(n)$ such that

$$
R=S\left(\begin{array}{rrrr}
\cos \theta_{1} & -\sin \theta_{1} & & \\
\sin \theta_{1} & \cos \theta_{1} & & \\
& & \ddots & \\
& & & \cos \theta_{k}-\sin \theta_{k} \\
& & & \sin \theta_{k} \\
& & \cos \theta_{k}
\end{array}\right) S^{-1}
$$

and that if $n=2 k+1$, there exists $S \in \mathrm{SO}(n)$ such that

$$
R=S\left(\begin{array}{ccccc}
\cos \theta_{1}-\sin \theta_{1} & & & \\
\sin \theta_{1} & \cos \theta_{1} & & & \\
& & \ddots & & \\
& & & \cos \theta_{k}-\sin \theta_{k} & \\
& & & \sin \theta_{k} & \cos \theta_{k} \\
& & & & 1
\end{array}\right) S^{-1}
$$

That is, in a suitable orthonormal basis, $R$ is block diagonal with $2 \times 2$ blocks of the indicated form, with a single $1 \times 1$ block if $n$ is odd.

Hint: Show that the number of eigenvalues of $R$ equal to -1 is even.\\
16. (a) Show that if a matrix $A$ commutes with every matrix $X$ of the form (1.14), then $A$ commutes with every element of $M_{2}(\mathbb{C})$. Conclude that $A$ must be a multiple of the identity.\\
(b) Show that the kernel of the map $U \mapsto \Phi_{U}$ in Proposition 1.19 is precisely the set $\{I,-I\}$.\\
17. Suppose $G \subset \mathrm{GL}\left(n_{1} ; \mathbb{C}\right)$ and $H \subset \mathrm{GL}\left(n_{2} ; \mathbb{C}\right)$ are matrix Lie groups and that $\Phi: G \rightarrow H$ is a Lie group homomorphism. Then the image of $G$ under $\Phi$ is a subgroup of $H$ and thus of $\mathrm{GL}\left(n_{2} ; \mathbb{C}\right)$. Is the image of $G$ under $\Phi$ necessarily a matrix Lie group? Prove or give a counter-example.\\
18. Show that every continuous homomorphism $\Phi$ from $\mathbb{R}$ to $S^{1}$ is of the form $\Phi(x)=e^{i a x}$ for some $a \in \mathbb{R}$.\\
Hint: Since $\Phi$ is continuous, there is some $\varepsilon>0$ such that if $|x|<\varepsilon$, then $\Phi(x)$ belongs to the right half of the unit circle.

\section*{Chapter 2 The Matrix Exponential}
\subsection*{The Exponential of a Matrix}
The exponential of a matrix plays a crucial role in the theory of Lie groups. The exponential enters into the definition of the Lie algebra of a matrix Lie group (Sect. 3.3) and is the mechanism for passing information from the Lie algebra to the Lie group.

If $X$ is an $n \times n$ matrix, we define the exponential of $X$, denoted $e^{X}$ or $\exp X$, by the usual power series

$$
e^{X}=\sum_{m=0}^{\infty} \frac{X^{m}}{m!},
$$

where $X^{0}$ is defined to be the identity matrix $I$ and where $X^{m}$ is the repeated matrix product of $X$ with itself.

Proposition 2.1. The series (2.1) converges for all $X \in M_{n}(\mathbb{C})$ and $e^{X}$ is a continuous function of $X$.

Our proof will use the notion of the norm of a matrix $X \in M_{n}(\mathbb{C})$, which we define by thinking of $M_{n}(\mathbb{C})$ as $\mathbb{C}^{n^{2}}$.

Definition 2.2. For any $X \in M_{n}(\mathbb{C})$, we define

$$
\|X\|=\left(\sum_{j, k=1}^{n}\left|X_{j k}\right|^{2}\right)^{1 / 2}
$$

The quantity $\|X\|$ is called the Hilbert-Schmidt norm of $X$.

The Hilbert-Schmidt norm may be computed in a basis-independent way as

$$
\|X\|=\left(\operatorname{trace}\left(X^{*} X\right)\right)^{1 / 2}
$$

(See Sect. A.6.) This norm satisfies the inequalities

$$
\begin{aligned}
\|X+Y\| & \leq\|X\|+\|Y\| \\
\|X Y\| & \leq\|X\|\|Y\|
\end{aligned}
$$

for all $X, Y \in M_{n}(\mathbb{C})$. The first of these inequalities is the triangle inequality for $\mathbb{C}^{n^{2}}$ and the second follows from the Cauchy-Schwarz inequality (Exercise 1). If $X_{m}$ is a sequence of matrices, then it is easy to see that $X_{m}$ converges to a matrix $X$ in the sense of Definition 1.3 if and only if $\left\|X_{m}-X\right\| \rightarrow 0$ as $m \rightarrow \infty$.

Proof of Proposition 2.1. In light of (2.4), we see that

$$
\left\|X^{m}\right\| \leq\|X\|^{m}
$$

for all $m \geq 1$, and, hence,

$$
\sum_{m=0}^{\infty}\left\|\frac{X^{m}}{m!}\right\| \leq\|I\|+\sum_{m=1}^{\infty} \frac{\|X\|^{m}}{m!}<\infty
$$

Thus, the series (2.1) converges absolutely.\\
To show continuity, note that since $X^{m}$ is a continuous function of $X$, the partial sums of (2.1) are continuous. By the Weierstrass M-test, the series (2.1) converges uniformly on each set of the form $\{\|X\| \leq R\}$. Thus, $e^{X}$ is continuous on each such set, and, thus, continuous on all of $M_{n}(\mathbb{C})$.

We now list some elementary properties of the matrix exponential.\\
Proposition 2.3 Let $X$ and $Y$ be arbitrary $n \times n$ matrices. Then we have the following:

\begin{enumerate}
  \item $e^{0}=I$.
  \item $\left(e^{X}\right)^{*}=e^{X^{*}}$.
  \item $e^{X}$ is invertible and $\left(e^{X}\right)^{-1}=e^{-X}$.
  \item $e^{(\alpha+\beta) X}=e^{\alpha X} e^{\beta X}$ for all $\alpha$ and $\beta$ in $\mathbb{C}$.
  \item If $X Y=Y X$, then $e^{X+Y}=e^{X} e^{Y}=e^{Y} e^{X}$.
  \item If $C$ is invertible, then $e^{C X C^{-1}}=C e^{X} C^{-1}$.
\end{enumerate}

Although $e^{X+Y}=e^{X} e^{Y}$ when $X$ and $Y$ commute, this identity fails in general. This is an important point, which we will return to in the Lie product formula in Sect. 2.4 and the Baker-Campbell-Hausdorff formula in Chapter 5.

Proof. Point 1 is obvious and Point 2 follows from taking term-by-term adjoints of the series for $e^{X}$. Points 3 and 4 are special cases of Point 5. To verify Point 5,\\
we simply multiply the two power series term by term, which is permitted because both series converge absolutely. Multiplying out $e^{X} e^{Y}$ and collecting terms where the power of $X$ plus the power of $Y$ equals $m$, we obtain

$$
e^{X} e^{Y}=\sum_{m=0}^{\infty} \sum_{k=0}^{m} \frac{X^{k}}{k!} \frac{Y^{m-k}}{(m-k)!}=\sum_{m=0}^{\infty} \frac{1}{m!} \sum_{k=0}^{m} \frac{m!}{k!(m-k)!} X^{k} Y^{m-k}
$$

Now, because (and only because) $X$ and $Y$ commute,

$$
(X+Y)^{m}=\sum_{k=0}^{m} \frac{m!}{k!(m-k)!} X^{k} Y^{m-k}
$$

and, thus, (2.5) becomes

$$
e^{X} e^{Y}=\sum_{m=0}^{\infty} \frac{1}{m!}(X+Y)^{m}=e^{X+Y}
$$

To prove Point 6, simply note that

$$
\left(C X C^{-1}\right)^{m}=C X^{m} C^{-1}
$$

and, thus, the two sides of Point 6 are equal term by term.\\
Proposition 2.4. Let $X$ be a $n \times n$ complex matrix. Then $e^{t X}$ is a smooth curve in $M_{n}(\mathbb{C})$ and

$$
\frac{d}{d t} e^{t X}=X e^{t X}=e^{t X} X
$$

In particular,

$$
\left.\frac{d}{d t} e^{t X}\right|_{t=0}=X
$$

Results that hold for the exponential of numbers may or may not hold for the matrix exponential. Although Proposition 2.4 is what one would expect from the scalar case, it should be noted that, in general, the derivative of $e^{X+t Y}$ is not equal to $e^{X+t Y} Y$. See Sect. 5.4.

Proof. Differentiate the power series for $e^{t X}$ term by term. This is permitted because, for each $j$ and $k,\left(e^{t X}\right)_{j k}$ is given by a convergent power series in $t$, and one can differentiate a power series term by term inside its radius of convergence (e.g., Theorem 12 in Chapter 4 of [Pugh]).

\subsection*{Computing the Exponential}
We consider here methods for exponentiating general matrices. A special method for exponentiating $2 \times 2$ matrices is described in Exercises 6 and 7. Suppose that $X \in M_{n}(\mathbb{C})$ has $n$ linearly independent eigenvectors $v_{1}, \ldots, v_{n}$ with eigenvalues $\lambda_{1}, \ldots, \lambda_{n}$. Let $C$ be the $n \times n$ matrix whose columns are $v_{1}, \ldots, v_{n}$ and let $D$ be the diagonal matrix with diagonal entries $\lambda_{1}, \ldots, \lambda_{n}$. Then $X=C D C^{-1}$. It is easily verified that $e^{D}$ is the diagonal matrix with diagonal entries $e^{\lambda_{1}}, \ldots, e^{\lambda_{n}}$, and thus, by Proposition 2.3, we have

$$
e^{X}=C\left(\begin{array}{ccc}
e^{\lambda_{1}} & & 0 \\
& \ddots & \\
& & e^{\lambda_{n}}
\end{array}\right) C^{-1}
$$

Meanwhile, if $X$ is nilpotent (i.e., $X^{k}=0$ for some $k$ ), then the series that defines $e^{X}$ terminates. Finally, according to Theorem A.6, every matrix $X$ can be written (uniquely) in the form $X=S+N$, with $S$ diagonalizable, $N$ nilpotent, and $S N=N S$. Then, since $N$ and $S$ commute,

$$
e^{X}=e^{S+N}=e^{S} e^{N}
$$

Example 2.5. Consider the matrices

$$
X_{1}=\left(\begin{array}{rr}
0 & -a \\
a & 0
\end{array}\right) ; \quad X_{2}=\left(\begin{array}{ccc}
0 & a & b \\
0 & 0 & c \\
0 & 0 & 0
\end{array}\right) ; \quad X_{3}=\left(\begin{array}{cc}
a & b \\
0 & a
\end{array}\right)
$$

Then

$$
e^{X_{1}}=\left(\begin{array}{rr}
\cos a & -\sin a \\
\sin a & \cos a
\end{array}\right)
$$

and

$$
e^{X_{2}}=\left(\begin{array}{ccc}
1 & a & b+a c / 2 \\
0 & 1 & c \\
0 & 0 & 1
\end{array}\right)
$$

and

$$
e^{X_{3}}=\left(\begin{array}{cc}
e^{a} & e^{a} b \\
0 & e^{a}
\end{array}\right)
$$

Proof. The eigenvectors of $X_{1}$ are ( $1, i$ ) and ( $i, 1$ ), with eigenvalues $-i a$ and $i a$, respectively. Thus,

$$
e^{X_{1}}=\left(\begin{array}{cc}
1 & i \\
i & 1
\end{array}\right)\left(\begin{array}{cc}
e^{-i a} & 0 \\
0 & e^{i a}
\end{array}\right)\left(\begin{array}{rr}
1 / 2 & -i / 2 \\
-i / 2 & 1 / 2
\end{array}\right)
$$

which simplifies to the claimed result. Meanwhile, $X_{2}^{2}$ has the value $a c$ in the upper right-hand corner and all other entries equal to zero, whereas $X_{2}^{3}=0$. Thus, $e^{X_{2}}= 1+X_{2}+X_{2}^{2} / 2$, which reduces to the claimed result. Finally,

$$
X_{3}=\left(\begin{array}{cc}
a & 0 \\
0 & a
\end{array}\right)+\left(\begin{array}{ll}
0 & b \\
0 & 0
\end{array}\right)
$$

where the two terms clearly commute and the second term is nilpotent. Thus, we obtain

$$
e^{X_{3}}=\left(\begin{array}{cc}
e^{a} & 0 \\
0 & e^{a}
\end{array}\right)\left(\begin{array}{ll}
1 & b \\
0 & 1
\end{array}\right)
$$

which reduces to the claimed result.\\
The matrix exponential is used the elementary theory of differential equations, to solve systems of linear equations. Consider a first-order differential equation of the form

$$
\begin{aligned}
\frac{d \mathbf{v}}{d t} & =X \mathbf{v} \\
\mathbf{v}(0) & =\mathbf{v}_{0}
\end{aligned}
$$

where $\mathbf{v}(t) \in \mathbb{R}^{n}$ and $X$ is a fixed $n \times n$ matrix. The (unique) solution of this equation is given by

$$
\mathbf{v}(t)=e^{t X} \mathbf{v}_{0}
$$

as may be easily verified using Proposition 2.4. Curves of the form $t \mapsto e^{t X} \mathbf{v}_{0}$, with $\mathbf{v}_{0}$ fixed, trace out the flow along the vector field $\mathbf{v} \mapsto X \mathbf{v}$.

Let us consider the two matrices

$$
X_{4}=\left(\begin{array}{rr}
1 & 2 \\
-2 & 1
\end{array}\right) ; \quad X_{5}=\left(\begin{array}{cc}
1 & 2 \\
2 & 1
\end{array}\right)
$$

Figure 2.1 plots several curves of this form for each matrix (see Exercise 8).

\includegraphics[scale=0.2, center]{2025_08_28_fc86fd4ade69f41b5842g-049}

Fig. 2.1 Curves of the form $t \textbackslash mapsto e^{t X_{4}} \mathbf{v}_{0}$ (left) and $t \mapsto e^{t X_{5}} \mathbf{v}_{0}$ (right)


\subsection*{The Matrix Logarithm}
We now wish to define a matrix logarithm, which should be an inverse function (to the extent possible) to the matrix exponential. Let us recall the situation for the logarithm of complex numbers, in order to see what is reasonable to expect in the matrix case. Since $e^{z}$ is never zero, only nonzero numbers can have a logarithm. Every nonzero complex number can be written as $e^{z}$ for some $z$, but the $z$ is not unique and cannot be defined continuously on $\mathbb{C}^{*}$. In the matrix case, $e^{X}$ is invertible for all $X \in M_{n}(\mathbb{C})$. We will see (Theorem 2.10) that every invertible matrix can be written as $e^{X}$, for some $X \in M_{n}(\mathbb{C})$, but the $X$ is not unique.

The simplest way to define the matrix logarithm is by a power series. We recall how this works in the complex case.

Lemma 2.6. The function

$$
\log z=\sum_{m=1}^{\infty}(-1)^{m+1} \frac{(z-1)^{m}}{m}
$$

is defined and analytic in a circle of radius 1 about $z=1$.\\
For all $z$ with $|z-1|<1$,

$$
e^{\log z}=z .
$$

For all $u$ with $|u|<\log 2$, we have $\left|e^{u}-1\right|<1$ and

$$
\log e^{u}=u .
$$

Proof. The usual logarithm for real, positive numbers satisfies

$$
\frac{d}{d x} \log (1-x)=\frac{-1}{1-x}=-\left(1+x+x^{2}+\cdots\right)
$$

for $|x|<1$. Integrating term by term and noting that $\log 1=0$ gives

$$
\log (1-x)=-\left(x+\frac{x^{2}}{2}+\frac{x^{3}}{3}+\cdots\right)
$$

Taking $z=1-x$ (so that $x=1-z$ ), we have

$$
\begin{aligned}
\log z & =-\left((1-z)+\frac{(1-z)^{2}}{2}+\frac{(1-z)^{3}}{3}+\cdots\right) \\
& =\sum_{m=1}^{\infty}(-1)^{m+1} \frac{(z-1)^{m}}{m}
\end{aligned}
$$

The series (2.8) has radius of convergence 1 and defines a holomorphic function on the set $\{|z-1|<1\}$, which coincides with the usual logarithm for real $z$ in the interval ( 0,2 ). Now, $\exp (\log z)=z$ for $z \in(0,2)$ and since both sides of this identity are holomorphic in $z$, the identity continues to hold on the whole set $\{|z-1|<1\}$.

On the other hand, if $|u|<\log 2$, then

$$
\left|e^{u}-1\right|=\left|u+\frac{u^{2}}{2!}+\cdots\right| \leq|u|+\frac{|u|^{2}}{2!}+\cdots=e^{|u|}-1<1 .
$$

Thus, $\log (\exp u)$ makes sense for all such $u$. Since $\log (\exp u)=u$ for real $u$ with $|u|<\log 2$, it follows by holomorphicity that $\log (\exp u)=u$ for all complex numbers with $|u|<\log 2$.

Definition 2.7. For an $n \times n$ matrix $A$, define $\log A$ by

$$
\log A=\sum_{m=1}^{\infty}(-1)^{m+1} \frac{(A-I)^{m}}{m}
$$

whenever the series converges.\\
Since the complex-valued series (2.7) has radius of convergence 1 and since $\left\|(A-I)^{m}\right\| \leq\|A-I\|^{m}$ for $m \geq 1$, the matrix-valued series (2.9) will converge if $\|A-I\|<1$. Even if $\|A-I\|>1$, the series might converge, for example, if $A-I$ is nilpotent (see Exercise 9).

Theorem 2.8. The function

$$
\log A=\sum_{m=1}^{\infty}(-1)^{m+1} \frac{(A-I)^{m}}{m}
$$

is defined and continuous on the set of all $n \times n$ complex matrices $A$ with $\|A-I\|<1$.

For all $A$ with $\|A-I\|<1$,

$$
e^{\log A}=A
$$

For all $X$ with $\|X\|<\log 2,\left\|e^{X}-I\right\|<1$ and

$$
\log e^{X}=X
$$

Although it might seem plausible that $\log \left(e^{X}\right)$ should be equal to $X$ whenever the series for the logarithm is convergent, this claim is false (even over $\mathbb{C}$ ). If, for example, $X=2 \pi i I$, then $e^{X}=e^{2 \pi i} I=I$. Then $e^{X}-I=0$, so that $\log \left(e^{X}\right)$ is defined and equal to 0 . In this case, $\log \left(e^{X}\right)$ is defined but not equal to $X$. Thus, the assumption that $\|X\|<\log 2$ cannot be replaced by, say, the assumption that $\left\|e^{X}-I\right\|<1$.\\
Proof. Since $\left\|(A-I)^{m}\right\| \leq\|(A-I)\|^{m}$ and since the series (2.7) has radius of convergence 1, the series (2.9) converges absolutely for all $A$ with $\|A-I\|<1$. The proof of continuity is essentially the same as for the exponential.

Suppose now that $A$ satisfies $\|A-I\|<1$. If $A$ is diagonalizable with eigenvalue $z_{1}, \ldots, z_{n}$, then we can express $A$ in the form $C D C^{-1}$ with $D$ diagonal, in which case

$$
(A-I)^{m}=C\left(\begin{array}{ccc}
\left(z_{1}-1\right)^{m} & & 0 \\
& \ddots & \\
0 & & \left(z_{n}-1\right)^{m}
\end{array}\right) C^{-1} .
$$

Since $\|A-I\|<1$, each eigenvalue $z_{j}$ of $A$ must satisfy $\left|z_{j}-1\right|<1$ (Exercise 2). Thus,

$$
\sum_{m=1}^{\infty}(-1)^{m+1} \frac{(A-I)^{m}}{m}=C\left(\begin{array}{ccc}
\log z_{1} & & 0 \\
& \ddots & \\
& & \log z_{n}
\end{array}\right) C^{-1}
$$

and by Lemma 2.6,

$$
e^{\log A}=C\left(\begin{array}{ccc}
e^{\log z_{1}} & & 0 \\
& \ddots & \\
& & e^{\log z_{n}}
\end{array}\right) C^{-1}=A
$$

If $A$ is not diagonalizable, we approximate $A$ by a sequence $A_{m}$ of diagonalizable matrices (Exercise 4) and appeal to the continuity of the logarithm and exponential functions. Thus, $\exp (\log A)=A$ for all $A$ with $\|A-I\|<1$.

Now, the same argument as in the complex case shows that if $\|X\|<\log 2$, then $\left\|e^{X}-I\right\|<1$. The proof that $\log \left(e^{X}\right)=X$ is then very similar to the proof that $\exp (\log A)=A$.

Proposition 2.9. There exists a constant $c$ such that for all $n \times n$ matrices $B$ with $\|B\|<\frac{1}{2}$,

$$
\|\log (I+B)-B\| \leq c\|B\|^{2}
$$

Proof. Note that

$$
\log (I+B)-B=\sum_{m=2}^{\infty}(-1)^{m+1} \frac{B^{m}}{m}=B^{2} \sum_{m=2}^{\infty}(-1)^{m+1} \frac{B^{m-2}}{m}
$$

so that

$$
\|\log (I+B)-B\| \leq\|B\|^{2} \sum_{m=2}^{\infty} \frac{\left(\frac{1}{2}\right)^{m-2}}{m}
$$

which is an estimate of the desired form.\\
We may restate the proposition in a more concise way by saying that

$$
\log (I+B)=B+O\left(\|B\|^{2}\right)
$$

where $O\left(\|B\|^{2}\right)$ denotes a quantity of order $\|B\|^{2}$ (i.e., a quantity that is bounded by a constant times $\|B\|^{2}$ for all sufficiently small values of $\|B\|$ ).

We conclude this section with a result that, although we will not use it elsewhere, is worth recording. The proof is sketched in Exercises 9 and 10.\\
Theorem 2.10. Every invertible $n \times n$ matrix can be expressed as $e^{X}$ for some $X \in M_{n}(\mathbb{C})$.

\subsection*{Further Properties of the Exponential}
In this section, we give several additional results involving the exponential of a matrix that will be important in our study of Lie algebras.

Theorem 2.11 (Lie Product Formula). For all $X, Y \in M_{n}(\mathbb{C})$, we have

$$
e^{X+Y}=\lim _{m \rightarrow \infty}\left(e^{\frac{X}{m}} e^{\frac{Y}{m}}\right)^{m}
$$

There is a version of this result, known as the Trotter product formula, which holds for suitable unbounded operators on an infinite-dimensional Hilbert space. See, for example, Theorem 20.1 in [Hall].\\
Proof. If we multiply the power series for $e^{\frac{X}{m}}$ and $e^{\frac{Y}{m}}$, all but three of the terms will involve $1 / m^{2}$ or higher powers of $1 / m$. Thus,

$$
e^{\frac{X}{m}} e^{\frac{Y}{m}}=I+\frac{X}{m}+\frac{Y}{m}+O\left(\frac{1}{m^{2}}\right)
$$

Now, since $e^{\frac{X}{m}} e^{\frac{Y}{m}} \rightarrow I$ as $m \rightarrow \infty, e^{\frac{X}{m}} e^{\frac{Y}{m}}$ is in the domain of the logarithm for all sufficiently large $m$. By Proposition 2.9,

$$
\begin{aligned}
\log \left(e^{\frac{X}{m}} e^{\frac{Y}{m}}\right) & =\log \left(I+\frac{X}{m}+\frac{Y}{m}+O\left(\frac{1}{m^{2}}\right)\right) \\
& =\frac{X}{m}+\frac{Y}{m}+O\left(\left\|\frac{X}{m}+\frac{Y}{m}+O\left(\frac{1}{m^{2}}\right)\right\|^{2}\right) \\
& =\frac{X}{m}+\frac{Y}{m}+O\left(\frac{1}{m^{2}}\right)
\end{aligned}
$$

Exponentiating the logarithm then gives

$$
e^{\frac{X}{m}} e^{\frac{Y}{m}}=\exp \left(\frac{X}{m}+\frac{Y}{m}+O\left(\frac{1}{m^{2}}\right)\right)
$$

and, therefore,

$$
\left(e^{\frac{X}{m}} e^{\frac{Y}{m}}\right)^{m}=\exp \left(X+Y+O\left(\frac{1}{m}\right)\right)
$$

Thus, by the continuity of the exponential, we conclude that

$$
\lim _{m \rightarrow \infty}\left(e^{\frac{X}{m}} e^{\frac{Y}{m}}\right)^{m}=\exp (X+Y)
$$

which is the Lie product formula.

Recall (Sect. A.5) that the trace of matrix is defined as the sum of its diagonal entries and that similar matrices have the same trace.

Theorem 2.12. For any $X \in M_{n}(\mathbb{C})$, we have

$$
\operatorname{det}\left(e^{X}\right)=e^{\operatorname{trace}(X)}
$$

Proof. If $X$ is diagonalizable with eigenvalues $\lambda_{1}, \ldots, \lambda_{n}$, then $e^{X}$ is diagonalizable with eigenvalues $e^{\lambda_{1}}, \ldots, e^{\lambda_{n}}$. Thus, trace $(X)=\sum_{j} \lambda_{j}$ and

$$
\operatorname{det}\left(e^{X}\right)=e^{\lambda_{1}} \cdots e^{\lambda_{n}}=e^{\lambda_{1}+\cdots \lambda_{n}}=e^{\operatorname{trace}(X)} .
$$

If $X$ is not diagonalizable, we can approximate it by matrices that are diagonalizable (Exercise 4).

Definition 2.13. A function $A: \mathbb{R} \rightarrow \mathrm{GL}(n ; \mathbb{C})$ is called a one-parameter subgroup of $\mathrm{GL}(n ; \mathbb{C})$ if

\begin{enumerate}
  \item $A$ is continuous,
  \item $A(0)=I$,
  \item $A(t+s)=A(t) A(s)$ for all $t, s \in \mathbb{R}$.
\end{enumerate}

Theorem 2.14 (One-Parameter Subgroups). If $A(\cdot)$ is a one-parameter subgroup of $\mathrm{GL}(n ; \mathbb{C})$, there exists a unique $n \times n$ complex matrix $X$ such that

$$
A(t)=e^{t X}
$$

By taking $n=1$, and noting that $\mathrm{GL}(1 ; \mathbb{C}) \cong \mathbb{C}^{*}$, this theorem provides a method of solving Exercise 18 in Chapter 1.

Lemma 2.15. Fix some $\varepsilon$ with $\varepsilon<\log 2$. Let $B_{\varepsilon / 2}$ be the ball of radius $\varepsilon / 2$ around the origin in $M_{n}(\mathbb{C})$, and let $U=\exp \left(B_{\varepsilon / 2}\right)$. Then every $B \in U$ has a unique square root $C$ in $U$, given by $C=\exp \left(\frac{1}{2} \log B\right)$.

Proof. It is evident that $C$ is a square root of $B$ and that $C$ is in $U$. To establish uniqueness, suppose $C^{\prime} \in U$ satisfies $\left(C^{\prime}\right)^{2}=B$. Let $Y=\log C^{\prime}$; then $\exp (Y)= C^{\prime}$ and

$$
\exp (2 Y)=\left(C^{\prime}\right)^{2}=B=\exp (\log B)
$$

We have that $Y \in B_{\varepsilon / 2}$ and, thus, $2 Y \in B_{\varepsilon}$, and also that $\log B \in B_{\varepsilon / 2} \subset B_{\varepsilon}$. Since, by Theorem 2.8, $\exp$ is injective on $B_{\varepsilon}$ and $\exp (2 Y)=\exp (\log B)$, we must have $2 Y=\log B$. Thus, $C^{\prime}=\exp \left(\frac{1}{2} \log B\right)=C$.

Proof of Theorem 2.14. The uniqueness is immediate, since if there is such an $X$, then $X=\left.\frac{d}{d t} A(t)\right|_{t=0}$. To prove existence, let $U$ be as in Lemma 2.15, which is an open set in $\operatorname{GL}(n ; \mathbb{C})$. The continuity of $A$ guarantees that there exists $t_{0}>0$ such that $A(t) \in U$ for all $t$ with $|t| \leq t_{0}$. Define

$$
X=\frac{1}{t_{0}} \log \left(A\left(t_{0}\right)\right),
$$

so that $t_{0} X=\log \left(A\left(t_{0}\right)\right)$. Then $t_{0} X \in B_{\varepsilon / 2}$ and

$$
e^{t_{0} X}=A\left(t_{0}\right) .
$$

Now, $A\left(t_{0} / 2\right)$ is again in $U$ and $A\left(t_{0} / 2\right)^{2}=A\left(t_{0}\right)$. But by Lemma 2.15, $A\left(t_{0}\right)$ has a unique square root in $U$, and that unique square root is $\exp \left(t_{0} X / 2\right)$. Thus,

$$
A\left(t_{0} / 2\right)=\exp \left(t_{0} X / 2\right)
$$

Applying this argument repeatedly, we conclude that

$$
A\left(t_{0} / 2^{k}\right)=\exp \left(t_{0} X / 2^{k}\right)
$$

for all positive integers $k$. Then for any integer $m$, we have

$$
A\left(m t_{0} / 2^{k}\right)=A\left(t_{0} / 2^{k}\right)^{m}=\exp \left(m t_{0} X / 2^{k}\right) .
$$

It follows that $A(t)=\exp (t X)$ for all real numbers $t$ of the form $t=m t_{0} / 2^{k}$, and the set of such $t$ 's is dense in $\mathbb{R}$. Since both $\exp (t X)$ and $A(t)$ are continuous, it follows that $A(t)=\exp (t X)$ for all real numbers $t$.

Proposition 2.16. The exponential map is an infinitely differentiable map of $M_{n}(\mathbb{C})$ into $M_{n}(\mathbb{C})$.

We will compute the derivative of the matrix exponential in Chapter 5.\\
Proof. Note that for each $j$ and $k$, the quantity $\left(X^{m}\right)_{j k}$ is a homogeneous polynomial of degree $m$ in the entries of $X$. Thus, the series for the function $\left(X^{m}\right)_{j k}$ has the form of a multivariable power series on $M_{n}(\mathbb{C}) \cong \mathbb{R}^{2 n^{2}}$. Since the series converges on all of $\mathbb{R}^{2 n^{2}}$, it is permissible to differentiate the series term by term as many times as we like. (Apply Theorem 12 in Chapter 4 of [Pugh] in each of the $n^{2}$ variables with the other variables fixed.)

\subsection*{The Polar Decomposition}
The polar decomposition for a nonzero complex number $z$ states that $z$ can be written uniquely as $z=u p$, where $|u|=1$ and $p$ is real and positive. (If $z=0$, the decomposition still exists, with $p=0$, but $u$ is not unique.) Since $p$ is real and positive, it can be written as $p=e^{x}$ for a unique real number $x$. This gives an unconventional form of the polar decomposition for $z$, namely

$$
z=u e^{x}
$$

with $x \in \mathbb{R}$ and $|u|=1$. Although it is customary to leave $p$ as a positive real number and to write $u$ as $u=e^{i \theta}$, the decomposition in (2.10) is more convenient for us because $x$, unlike $\theta$, is unique.

We wish to establish a similar polar decomposition first for $\operatorname{GL}(n ; \mathbb{C})$ and then for various subgroups thereof. If $P$ is a self-adjoint $n \times n$ matrix (i.e., $P^{*}=P$ ), we say that $P$ is positive if $\langle v, P v\rangle>0$ for all nonzero $v \in \mathbb{C}^{n}$. It is easy to check that a self-adjoint matrix $P$ is positive if and only if all the eigenvalues of $P$ are positive. Suppose now that $A$ is an invertible $n \times n$ matrix. We wish to write $A$ as $A=U P$ where $U$ is unitary and $P$ is self-adjoint and positive. We will then write the selfadjoint, positive matrix $P$ as $P=e^{X}$ where $X$ is self-adjoint but not necessarily positive.

\section*{Theorem 2.17.}
\begin{enumerate}
  \item Every $A \in \mathrm{GL}(n ; \mathbb{C})$ can be written uniquely in the form
\end{enumerate}

$$
A=U P
$$

where $U$ is unitary and $P$ is self-adjoint and positive.\\
2. Every self-adjoint positive matrix $P$ can be written uniquely in the form

$$
P=e^{X}
$$

with $X$ self-adjoint. Conversely, if $X$ is self-adjoint, then $e^{X}$ is self-adjoint and positive.\\
3. If we decompose each $A \in \mathrm{GL}(n ; \mathbb{C})$ (uniquely) as

$$
A=U e^{X}
$$

with $U$ unitary and $X$ self-adjoint, then $U$ and $X$ depend continuously on $A$.\\
Lemma 2.18. If $Q$ is a self-adjoint, positive matrix, then $Q$ has a unique positive, self-adjoint square root.

Proof. Since $Q$ has an orthonormal basis of eigenvectors, $Q$ can be written as

$$
Q=U\left(\begin{array}{lll}
\lambda_{1} & & \\
& \ddots & \\
& & \lambda_{n}
\end{array}\right) U^{-1}
$$

with $U$ unitary. Since $Q$ is self-adjoint and positive, each $\lambda_{j}$ is positive. Thus, we can construct a square root of $Q$ as

$$
Q^{1 / 2}=U\left(\begin{array}{lll}
\lambda_{1}^{1 / 2} & & \\
& \ddots & \\
& & \lambda_{n}^{1 / 2}
\end{array}\right) U^{-1}
$$

and $Q^{1 / 2}$ will still be self-adjoint and positive, establishing the existence of the square root.

If $P$ is a self-adjoint, positive matrix, the eigenspaces of $P^{2}$ are precisely the same as the eigenspaces of $P$, with the eigenvalues of $P^{2}$ being, of course, the squares of the eigenvalues of $P$. The point here is that because the function $x \mapsto x^{2}$ is injective on positive real numbers, eigenspaces with distinct eigenvalues remain with distinct eigenvalues after squaring. Looking at this claim the other way around, if a positive, self-adjoint matrix $Q$ is to have a positive self-adjoint square root $P$, the eigenspaces of $P$ must be the same as the eigenspaces of $Q$, and the eigenvalues of $P$ must be the positive square roots of the eigenvalues of $Q$. Thus, $P$ is uniquely determined by $Q$.

Proof of Theorem 2.17. For the existence of the decomposition in Point 1, note that if $A=U P$, then $A^{*} A=P U^{*} U P=P^{2}$. Now, for any matrix $A$, the matrix $A^{*} A$ is self-adjoint. If, in addition, $A$ is invertible, then for all nonzero $v \in \mathbb{C}^{n}$, we have

$$
\left\langle v, A^{*} A v\right\rangle=\langle A v, A v\rangle>0
$$

showing that $A$ is positive. For all invertible $A$, then, let us define $P$ by

$$
P=\left(A^{*} A\right)^{1 / 2}
$$

where $(\cdot)^{1 / 2}$ is the unique positive square root of Lemma 2.18. We then define

$$
U=A P^{-1}=A\left[\left(A^{*} A\right)^{1 / 2}\right]^{-1} .
$$

Since $P$ is, by construction, self-adjoint and positive, and since $A=U P$ by the definition of $U$, it remains only to check that $U$ is unitary. To that end, we check that

$$
U^{*} U=\left[\left(A^{*} A\right)^{1 / 2}\right]^{-1} A^{*} A\left[\left(A^{*} A\right)^{1 / 2}\right]^{-1}
$$

since the inverse of a positive self-adjoint matrix is self-adjoint. Since $A^{*} A$ is the square of $\left(A^{*} A\right)^{1 / 2}$, we see that $U^{*} U=I$, showing that $U$ is unitary.

For the uniqueness of the decomposition, we have already noted that if $A=U P$, then $P^{2}=A^{*} A$, where $A^{*} A$ is self-adjoint and positive. Thus, the uniqueness of $P$ follows from the uniqueness in Lemma 2.18. The uniqueness of $U$ then follows, since if $A=U P$, then $U=A P^{-1}$.

The existence and uniqueness of the decomposition in Point 2 are proved in precisely the same way as in Lemma 2.18, with the logarithm function (which is a bijection between $(0, \infty)$ and $\mathbb{R}$ ) replacing the square root function. The same sort\\
of reasoning shows that for any self-adjoint $X$, the matrix $e^{X}$ is self-adjoint and positive.

Finally, we address the continuity claim in Point 3. From the formulas for $P$ and $U$ in terms of $A$ in the proof of Point 1 , we see that $U$ and $P$ depend continuously on $A$. It remains only to show that the logarithm $X$ of $P$ depends continuously on $P$. To see this, note that if the eigenvalues of $P$ are between 0 and 2 , then the power series for $\log P$ will converge to $X$, in which case, continuity follows by the same argument as in the proof of Proposition 2.1. In general, fix some positive, self-adjoint matrix $P_{0}$. Choose some large positive number $a$, and for $P$ in a small neighborhood $V$ of $P_{0}$, write $P=e^{a}\left(e^{-a} P\right)$. Then $P=e^{X}$, where

$$
X=a I+\log \left(e^{-a} P\right)
$$

Since $a$ is large, the eigenvalues of $e^{-a} P$ will all be less than $\log 2$, and the series for $\log \left(e^{-a} P\right)$ will converge and depend continuously on $P$.

We now establish polar decompositions for $\mathrm{GL}(n ; \mathbb{R}), \operatorname{SL}(n ; \mathbb{C})$, and $\operatorname{SL}(n ; \mathbb{R})$.

\section*{Proposition 2.19.}
\begin{enumerate}
  \item Every $A \in \mathrm{GL}(n ; \mathbb{R})$ can be written uniquely as
\end{enumerate}

$$
A=\operatorname{Re}^{X}
$$

where $R$ is in $\mathrm{O}(n)$ and $X$ is real and symmetric.\\
2. Every $A \in \mathrm{SL}(n ; \mathbb{C})$ can be written uniquely as

$$
A=U e^{X}
$$

where $U$ is in $\mathrm{SU}(n)$ and $X$ is self-adjoint with trace zero.\\
3. Every $A \in \mathrm{SL}(n ; \mathbb{R})$ can be written uniquely as

$$
A=R e^{X}
$$

where $R$ is in $\mathrm{SO}(n)$ and $X$ is real and symmetric and has trace zero.\\
Proof. If $A$ is real, then $A^{*} A$ is real and symmetric. Now, a real, symmetric matrix can be diagonalized over $\mathbb{R}$. Thus, $P$, which is the unique self-adjoint positive square root of $A^{*} A$ (constructed as in (2.11)), is real. Then $U=A P^{-1}$ is real and unitary, hence in $\mathrm{O}(n)$.

Meanwhile, if $A \in \mathrm{SL}(n ; \mathbb{C})$ and we write $A=U e^{X}$ with $U \in \mathrm{U}(n)$ and $X$ selfadjoint, then $\operatorname{det}(A)=\operatorname{det}(U) e^{\operatorname{trace}(X)}$. Now, $|\operatorname{det}(U)|=1$, and $e^{\operatorname{trace}(X)}$ is real and positive. Thus, by the uniqueness of the polar decomposition for nonzero complex numbers, we must have $\operatorname{det}(U)=1$ and $\operatorname{trace}(X)=0$. The case of $A \in \operatorname{SL}(n ; \mathbb{R})$ follows by combining the arguments in the two previous cases.

\subsection*{Exercises}
\begin{enumerate}
  \item The Cauchy-Schwarz inequality from elementary analysis tells us that for all $u=\left(u_{1}, \ldots, u_{n}\right)$ and $v=\left(v_{1}, \ldots, v_{n}\right)$ in $\mathbb{C}^{n}$, we have
\end{enumerate}

$$
\left|u_{1} v_{1}+\cdots+u_{n} v_{n}\right|^{2} \leq\left(\sum_{j=1}^{n}\left|u_{j}\right|^{2}\right)\left(\sum_{k=1}^{n}\left|v_{k}\right|^{2}\right) .
$$

Use this to verify that $\|X Y\| \leq\|X\|\|Y\|$ for all $X, Y \in M_{n}(\mathbb{C})$, where $\|\cdot\|$ is the Hilbert-Schmidt norm in Definition 2.2.\\
2. Show that for $X \in M_{n}(\mathbb{C})$ and any orthonormal basis $\left\{u_{1}, \ldots, u_{n}\right\}$ of $\mathbb{C}^{n}$, $\|X\|^{2}=\sum_{j, k=1}^{n}\left|\left\langle u_{j}, X u_{k}\right\rangle\right|^{2}$, where $\|X\|$ is as in Definition 2.2. Now show that if $v$ is an eigenvector for $X$ with eigenvalue $\lambda$, then $|\lambda| \leq\|X\|$.\\
3. The product rule. Recall that a matrix-valued function $A(t)$ is said to be smooth if each $A_{j k}(t)$ is smooth. The derivative of such a function is defined as

$$
\left(\frac{d A}{d t}\right)_{j k}=\frac{d A_{j k}}{d t}
$$

or, equivalently,

$$
\frac{d}{d t} A(t)=\lim _{h \rightarrow 0} \frac{A(t+h)-A(t)}{h} .
$$

Let $A(t)$ and $B(t)$ be two such functions. Prove that $A(t) B(t)$ is again smooth and that

$$
\frac{d}{d t}[A(t) B(t)]=\frac{d A}{d t} B(t)+A(t) \frac{d B}{d t}
$$

\begin{enumerate}
  \setcounter{enumi}{3}
  \item Using Theorem A.4, show that every $n \times n$ complex matrix $A$ is the limit of a sequence of diagonalizable matrices.\\
Hint: If an $n \times n$ matrix has $n$ distinct eigenvalues, it is necessarily diagonalizable.
  \item For any $a$ and $d$ in $\mathbb{C}$, define the expression $\left(e^{a}-e^{d}\right) /(a-d)$ in the obvious way for $a \neq d$ and by means of the limit
\end{enumerate}

$$
\lim _{a \rightarrow d} \frac{e^{a}-e^{d}}{a-d}=e^{a}
$$

when $a=d$. Show that for any $a, b, c \in \mathbb{C}$, we have

$$
\exp \left(\begin{array}{ll}
a & b \\
0 & d
\end{array}\right)=\left(\begin{array}{cc}
e^{a} & b \frac{e^{a}-e^{d}}{a-d} \\
0 & e^{d}
\end{array}\right)
$$

Hint: Show that if $a \neq d$, then

$$
\left(\begin{array}{ll}
a & b \\
0 & d
\end{array}\right)^{m}=\left(\begin{array}{cc}
a^{m} & b \frac{a^{m}-d^{m}}{a-d} \\
0 & b^{m}
\end{array}\right)
$$

for every positive integer $m$.\\
6. Show that every $2 \times 2$ matrix $X$ with trace $(X)=0$ satisfies

$$
X^{2}=-\operatorname{det}(X) I .
$$

If $X$ is $2 \times 2$ with trace zero, show by direct calculation using the power series for the exponential that

$$
e^{X}=\cos (\sqrt{\operatorname{det} X}) I+\frac{\sin \sqrt{\operatorname{det} X}}{\sqrt{\operatorname{det} X}} X
$$

where $\sqrt{\operatorname{det} X}$ is either of the two (possibly complex) square roots of $\operatorname{det} X$. Use this to give an alternative computation of the exponential $e^{X_{1}}$ in Example 2.5. Note: The value of the coefficient of $X$ in (2.12) is to be interpreted as 1 when $\operatorname{det} X=0$, in accordance with the limit $\lim _{\theta \rightarrow 0} \sin \theta / \theta=1$.\\
7. Use the result of Exercise 6 to compute the exponential of the matrix

$$
X=\left(\begin{array}{rr}
4 & 3 \\
-1 & 2
\end{array}\right)
$$

Hint: Reduce the calculation to the trace-zero case.\\
8. Consider the two matrices $X_{4}$ and $X_{5}$ in (2.6). Compute $e^{t X_{4}}$ and $e^{t X_{5}}$ either by diagonalization or by the method in Exercises 6 and 7 . Show that curves of the form $t \mapsto e^{t X_{4}} \mathbf{v}_{0}$, with $\mathbf{v}_{0} \neq 0$, spiral out to infinity. Show that for $\mathbf{v}_{0}$ outside of a certain one-dimensional subspace of $\mathbb{R}^{2}$, curves of the form $t \mapsto e^{t X_{5}} \mathbf{v}_{0}$ tend to infinity in the direction of $(1,1)$ or the direction of $(-1,-1)$.\\
9. A matrix $A$ is said to be unipotent if $A-I$ is nilpotent (i.e., if $A$ is of the form $A=I+N$, with $N$ nilpotent). Note that $\log A$ is defined whenever $A$ is unipotent, because the series in Definition 2.7 terminates.\\
(a) Show that if $A$ is unipotent, then $\log A$ is nilpotent.\\
(b) Show that if $X$ is nilpotent, then $e^{X}$ is unipotent.\\
(c) Show that if $A$ is unipotent, then $\exp (\log A)=A$ and that if $X$ is nilpotent, then $\log (\exp X)=X$.

Hint: Let $A(t)=I+t(A-I)$. Show that $\exp (\log (A(t)))$ depends polynomially on $t$ and that $\exp (\log (A(t)))=A(t)$ for all sufficiently small $t$.\\
10. Show that every invertible $n \times n$ matrix $A$ can be written as $A=e^{X}$ for some $X \in M_{n}(\mathbb{C})$.\\
Hint: Theorem A. 5 implies that $A$ is similar to a block-diagonal matrix in which each block is of the form $\lambda I+N_{\lambda}$, with $N_{\lambda}$ being nilpotent. Use this result and Exercise 9.\\
11. Show that for all $X \in M_{n}(\mathbb{C})$, we have

$$
\lim _{m \rightarrow \infty}\left[I+\frac{X}{m}\right]^{m}=e^{X}
$$

Hint: Use the matrix logarithm.

\section*{Chapter 3 Lie Algebras}
\subsection*{Definitions and First Examples}
We now introduce the "abstract" notion of a Lie algebra. In Sect. 3.3, we will associate to each matrix Lie group a Lie algebra. It is customary to use lowercase Gothic (Fraktur) characters such as $\mathfrak{g}$ and $\mathfrak{h}$ to refer to Lie algebras.

Definition 3.1. A finite-dimensional real or complex Lie algebra is a finitedimensional real or complex vector space $\mathfrak{g}$, together with a map $[\cdot, \cdot]$ from $\mathfrak{g} \times \mathfrak{g}$ into $\mathfrak{g}$, with the following properties:

\begin{enumerate}
  \item $[\cdot, \cdot]$ is bilinear.
  \item $[\cdot, \cdot]$ is skew symmetric: $[X, Y]=-[Y, X]$ for all $X, Y \in \mathfrak{g}$.
  \item The Jacobi identity holds:
\end{enumerate}

$$
[X,[Y, Z]]+[Y,[Z, X]]+[Z,[X, Y]]=0
$$

for all $X, Y, Z \in \mathfrak{g}$.\\
Two elements $X$ and $Y$ of a Lie algebra $\mathfrak{g}$ commute if $[X, Y]=0$. A Lie algebra $\mathfrak{g}$ is commutative if $[X, Y]=0$ for all $X, Y \in \mathfrak{g}$.

The map $[\cdot, \cdot]$ is referred to as the bracket operation on $\mathfrak{g}$. Note also that Condition 2 implies that $[X, X]=0$ for all $X \in \mathfrak{g}$. The bracket operation on a Lie algebra is not, in general associative; nevertheless, the Jacobi identity can be viewed as a substitute for associativity.

Example 3.2. Let $\mathfrak{g}=\mathbb{R}^{3}$ and let $[\cdot, \cdot]: \mathbb{R}^{3} \times \mathbb{R}^{3} \rightarrow \mathbb{R}^{3}$ be given by

$$
[x, y]=x \times y,
$$

where $x \times y$ is the cross product (or vector product). Then $\mathfrak{g}$ is a Lie algebra.

Proof. Bilinearity and skew symmetry are standard properties of the cross product. To verify the Jacobi identity, it suffices (by bilinearity) to verify it when $x=e_{j}$, $y=e_{k}$, and $z=e_{l}$, where $e_{1}, e_{2}$, and $e_{3}$ are the standard basis elements for $\mathbb{R}^{3}$. If $j, k$, and $l$ are all equal, each term in the Jacobi identity is zero. If $j, k$, and $l$ are all different, the cross product of any two of $e_{j}, e_{k}$, and $e_{l}$ is equal to a multiple of the third, so again, each term in the Jacobi identity is zero. It remains to consider the case in which two of $j, k, l$ are equal and the third is different. By re-ordering the terms in the Jacobi identity as necessary, it suffices to verify the identity

$$
\left[e_{j},\left[e_{j}, e_{k}\right]\right]+\left[e_{j},\left[e_{k}, e_{j}\right]\right]+\left[e_{k},\left[e_{j}, e_{j}\right]\right] .
$$

The first two terms in (3.1) are negatives of each other and the third is zero.\\
Example 3.3. Let $\mathcal{A}$ be an associative algebra and let $\mathfrak{g}$ be a subspace of $\mathcal{A}$ such that $X Y-Y X \in \mathfrak{g}$ for all $X, Y \in \mathfrak{g}$. Then $\mathfrak{g}$ is a Lie algebra with bracket operation given by

$$
[X, Y]=X Y-Y X
$$

Proof. The bilinearity and skew symmetry of the bracket are evident. To verify the Jacobi identity, note that each double bracket generates four terms, for a total of 12 terms. It is left to the reader to verify that the product of $X, Y$, and $Z$ in each of the six possible orderings occurs twice, once with a plus sign and once with a minus sign.

If we look carefully at the proof of the Jacobi identity, we see that the two occurrences of, say, $X Y Z$ occur with different groupings, once as $X(Y Z)$ and once as $(X Y) Z$. Thus, associativity of the algebra $\mathcal{A}$ is essential. For any Lie algebra, the Jacobi identity means that the bracket operation behaves as if it were $X Y-Y X$ in some associative algebra, even if it is not actually defined this way. Indeed, we will prove in Chapter 9 that every Lie algebra $\mathfrak{g}$ can be embedded into an associative algebra $\mathcal{A}$ in such a way that the bracket becomes $X Y-Y X$. (This claim follows from Theorem 6.7, the Poincaré-Birkhoff-Witt theorem.)

Of particular interest to us is the case in which $\mathcal{A}$ is the space $M_{n}(\mathbb{C})$ of all $n \times n$ complex matrices.

Example 3.4. Let $\operatorname{sI}(n ; \mathbb{C})$ the space $X \in M_{n}(\mathbb{C})$ for which trace $(X)=0$. Then $\operatorname{sl}(n ; \mathbb{C})$ is a Lie algebra with bracket $[X, Y]=X Y-Y X$.

Proof. For any $X$ and $Y$ in $M_{n}(\mathbb{C})$, we have

$$
\operatorname{trace}(X Y-Y X)=\operatorname{trace}(X Y)-\operatorname{trace}(Y X)=0 .
$$

This holds, in particular, if $X$ and $Y$ have trace zero. Thus, Example 3.3 applies.\\
Definition 3.5. A subalgebra of a real or complex Lie algebra $\mathfrak{g}$ is a subspace $\mathfrak{h}$ of $\mathfrak{g}$ such that $\left[H_{1}, H_{2}\right] \in \mathfrak{h}$ for all $H_{1}$ and $H_{2} \in \mathfrak{h}$. If $\mathfrak{g}$ is a complex Lie algebra and $\mathfrak{h}$ is a real subspace of $\mathfrak{g}$ which is closed under brackets, then $\mathfrak{h}$ is said to be a real subalgebra of $\mathfrak{g}$.

A subalgebra $\mathfrak{h}$ of a Lie algebra $\mathfrak{g}$ is said to be an ideal in $\mathfrak{g}$ if $[X, H] \in \mathfrak{h}$ for all $X$ in $\mathfrak{g}$ and $H$ in $\mathfrak{h}$.

The center of a Lie algebra $\mathfrak{g}$ is the set of all $X \in \mathfrak{g}$ for which $[X, Y]=0$ for all $Y \in \mathfrak{g}$.

Definition 3.6. If $\mathfrak{g}$ and $\mathfrak{h}$ are Lie algebras, then a linear map $\phi: \mathfrak{g} \rightarrow \mathfrak{h}$ is called a Lie algebra homomorphism if $\phi([X, Y])=[\phi(X), \phi(Y)]$ for all $X, Y \in \mathfrak{g}$. If, in addition, $\phi$ is one-to-one and onto, then $\phi$ is called a Lie algebra isomorphism. A Lie algebra isomorphism of a Lie algebra with itself is called a Lie algebra automorphism.

Definition 3.7. If $\mathfrak{g}$ is a Lie algebra and $X$ is an element of $\mathfrak{g}$, define a linear map $\operatorname{ad}_{X}: \mathfrak{g} \rightarrow \mathfrak{g}$ by

$$
\operatorname{ad}_{X}(Y)=[X, Y] .
$$

The map $X \mapsto \operatorname{ad}_{X}$ is the adjoint map or adjoint representation.\\
Although $\operatorname{ad}_{X}(Y)$ is just $[X, Y]$, the alternative "ad" notation can be useful. For example, instead of writing

$$
[X,[X,[X,[X, Y]]]],
$$

we can now write

$$
\left(\operatorname{ad}_{X}\right)^{4}(Y) .
$$

This sort of notation will be essential in Chapter 5. We can view ad (that is, the map $\left.X \mapsto \mathrm{ad}_{X}\right)$ as a linear map of $\mathfrak{g}$ into $\operatorname{End}(\mathfrak{g})$, the space of linear operators on $\mathfrak{g}$. The Jacobi identity is then equivalent to the assertion that $\mathrm{ad}_{X}$ is a derivation of the bracket:

$$
\operatorname{ad}_{X}([Y, Z])=\left[\operatorname{ad}_{X}(Y), Z\right]+\left[Y, \operatorname{ad}_{X}(Z)\right] .
$$

Proposition 3.8. If $\mathfrak{g}$ is a Lie algebra, then

$$
\operatorname{ad}_{[X, Y]}=\operatorname{ad}_{X} \operatorname{ad}_{Y}-\operatorname{ad}_{Y} \operatorname{ad}_{X}=\left[\operatorname{ad}_{X}, \operatorname{ad}_{Y}\right] ;
$$

that is, ad: $\mathfrak{g} \rightarrow \operatorname{End}(\mathfrak{g})$ is a Lie algebra homomorphism.\\
Proof. Observe that

$$
\operatorname{ad}_{[X, Y]}(Z)=[[X, Y], Z],
$$

whereas

$$
\left[\operatorname{ad}_{X}, \operatorname{ad}_{Y}\right](Z)=[X,[Y, Z]]-[Y,[X, Z]] .
$$

Thus, we want to show that

$$
[[X, Y], Z]=[X,[Y, Z]]-[Y,[X, Z]]
$$

which is equivalent to the Jacobi identity.\\
Definition 3.9. If $\mathfrak{g}_{1}$ and $\mathfrak{g}_{2}$ are Lie algebras, the direct sum of $\mathfrak{g}_{1}$ and $\mathfrak{g}_{2}$ is the vector space direct sum of $\mathfrak{g}_{1}$ and $\mathfrak{g}_{2}$, with bracket given by

$$
\left[\left(X_{1}, X_{2}\right),\left(Y_{1}, Y_{2}\right)\right]=\left(\left[X_{1}, Y_{1}\right],\left[X_{2}, Y_{2}\right]\right) .
$$

If $\mathfrak{g}$ is a Lie algebra and $\mathfrak{g}_{1}$ and $\mathfrak{g}_{2}$ are subalgebras, we say that $\mathfrak{g}$ decomposes as the Lie algebra direct sum of $\mathfrak{g}_{1}$ and $\mathfrak{g}_{2}$ if $\mathfrak{g}$ is the direct sum of $\mathfrak{g}_{1}$ and $\mathfrak{g}_{2}$ as vector spaces and $\left[X_{1}, X_{2}\right]=0$ for all $X_{1} \in \mathfrak{g}_{1}$ and $X_{2} \in \mathfrak{g}_{2}$.

It is straightforward to verify that the bracket in (3.3) makes $\mathfrak{g}_{1} \oplus \mathfrak{g}_{2}$ into a Lie algebra. If $\mathfrak{g}$ decomposes as a Lie algebra direct sum of subalgebras $\mathfrak{g}_{1}$ and $\mathfrak{g}_{2}$, it is easy to check that $\mathfrak{g}$ is isomorphic as a Lie algebra to the "abstract" direct sum of $\mathfrak{g}_{1}$ and $\mathfrak{g}_{2}$. (This would not be the case without the assumption that every element of $\mathfrak{g}_{1}$ commutes with every element of $\mathfrak{g}_{2}$.)

Definition 3.10. Let $\mathfrak{g}$ be a finite-dimensional real or complex Lie algebra, and let $X_{1}, \ldots, X_{N}$ be a basis for $\mathfrak{g}$ (as a vector space). Then the unique constants $c_{j k l}$ such that

$$
\left[X_{j}, X_{k}\right]=\sum_{l=1}^{N} c_{j k l} X_{l}
$$

are called the structure constants of $\mathfrak{g}$ (with respect to the chosen basis).\\
Although we will not have much occasion to use them, structure constants do appear frequently in the physics literature. The structure constants satisfy the following two conditions:

$$
\begin{aligned}
c_{j k l}+c_{k j l} & =0 \\
\sum_{n}\left(c_{j k n} c_{n l m}+c_{k l n} c_{n j m}+c_{l n} c_{n k m}\right) & =0
\end{aligned}
$$

for all $j, k, l, m$. The first of these conditions comes from the skew symmetry of the bracket, and the second comes from the Jacobi identity.

\subsection*{Simple, Solvable, and Nilpotent Lie Algebras}
In this section, we consider various special types of Lie algebras. Recall from Definition 3.5 the notion of an ideal in a Lie algebra.

Definition 3.11. A Lie algebra $\mathfrak{g}$ is called irreducible if the only ideals in $\mathfrak{g}$ are $\mathfrak{g}$ and $\{0\}$. A Lie algebra $\mathfrak{g}$ is called simple if it is irreducible and $\operatorname{dim} \mathfrak{g} \geq 2$.

A one-dimensional Lie algebra is certainly irreducible, since it is has no nontrivial subspaces and therefore no nontrivial subalgebras and no nontrivial ideals. Nevertheless, such a Lie algebra is, by definition, not considered simple.

Note that a one-dimensional Lie algebra $\mathfrak{g}$ is necessarily commutative, since $[a X, b X]=0$ for any $X \in \mathfrak{g}$ and any scalars $a$ and $b$. On the other hand, if $\mathfrak{g}$ is commutative, then any subspace of $\mathfrak{g}$ is an ideal. Thus, the only way a commutative Lie algebra can be irreducible is if it is one dimensional. Thus, an equivalent definition of "simple" is that a Lie algebra is simple if it is irreducible and noncommutative.

There is an analogy between groups and Lie algebras, in which the role of subgroups is played by subalgebras and the role of normal subgroups is played by ideals. (For example, the kernel of a Lie algebra homomorphism is always an ideal, just as the kernel of a Lie group homomorphism is always a normal subgroup.) There is, however, an inconsistency in the terminology in the two fields. On the group side, any group with no nontrivial normal subgroups is called simple, including the most obvious example, a cyclic group of prime order. On the Lie algebra side, by contrast, the most obvious example of an algebra with no nontrivial ideals-namely, a onedimensional algebra-is not called simple.

We will eventually see many examples of simple Lie algebras, but for now we content ourselves with a single example. Recall the Lie algebra $\operatorname{sI}(n ; \mathbb{C})$ in Example 3.4.

Proposition 3.12. The Lie algebra $\mathrm{SI}(2 ; \mathbb{C})$ is simple.\\
Proof. We use the following basis for $\operatorname{sI}(2 ; \mathbb{C})$ :

$$
X=\left(\begin{array}{ll}
0 & 1 \\
0 & 0
\end{array}\right) ; \quad Y=\left(\begin{array}{ll}
0 & 0 \\
1 & 0
\end{array}\right) ; \quad H=\left(\begin{array}{rr}
1 & 0 \\
0 & -1
\end{array}\right)
$$

Direct calculation shows that these basis elements have the following commutation relations: $[X, Y]=H,[H, X]=2 X$, and $[H, Y]=-2 Y$. Suppose $\mathfrak{h}$ is an ideal in $\mathrm{sl}(2 ; \mathbb{C})$ and that $\mathfrak{h}$ contains an element $Z=a X+b H+c Y$, where $a, b$, and $c$ are not all zero. We will show, then, that $\mathfrak{h}=\mathrm{sl}(2 ; \mathbb{C})$. Suppose first that $c \neq 0$. Then the element

$$
[X,[X, Z]]=[X,[-2 b X+c H]]=-2 c X
$$

is a nonzero multiple of $X$. Since $\mathfrak{h}$ is an ideal, we conclude that $X \in \mathfrak{h}$. But $[Y, X]$ is a nonzero multiple of $H$ and $[Y,[Y, X]]$ is a nonzero multiple of $Y$, showing that $Y$ and $H$ also belong to $\mathfrak{h}$, from which we conclude that $\mathfrak{h}=\mathrm{sl}(2 ; \mathbb{C})$.

Suppose next that $c=0$ but $b \neq 0$. Then $[X, Z]$ is a nonzero multiple of $X$ and we may then apply the same argument in the previous paragraph to show that $\mathfrak{h}=\mathrm{sl}(2 ; \mathbb{C})$. Finally, if $c=0$ and $b=0$ but $a \neq 0$, then $Z$ itself is a nonzero multiple of $X$ and we again conclude that $\mathfrak{h}=\mathrm{sl}(2 ; \mathbb{C})$.

Definition 3.13. If $\mathfrak{g}$ is a Lie algebra, then the commutator ideal in $\mathfrak{g}$, denoted $[\mathfrak{g}, \mathfrak{g}]$, is the space of linear combinations of commutators, that is, the space of elements $Z$ in $\mathfrak{g}$ that can be expressed as

$$
Z=c_{1}\left[X_{1}, Y_{1}\right]+\cdots+c_{m}\left[X_{m}, Y_{m}\right]
$$

for some constants $c_{j}$ and vectors $X_{j}, Y_{j} \in \mathfrak{g}$.\\
For any $X$ and $Y$ in $\mathfrak{g}$, the commutator $[X, Y]$ is in $[\mathfrak{g}, \mathfrak{g}]$. This holds, in particular, if $X$ is in $[\mathfrak{g}, \mathfrak{g}]$, showing that $[\mathfrak{g}, \mathfrak{g}]$ is an ideal in $\mathfrak{g}$.

Definition 3.14. For any Lie algebra $\mathfrak{g}$, we define a sequence of subalgebras $\mathfrak{g}_{0}, \mathfrak{g}_{1}, \mathfrak{g}_{2}, \ldots$ of $\mathfrak{g}$ inductively as follows: $\mathfrak{g}_{0}=\mathfrak{g}, \mathfrak{g}_{1}=\left[\mathfrak{g}_{0}, \mathfrak{g}_{0}\right], \mathfrak{g}_{2}=\left[\mathfrak{g}_{1}, \mathfrak{g}_{1}\right]$, etc. These subalgebras are called the derived series of $\mathfrak{g}$. A Lie algebra $\mathfrak{g}$ is called solvable if $\mathfrak{g}_{j}=\{0\}$ for some $j$.

In light of the comments following Definition 3.13, each derived algebra $\mathfrak{g}_{j}$ is an ideal in $\mathfrak{g}_{j-1}$, but not necessarily an ideal in $\mathfrak{g}$.

Definition 3.15. For any Lie algebra $\mathfrak{g}$, we define a sequence of ideals $\mathfrak{g}^{j}$ in $\mathfrak{g}$ inductively as follows. We set $\mathfrak{g}^{0}=\mathfrak{g}$ and then define $\mathfrak{g}^{j+1}$ to be the space of linear combinations of commutators of the form $[X, Y]$ with $X \in \mathfrak{g}$ and $Y \in \mathfrak{g}^{j}$. These algebras are called the upper central series of $\mathfrak{g}$. A Lie algebra $\mathfrak{g}$ is said to be nilpotent if $\mathfrak{g}^{j}=\{0\}$ for some $j$.

Equivalently, $\mathfrak{g}^{j}$ is the space spanned by all $j$ th-order commutators,

$$
\left[X_{1},\left[X_{2},\left[X_{3}, \ldots\left[X_{j}, X_{j+1}\right] \ldots\right]\right]\right] .
$$

Note that every $j$ th-order commutator is also a ( $j-1$ )th-order commutator, by setting $\tilde{X}_{j}=\left[X_{j}, X_{j+1}\right]$. Thus, $\mathfrak{g}^{j-1} \subset \mathfrak{g}^{j}$. For every $X \in \mathfrak{g}$ and $Y \in \mathfrak{g}^{j}$, we have $[X, Y] \in \mathfrak{g}^{j+1} \subset \mathfrak{g}^{j}$, showing that $\mathfrak{g}^{j}$ is an ideal in $\mathfrak{g}$. Furthermore, it is clear that $\mathfrak{g}_{j} \subset \mathfrak{g}^{j}$ for all $j$; thus, if $\mathfrak{g}$ is nilpotent, $\mathfrak{g}$ is also solvable.

Proposition 3.16. If $\mathfrak{g} \subset M_{3}(\mathbb{R})$ denotes the space of $3 \times 3$ upper triangular matrices with zeros on the diagonal, then $\mathfrak{g}$ satisfies the assumptions of Example 3.3. The Lie algebra $\mathfrak{g}$ is a nilpotent Lie algebra.

Proof. We will use the following basis for $\mathfrak{g}$,

$$
X=\left(\begin{array}{ccc}
0 & 1 & 0 \\
0 & 0 & 0 \\
0 & 0 & 0
\end{array}\right) ; \quad Y=\left(\begin{array}{ccc}
0 & 0 & 0 \\
0 & 0 & 1 \\
0 & 0 & 0
\end{array}\right) ; \quad Z=\left(\begin{array}{ccc}
0 & 0 & 1 \\
0 & 0 & 0 \\
0 & 0 & 0
\end{array}\right)
$$

Direct calculation then establishes the following commutation relations: $[X, Y]= Z$ and $[X, Z]=[Y, Z]=0$. In particular, the bracket of two elements of $\mathfrak{g}$ is again in $\mathfrak{g}$, so that $\mathfrak{g}$ is a Lie algebra. Then $[\mathfrak{g}, \mathfrak{g}]$ is the span of $Z$ and $[\mathfrak{g},[\mathfrak{g}, \mathfrak{g}]]=0$, showing that $\mathfrak{g}$ is nilpotent.

Proposition 3.17. If $\mathfrak{g} \subset M_{2}(\mathbb{C})$ denotes the space of $2 \times 2$ matrices of the form

$$
\left(\begin{array}{ll}
a & b \\
0 & c
\end{array}\right)
$$

with $a, b$, and $c$ in $\mathbb{C}$, then $\mathfrak{g}$ satisfies the assumptions of Example 3.3. The Lie algebra $\mathfrak{g}$ is solvable but not nilpotent.

Proof. Direct calculation shows that

$$
\left[\left(\begin{array}{ll}
a & b \\
0 & c
\end{array}\right),\left(\begin{array}{ll}
d & e \\
0 & f
\end{array}\right)\right]=\left(\begin{array}{ll}
0 & h \\
0 & 0
\end{array}\right)
$$

where $h=a e+b f-b d-c e$, showing that $\mathfrak{g}$ is a Lie subalgebra of $M_{2}(\mathbb{C})$. Furthermore, the commutator ideal $[\mathfrak{g}, \mathfrak{g}]$ is one dimensional and hence commutative. Thus, $\mathfrak{g}_{2}=\{0\}$, showing that $\mathfrak{g}$ is solvable. On the other hand, consider the following elements of $\mathfrak{g}$ :

$$
H=\left(\begin{array}{rr}
1 & 0 \\
0 & -1
\end{array}\right) ; \quad X=\left(\begin{array}{ll}
0 & 1 \\
0 & 0
\end{array}\right)
$$

Using (3.6), we can see that $[H, X]=2 X$, and thus that

$$
[H,[H,[H, \cdots[H, X] \cdots]]]
$$

is a nonzero multiple of $X$, showing that $\mathfrak{g}^{j} \neq\{0\}$ for all $j$.

\subsection*{The Lie Algebra of a Matrix Lie Group}
In this section, we associate to each matrix Lie group $G$ a Lie algebra $\mathfrak{g}$. Many questions involving a group can be studied by transferring them to the Lie algebra, where we can use tools of linear algebra. We begin by defining $\mathfrak{g}$ as a set, and then proceed to give $\mathfrak{g}$ the structure of a Lie algebra.

Definition 3.18. Let $G$ be a matrix Lie group. The Lie algebra of $G$, denoted $\mathfrak{g}$, is the set of all matrices $X$ such that $e^{t X}$ is in $G$ for all real numbers $t$.

Equivalently, $X$ is in $\mathfrak{g}$ if and only if the entire one-parameter subgroup (Definition 2.13) generated by $X$ lies in $G$. Note that merely having $e^{X}$ in $G$ does not guarantee that $X$ is in $\mathfrak{g}$. Even though $G$ is a subgroup of $\operatorname{GL}(n ; \mathbb{C})$ (and not necessarily of $\mathrm{GL}(n ; \mathbb{R})$ ), we do not require that $e^{t X}$ be in $G$ for all complex numbers $t$, but only for all real numbers $t$. We will show in Sect. 3.7 that every matrix Lie group is an embedded submanifold of $\mathrm{GL}(n ; \mathbb{C})$. We will then show (Corollary 3.46) that $\mathfrak{g}$ is the tangent space to $G$ at the identity.

We will now establish various basic properties of the Lie algebra $\mathfrak{g}$ of a matrix Lie group $G$. In particular, we will see that there is a bracket operation on $\mathfrak{g}$ that makes $\mathfrak{g}$ into a Lie algebra in the sense of Definition 3.1.

Proposition 3.19. Let $G$ be a matrix Lie group, and $X$ an element of its Lie algebra. Then $e^{X}$ is an element of the identity component $G_{0}$ of $G$.

Proof. By definition of the Lie algebra, $e^{t X}$ lies in $G$ for all real $t$. However, as $t$ varies from 0 to $1, e^{t X}$ is a continuous path connecting the identity to $e^{X}$.

Theorem 3.20. Let $G$ be a matrix Lie group with Lie algebra $\mathfrak{g}$. If $X$ and $Y$ are elements of $\mathfrak{g}$, the following results hold.

\begin{enumerate}
  \item $A X A^{-1} \in \mathfrak{g}$ for all $A \in G$.
  \item $s X \in \mathfrak{g}$ for all real numbers $s$.
  \item $X+Y \in \mathfrak{g}$.
  \item $X Y-Y X \in \mathfrak{g}$.
\end{enumerate}

It follows from this result and Example 3.3 that the Lie algebra of a matrix Lie group is a real Lie algebra, with bracket given by $[X, Y]=X Y-Y X$. For $X$ and $Y$ in $\mathfrak{g}$, we refer to $[X, Y]=X Y-Y X \in \mathfrak{g}$ as the bracket or commutator of $X$ and $Y$.

Proof. For Point 1, we observe that, by Proposition 2.3,

$$
e^{t\left(A X A^{-1}\right)}=A e^{t X} A^{-1} \in G
$$

for all $t$, showing that $A X A^{-1}$ is in $\mathfrak{g}$. For Point 2 , we observe that $e^{t(s X)}=e^{(t s) X}$, which must be in $G$ for all $t \in \mathbb{R}$ if $X$ is in $\mathfrak{g}$. For Point 3 we use the Lie product formula, which says that

$$
e^{t(X+Y)}=\lim _{m \rightarrow \infty}\left(e^{t X / m} e^{t Y / m}\right)^{m}
$$

Thus, $\left(e^{t X / m} e^{t Y / m}\right)^{m}$ is in $G$ for all $m$. Since $G$ is closed, the limit (which is invertible) must be again in $G$. This shows that $X+Y$ is again in $\mathfrak{g}$.

Finally, for Point 4, we use the product rule (Exercise 3) and Proposition 2.4 to compute

$$
\begin{aligned}
\left.\frac{d}{d t}\left(e^{t X} Y e^{-t X}\right)\right|_{t=0} & =(X Y) e^{0}+\left(e^{0} Y\right)(-X) \\
& =X Y-Y X
\end{aligned}
$$

Now, by Point $1, e^{t X} Y e^{-t X}$ is in $\mathfrak{g}$ for all $t$. Furthermore, by Points 2 and $3, \mathfrak{g}$ is a real subspace of $M_{n}(\mathbb{C})$, from which it follows that $\mathfrak{g}$ is a (topologically) closed subset of $M_{n}(\mathbb{C})$. Thus,

$$
X Y-Y X=\lim _{h \rightarrow 0} \frac{e^{h X} Y e^{-h X}-Y}{h}
$$

belongs to $\mathfrak{g}$.\\
Note that even if the elements of $G$ have complex entries, the Lie algebra $\mathfrak{g}$ of $G$ is not necessarily a complex vector space, since Point 2 holds, in general, only for $s \in \mathbb{R}$. Nevertheless, it may happen in certain cases that $\mathfrak{g}$ is a complex vector space.

Definition 3.21. A matrix Lie group $G$ is said to be complex if its Lie algebra $\mathfrak{g}$ is a complex subspace of $M_{n}(\mathbb{C})$, that is, if $i X \in \mathfrak{g}$ for all $X \in \mathfrak{g}$.

Examples of complex groups are $\operatorname{GL}(n ; \mathbb{C}), \operatorname{SL}(n ; \mathbb{C}), \operatorname{SO}(n ; \mathbb{C})$, and $\operatorname{Sp}(n ; \mathbb{C})$, as the calculations in Sect. 3.4 will show.

Proposition 3.22. If $G$ is commutative then $\mathfrak{g}$ is commutative.\\
We will see in Sect. 3.7 that if $G$ is connected and $\mathfrak{g}$ is commutative, $G$ must be commutative.

Proof. For any two matrices $X, Y \in M_{n}(\mathbb{C})$, the commutator of $X$ and $Y$ may be computed as

$$
[X, Y]=\left.\frac{d}{d t}\left(\left.\frac{d}{d s} e^{t X} e^{s Y} e^{-t X}\right|_{s=0}\right)\right|_{t=0} .
$$

If $G$ is commutative and $X$ and $Y$ belong to $\mathfrak{g}$, then $e^{t X}$ commutes with $e^{s Y}$ and the expression in parentheses on the right hand side of (3.7) is independent of $t$, so that $[X, Y]=0$.

\subsection*{Examples}
Physicists are accustomed to using the map $t \mapsto e^{i t X}$ rather than $t \mapsto e^{t X}$. Thus, the physicists' expressions for the Lie algebras of matrix Lie groups will differ by a factor of $i$ from the expressions we now derive.

Proposition 3.23. The Lie algebra of $\mathrm{GL}(n ; \mathbb{C})$ is the space $M_{n}(\mathbb{C})$ of all $n \times n$ matrices with complex entries. Similarly, the Lie algebra of $\mathrm{GL}(n ; \mathbb{R})$ is equal to $M_{n}(\mathbb{R})$. The Lie algebra of $\operatorname{SL}(n ; \mathbb{C})$ consists of all $n \times n$ complex matrices with trace zero, and the Lie algebra of $\operatorname{SL}(n ; \mathbb{R})$ consists of all $n \times n$ real matrices with trace zero.

We denote the Lie algebras of these groups as $\operatorname{gl}(n ; \mathbb{C}), \operatorname{gl}(n ; \mathbb{R}), \operatorname{sl}(n ; \mathbb{C})$, and $\mathrm{sI}(n ; \mathbb{R})$, respectively.\\
Proof. If $X \in M_{n}(\mathbb{C})$, then $e^{t X}$ is invertible, so that $X$ belongs to the Lie algebra of $\operatorname{GL}(n ; \mathbb{C})$. If $X \in M_{n}(\mathbb{R})$, then $e^{t X}$ is invertible and real, so that $X$ is in the Lie algebra of $\mathrm{GL}(n ; \mathbb{R})$. Conversely, if $e^{t X}$ is real for all real $t$, then $X= d\left(e^{t X}\right) /\left.d t\right|_{t=0}$ must also real. If $X \in M_{n}(\mathbb{C})$ has trace zero, then by Theorem 2.12, $\operatorname{det}\left(e^{t X}\right)=1$, showing that $X$ is in the Lie algebra of $\operatorname{SL}(n ; \mathbb{C})$. Conversely, if $\operatorname{det}\left(e^{t X}\right)=e^{t \operatorname{trace}(X)}=1$ for all real $t$, then

$$
\operatorname{trace}(X)=\left.\frac{d}{d t} e^{t \operatorname{trace}(X)}\right|_{t=0}=0
$$

Finally, if $X$ is real and has trace zero, then $e^{t X}$ is real and has determinant 1 for all real $t$, showing that $X$ is in the Lie algebra of $\operatorname{SL}(n ; \mathbb{R})$. Conversely, if $e^{t X}$ is real and has determinant 1 for all real $t$, the preceding arguments show that $X$ must be real and have trace zero.

Proposition 3.24. The Lie algebra of $\mathbf{U}(n)$ consists of all complex matrices satisfying $X^{*}=-X$ and the Lie algebra of $\mathrm{SU}(n)$ consists of all complex matrices satisfying $X^{*}=-X$ and trace $(X)=0$. The Lie algebra of the orthogonal group $\mathrm{O}(n)$ consists of all real matrices $X$ satisfying $X^{t r}=-X$ and the Lie algebra of $\mathrm{SO}(n)$ is the same as that of $\mathrm{O}(n)$.

The Lie algebras of $\mathrm{U}(n)$ and $\mathrm{SU}(n)$ are denoted $\mathrm{u}(n)$ and $\mathrm{su}(n)$, respectively. The Lie algebra of $\mathrm{SO}(n)$ (which is the same as that of $\mathrm{O}(n)$ ) is denoted $\operatorname{so}(n)$.\\
Proof. A matrix $U$ is unitary if and only if $U^{*}=U^{-1}$. Thus, $e^{t X}$ is unitary if and only if

$$
\left(e^{t X}\right)^{*}=\left(e^{t X}\right)^{-1}=e^{-t X}
$$

By Point 2 of Proposition 2.3, $\left(e^{t X}\right)^{*}=e^{t X^{*}}$, and so (3.8) becomes

$$
e^{t X^{*}}=e^{-t X}
$$

The condition (3.9) holds for all real $t$ if and only if $X^{*}=-X$. Thus, the Lie algebra of $\mathrm{U}(n)$ consists precisely of matrices $X$ such that $X^{*}=-X$. As in the proof of Proposition 3.23, adding the "determinant 1" condition at the group level adds the "trace 0 " condition at the Lie algebra level.

An exactly similar argument over $\mathbb{R}$ shows that a real matrix $X$ belongs to the Lie algebra of $\mathrm{O}(n)$ if and only if $X^{t r}=-X$. Since any such matrix has trace $(X)=0$ (since the diagonal entries of $X$ are all zero), we see that every element of the Lie algebra of $\mathrm{O}(n)$ is also in the Lie algebra of $\mathrm{SO}(n)$.

Proposition 3.25. If $g$ is the matrix in Exercise 1 of Chapter 1, then the Lie algebra of $\mathrm{O}(n ; k)$ consists precisely of those real matrices $X$ such that

$$
g X^{t r} g=-X,
$$

and the Lie algebra of $\mathrm{SO}(n ; k)$ is the same as that of $\mathrm{O}(n ; k)$. If $\Omega$ is the matrix (1.8), then the Lie algebra of $\operatorname{Sp}(n ; \mathbb{R})$ consists precisely of those real matrices $X$ such that

$$
\Omega X^{t r} \Omega=X,
$$

and the Lie algebra of $\operatorname{Sp}(n ; \mathbb{C})$ consists precisely of those complex matrices $X$ satisfying the same condition. The Lie algebra of $\operatorname{Sp}(n)$ consists precisely of those complex matrices $X$ such that $\Omega X^{t r} \Omega=X$ and $X^{*}=-X$.

The verification of Proposition 3.25 is similar to our previous computations and is omitted. The Lie algebra of $\mathrm{SO}(n ; k)$ (which is the same as that of $\mathrm{O}(n ; k)$ ) is denoted $\operatorname{so}(n ; k)$, whereas the Lie algebras of the symplectic groups are denoted $\operatorname{sp}(n ; \mathbb{R}), \operatorname{sp}(n ; \mathbb{C})$, and $\operatorname{sp}(n)$.

Proposition 3.26. The Lie algebra of the Heisenberg group $H$ in Sect. 1.2.6 is the space of all matrices of the form

$$
X=\left(\begin{array}{lll}
0 & a & b \\
0 & 0 & c \\
0 & 0 & 0
\end{array}\right),
$$

with $a, b, c \in \mathbb{R}$.\\
Proof. If $X$ is strictly upper triangular, it is easy to verify that $X^{m}$ will be strictly upper triangular for all positive integers $m$. Thus, for $X$ as in (3.10), we will have $e^{t X}=I+B$ with $B$ strictly upper triangular, showing that $e^{t X} \in H$. Conversely, if $e^{t X}$ belongs to $H$ for all real $t$, then all of the entries of $e^{t X}$ on or below the diagonal are independent of $t$. Thus, $X=d\left(e^{t X}\right) /\left.d t\right|_{t=0}$ will be of the form in (3.10).

We leave it as an exercise to determine the Lie algebras of the Euclidean and Poincaré groups.

Example 3.27. The following elements form a basis for the Lie algebra su(2):

$$
E_{1}=\frac{1}{2}\left(\begin{array}{rr}
i & 0 \\
0 & -i
\end{array}\right) ; \quad E_{2}=\frac{1}{2}\left(\begin{array}{ll}
0 & i \\
i & 0
\end{array}\right) ; \quad E_{3}=\frac{1}{2}\left(\begin{array}{rr}
0 & -1 \\
1 & 0
\end{array}\right) .
$$

These elements satisfy the commutation relations $\left[E_{1}, E_{2}\right]=E_{3},\left[E_{2}, E_{3}\right]=E_{1}$, and $\left[E_{3}, E_{1}\right]=E_{2}$. The following elements form a basis for the Lie algebra so(3):

$$
F_{1}=\left(\begin{array}{rrr}
0 & 0 & 0 \\
0 & 0 & -1 \\
0 & 1 & 0
\end{array}\right), \quad F_{2}=\left(\begin{array}{rrr}
0 & 0 & 1 \\
0 & 0 & 0 \\
-1 & 0 & 0
\end{array}\right), \quad F_{3}=\left(\begin{array}{rrr}
0 & -1 & 0 \\
1 & 0 & 0 \\
0 & 0 & 0
\end{array}\right),
$$

These elements satisfy the commutation relations $\left[F_{1}, F_{2}\right]=F_{3},\left[F_{2}, F_{3}\right]=F_{1}$, and $\left[F_{3}, F_{1}\right]=F_{2}$.

Note that the listed relations completely determine all commutation relations among, say, $E_{1}, E_{2}$, and $E_{3}$, since by the skew symmetry of the bracket, we must have $\left[E_{1}, E_{1}\right]=0,\left[E_{2}, E_{1}\right]=-E_{3}$, and so on. Since $E_{1}, E_{2}$, and $E_{3}$ satisfy the same commutation relations as $F_{1}, F_{2}$, and $F_{3}$, the two Lie algebras are isomorphic.

Proof. Direct calculation from Proposition 3.24.

\subsection*{Lie Group and Lie Algebra Homomorphisms}
The following theorem tells us that a Lie group homomorphism between two Lie groups gives rise in a natural way to a map between the corresponding Lie algebras. It will follow (Exercise 8) that isomorphic Lie groups have isomorphic Lie algebras.

Theorem 3.28. Let $G$ and $H$ be matrix Lie groups, with Lie algebras $\mathfrak{g}$ and $\mathfrak{h}$, respectively. Suppose that $\Phi: G \rightarrow H$ is a Lie group homomorphism. Then there exists a unique real-linear map $\phi: \mathfrak{g} \rightarrow \mathfrak{h}$ such that

$$
\Phi\left(e^{X}\right)=e^{\phi(X)}
$$

for all $X \in \mathfrak{g}$. The map $\phi$ has following additional properties:

\begin{enumerate}
  \item $\phi\left(A X A^{-1}\right)=\Phi(A) \phi(X) \Phi(A)^{-1}$, for all $X \in \mathfrak{g}, A \in G$.
  \item $\phi([X, Y])=[\phi(X), \phi(Y)]$, for all $X, Y \in \mathfrak{g}$.
  \item $\phi(X)=\left.\frac{d}{d t} \Phi\left(e^{t X}\right)\right|_{t=0}$, for all $X \in \mathfrak{g}$.
\end{enumerate}

In practice, given a Lie group homomorphism $\Phi$, the way one goes about computing $\phi$ is by using Property 3. In the language of manifolds, Property 3 says that $\phi$ is the derivative (or differential) of $\Phi$ at the identity. By Point $2, \phi: \mathfrak{g} \rightarrow \mathfrak{h}$ is a Lie algebra homomorphism. Thus, every Lie group homomorphism gives rise to a Lie algebra homomorphism. In Chapter 5, we will investigate the reverse question: If $\phi$ is a homomorphism between the Lie algebras of two Lie groups, is there an associated Lie group homomorphism $\Phi$ ?

Proof. The proof is similar to the proof of Theorem 3.20. Since $\Phi$ is a continuous group homomorphism, $\Phi\left(e^{t X}\right)$ will be a one-parameter subgroup of $H$, for each $X \in \mathfrak{g}$. Thus, by Theorem 2.14, there is a unique matrix $Z$ such that

$$
\Phi\left(e^{t X}\right)=e^{t Z}
$$

for all $t \in \mathbb{R}$. We define $\phi(X)=Z$ and check that $\phi$ has the required properties. First, by putting $t=1$ in (3.12), we see that $\Phi\left(e^{X}\right)=e^{\phi(X)}$ for all $X \in \mathfrak{g}$. Next, if $\Phi\left(e^{t X}\right)=e^{t Z}$ for all $t$, then $\Phi\left(e^{t s X}\right)=e^{t s Z}$, showing that $\phi(s X)=s \phi(X)$. Using the Lie product formula and the continuity of $\Phi$, we then compute that

$$
\begin{aligned}
e^{t \phi(X+Y)} & =\Phi\left(\lim _{m \rightarrow \infty}\left(e^{t X / m} e^{t Y / m}\right)^{m}\right) \\
& =\lim _{m \rightarrow \infty}\left(\Phi\left(e^{t X / m}\right) \Phi\left(e^{t Y / m}\right)\right)^{m}
\end{aligned}
$$

Thus,

$$
e^{t \phi(X+Y)}=\lim _{m \rightarrow \infty}\left(e^{t \phi(X) / m} e^{t \phi(Y) / m}\right)^{m}=e^{t(\phi(X)+\phi(Y))}
$$

Differentiating this result at $t=0$ shows that $\phi(X+Y)=\phi(X)+\phi(Y)$.\\
We have thus obtained a real-linear map $\phi$ satisfying (3.11). If there were another real-linear map $\phi^{\prime}$ with this property, we would have

$$
e^{t \phi(X)}=e^{t \phi^{\prime}(X)}=\Phi\left(e^{t X}\right)
$$

for all $t \in \mathbb{R}$. Differentiating this result at $t=0$ shows that $\phi(X)=\phi^{\prime}(X)$.\\
We now verify the remaining claimed properties of $\phi$. For any $A \in G$, we have

$$
e^{t \phi\left(A X A^{-1}\right)}=e^{\phi\left(t A X A^{-1}\right)}=\Phi\left(e^{t A X A^{-1}}\right) .
$$

Thus,

$$
\begin{aligned}
e^{t \phi\left(A X A^{-1}\right)} & =\Phi(A) \Phi\left(e^{t X}\right) \Phi(A)^{-1} \\
& =\Phi(A) e^{t \phi(X)} \Phi(A)^{-1} .
\end{aligned}
$$

Differentiating this identity at $t=0$ gives Point 1 .\\
Meanwhile, for any $X$ and $Y$ in $\mathfrak{g}$, we have, as in the proof of Theorem 3.20,

$$
\begin{aligned}
\phi([X, Y]) & =\phi\left(\left.\frac{d}{d t} e^{t X} Y e^{-t X}\right|_{t=0}\right) \\
& =\left.\frac{d}{d t} \phi\left(e^{t X} Y e^{-t X}\right)\right|_{t=0},
\end{aligned}
$$

where we have used the fact that a derivative commutes with a linear transformation. Thus,

$$
\begin{aligned}
\phi([X, Y]) & =\left.\frac{d}{d t} \Phi\left(e^{t X}\right) \phi(Y) \Phi\left(e^{-t X}\right)\right|_{t=0} \\
& =\left.\frac{d}{d t} e^{t \phi(X)} \phi(Y) e^{-t \phi(X)}\right|_{t=0} \\
& =[\phi(X), \phi(Y)]
\end{aligned}
$$

establishing Point 2. Finally, since $\Phi\left(e^{t X}\right)=e^{\phi(t X)}=e^{t \phi(X)}$, we can compute $\phi(X)$ as in Point 3.

Example 3.29. Let $\Phi: \mathrm{SU}(2) \rightarrow \mathrm{SO}(3)$ be the homomorphism in Proposition 1.19. Then the associated Lie algebra homomorphism $\phi: \operatorname{su}(2) \rightarrow \operatorname{so}(3)$ satisfies

$$
\phi\left(E_{j}\right)=F_{j}, \quad j=1,2,3,
$$

where $\left\{E_{1}, E_{2}, E_{3}\right\}$ and $\left\{F_{1}, F_{2}, F_{3}\right\}$ are the bases for $\mathrm{su}(2)$ and so(3), respectively, given in Example 3.27.

Since $\phi$ maps a basis for su(2) to a basis for so(3), we see that $\phi$ is a Lie algebra isomorphism, even though $\Phi$ is not a Lie group isomorphism (since $\operatorname{ker}(\Phi)=\{I,-I\})$.

Proof. If $X$ is in su(2) and $Y$ is in the space $V$ in (1.14), then

$$
\left.\frac{d}{d t} \Phi\left(e^{t X}\right) Y\right|_{t=0}=\left.\frac{d}{d t} e^{t X} Y e^{-t X}\right|_{t=0}=[X, Y]
$$

Thus, $\phi(X)$ is the linear map of $V \cong \mathbb{R}^{3}$ to itself given by $Y \mapsto[X, Y]$. If, say, $X=E_{1}$, then direct computation shows that

$$
\left[E_{1},\left(\begin{array}{cc}
x_{1} & x_{2}+i x_{3} \\
x_{2}-i x_{3} & -x_{1}
\end{array}\right)\right]=\left(\begin{array}{cc}
x_{1}^{\prime} & x_{2}^{\prime}+i x_{3}^{\prime} \\
x_{2}^{\prime}-i x_{3}^{\prime} & -x_{1}^{\prime}
\end{array}\right),
$$

where $\left(x_{1}^{\prime}, x_{2}^{\prime}, x_{3}^{\prime}\right)=\left(0,-x_{3}, x_{2}\right)$. Since

$$
\left(\begin{array}{c}
0 \\
-x_{3} \\
x_{2}
\end{array}\right)=\left(\begin{array}{rrr}
0 & 0 & 0 \\
0 & 0 & -1 \\
0 & 1 & 0
\end{array}\right)\left(\begin{array}{l}
x_{1} \\
x_{2} \\
x_{3}
\end{array}\right),
$$

we conclude that $\phi\left(E_{1}\right)$ is the $3 \times 3$ matrix appearing on the right-hand side of (3.13), which is precisely $F_{1}$. The computation of $\phi\left(E_{2}\right)$ and $\phi\left(E_{3}\right)$ is similar and is left to the reader.

Proposition 3.30. Suppose that $G, H$, and $K$ are matrix Lie groups and $\Phi: H \rightarrow K$ and $\Psi: G \rightarrow H$ are Lie group homomorphisms. Let $\Lambda: G \rightarrow K$ be the composition of $\Phi$ and $\Psi$ and let $\phi, \psi$, and $\lambda$ be the Lie algebra maps associated to $\Phi, \Psi$, and $\Lambda$, respectively. Then we have

$$
\lambda=\phi \circ \psi .
$$

Proof. For any $X \in \mathfrak{g}$,

$$
\Lambda\left(e^{t X}\right)=\Phi\left(\Psi\left(e^{t X}\right)\right)=\Phi\left(e^{t \psi(X)}\right)=e^{t \phi(\psi(X))}
$$

Thus, $\lambda(X)=\phi(\psi(X))$.\\
Proposition 3.31. If $\Phi: G \rightarrow H$ is a Lie group homomorphism and $\phi: \mathfrak{g} \rightarrow \mathfrak{h}$ is the associated Lie algebra homomorphism, then the kernel of $\Phi$ is a closed, normal subgroup of $G$ and the Lie algebra of the kernel is given by

$$
\operatorname{Lie}(\operatorname{ker}(\Phi))=\operatorname{ker}(\phi)
$$

Proof. The usual algebraic argument shows that $\operatorname{ker}(\Phi)$ is normal subgroup of $G$. Since, also, $\Phi$ is continuous, $\operatorname{ker}(\Phi)$ is closed. If $X \in \operatorname{ker}(\phi)$, then

$$
\Phi\left(e^{t X}\right)=e^{t \phi(X)}=I
$$

for all $t \in \mathbb{R}$, showing that $X$ is in the Lie algebra of $\operatorname{ker}(\Phi)$. In the other direction, if $e^{t X}$ lies in $\operatorname{ker}(\Phi)$ for all $t \in \mathbb{R}$, then

$$
e^{t \phi(X)}=\Phi\left(e^{t X}\right)=I
$$

for all $t$. Differentiating this relation with respect to $t$ at $t=0$ gives that $\phi(X)=0$, showing that $X \in \operatorname{ker}(\phi)$.

Definition 3.32 (The Adjoint Map). Let $G$ be a matrix Lie group, with Lie algebra $\mathfrak{g}$. Then for each $A \in G$, define a linear map $\operatorname{Ad}_{A}: \mathfrak{g} \rightarrow \mathfrak{g}$ by the formula

$$
\operatorname{Ad}_{A}(X)=A X A^{-1}
$$

Proposition 3.33. Let $G$ be a matrix Lie group, with Lie algebra $\mathfrak{g}$. Let $\mathrm{GL}(\mathfrak{g})$ denote the group of all invertible linear transformations of $\mathfrak{g}$. Then the map $A \rightarrow \mathrm{Ad}_{A}$ is a homomorphism of $G$ into $\mathrm{GL}(\mathfrak{g})$. Furthermore, for each $A \in G, \operatorname{Ad}_{A}$ satisfies $\operatorname{Ad}_{A}([X, Y])=\left[\operatorname{Ad}_{A}(X), \operatorname{Ad}_{A}(Y)\right]$ for all $X, Y \in \mathfrak{g}$.

Proof. Easy. Note that Point 1 of Theorem 3.20 guarantees that $\operatorname{Ad}_{A}(X)$ is actually in $\mathfrak{g}$ for all $X \in \mathfrak{g}$.

Since $\mathfrak{g}$ is a real vector space with some dimension $k, \mathrm{GL}(\mathfrak{g})$ is essentially the same as $\mathrm{GL}(k ; \mathbb{R})$. Thus, we will regard $\mathrm{GL}(\mathfrak{g})$ as a matrix Lie group. It is easy to\\
show that $\mathrm{Ad}: G \rightarrow \mathrm{GL}(\mathfrak{g})$ is continuous, and so is a Lie group homomorphism. By Theorem 3.28, there is an associated real linear map $X \rightarrow \mathrm{ad}_{X}$ from the Lie algebra of $G$ to the Lie algebra of $\mathrm{GL}(\mathfrak{g})$ (i.e., from $\mathfrak{g}$ to $\mathrm{gl}(\mathfrak{g})$ ), with the property that

$$
e^{\mathrm{ad}_{X}}=\operatorname{Ad}_{e^{X}} .
$$

Here, $\operatorname{gl}(\mathfrak{g})$ is the Lie algebra of $\operatorname{GL}(\mathfrak{g})$, namely the space of all linear maps of $\mathfrak{g}$ to itself.

Proposition 3.34. Let $G$ be a matrix Lie group, let $\mathfrak{g}$ be its Lie algebra, and let $\mathrm{Ad}: G \rightarrow \mathrm{GL}(\mathfrak{g})$ be as in Proposition 3.33. Let $\mathrm{ad}: \mathfrak{g} \rightarrow \mathrm{gl}(\mathfrak{g})$ be the associated Lie algebra map. Then for all $X, Y \in \mathfrak{g}$

$$
\operatorname{ad}_{X}(Y)=[X, Y] .
$$

The proposition shows that our usage of the notation $\mathrm{ad}_{X}$ in this section is consistent with that in Definition 3.7.

Proof. By Point 3 of Theorem 3.28, ad can be computed as follows:

$$
\operatorname{ad}_{X}=\left.\frac{d}{d t} \operatorname{Ad}_{e^{t X}}\right|_{t=0} .
$$

Thus,

$$
\operatorname{ad}_{X}(Y)=\left.\frac{d}{d t} e^{t X} Y e^{-t X}\right|_{t=0}=[X, Y],
$$

as claimed.\\
We have proved, as a consequence of Theorem 3.28 and Proposition 3.34, the following result, which we will make use of later.

Proposition 3.35. For any $X$ in $M_{n}(\mathbb{C})$, let $\operatorname{ad}_{X}: M_{n}(\mathbb{C}) \rightarrow M_{n}(\mathbb{C})$ be given by $\operatorname{ad}_{X} Y=[X, Y]$. Then for any $Y$ in $M_{n}(\mathbb{C})$, we have

$$
e^{X} Y e^{-X}=\operatorname{Ad}_{e^{X}}(Y)=e^{\operatorname{ad}_{X}}(Y),
$$

where

$$
e^{\operatorname{ad}_{X}}(Y)=Y+[X, Y]+\frac{1}{2}[X,[X, Y]]+\cdots .
$$

This result can also be proved by direct calculation-see Exercise 14.

\subsection*{The Complexification of a Real Lie Algebra}
In studying the representations of a matrix Lie group $G$ (as we will do in later chapters), it is often useful to pass to the Lie algebra $\mathfrak{g}$ of $G$, which is, in general, only a real Lie algebra. It is then often useful to pass to an associated complex Lie algebra, called the complexification of $\mathfrak{g}$.

Definition 3.36. If $V$ is a finite-dimensional real vector space, then the complexification of $V$, denoted $V_{\mathbb{C}}$, is the space of formal linear combinations

$$
v_{1}+i v_{2}
$$

with $v_{1}, v_{2} \in V$. This becomes a real vector space in the obvious way and becomes a complex vector space if we define

$$
i\left(v_{1}+i v_{2}\right)=-v_{2}+i v_{1} .
$$

We could more pedantically define $V_{\mathbb{C}}$ to be the space of ordered pairs ( $v_{1}, v_{2}$ ) with $v_{1}, v_{2} \in V$, but this is notationally cumbersome. It is straightforward to verify that the above definition really makes $V_{\mathbb{C}}$ into a complex vector space. We will regard $V$ as a real subspace of $V_{\mathbb{C}}$ in the obvious way.

Proposition 3.37. Let $\mathfrak{g}$ be a finite-dimensional real Lie algebra and $\mathfrak{g}_{\mathbb{C}}$ its complexification. Then the bracket operation on $\mathfrak{g}$ has a unique extension to $\mathfrak{g}_{\mathbb{C}}$ that makes $\mathfrak{g}_{\mathbb{C}}$ into a complex Lie algebra. The complex Lie algebra $\mathfrak{g}_{\mathbb{C}}$ is called the complexification of the real Lie algebra $\mathfrak{g}$.

Proof. The uniqueness of the extension is obvious, since if the bracket operation on $\mathfrak{g}_{\mathbb{C}}$ is to be bilinear, then it must be given by

$$
\left[X_{1}+i X_{2}, Y_{1}+i Y_{2}\right]=\left(\left[X_{1}, Y_{1}\right]-\left[X_{2}, Y_{2}\right]\right)+i\left(\left[X_{1}, Y_{2}\right]+\left[X_{2}, Y_{1}\right]\right) .
$$

To show existence, we must now check that (3.15) is really bilinear and skew symmetric and that it satisfies the Jacobi identity. It is clear that (3.15) is real bilinear, and skew-symmetric. The skew symmetry means that if (3.15) is complex linear in the first factor, it is also complex linear in the second factor. Thus, we need only show that

$$
\left[i\left(X_{1}+i X_{2}\right), Y_{1}+i Y_{2}\right]=i\left[X_{1}+i X_{2}, Y_{1}+i Y_{2}\right] .
$$

The left-hand side of (3.16) is

$$
\left[-X_{2}+i X_{1}, Y_{1}+i Y_{2}\right]=\left(-\left[X_{2}, Y_{1}\right]-\left[X_{1}, Y_{2}\right]\right)+i\left(\left[X_{1}, Y_{1}\right]-\left[X_{2}, Y_{2}\right]\right),
$$

whereas the right-hand side of (3.16) is

$$
\begin{aligned}
i\{ & \left.\left(\left[X_{1}, Y_{1}\right]-\left[X_{2}, Y_{2}\right]\right)+i\left(\left[X_{2}, Y_{1}\right]+\left[X_{1}, Y_{2}\right]\right)\right\} \\
& =\left(-\left[X_{2}, Y_{1}\right]-\left[X_{1}, Y_{2}\right]\right)+i\left(\left[X_{1}, Y_{1}\right]-\left[X_{2}, Y_{2}\right]\right)
\end{aligned}
$$

and, indeed, these expressions are equal.\\
It remains to check the Jacobi identity. Of course, the Jacobi identity holds if $X, Y$, and $Z$ are in $\mathfrak{g}$. Furthermore, for all $X, Y, Z \in \mathfrak{g}_{\mathbb{C}}$, the expression

$$
[X,[Y, Z]]+[Y,[Z, X]]+[Z,[X, Y]]
$$

is complex-linear in $X$ with $Y$ and $Z$ fixed. Thus, the Jacobi identity continues to hold if $X$ is in $\mathfrak{g}_{\mathbb{C}}$ and $Y$ and $Z$ are in $\mathfrak{g}$. The same argument then shows that the Jacobi identity holds when $X$ and $Y$ in $\mathfrak{g}_{\mathbb{C}}$ and $Z$ is in $\mathfrak{g}$. Applying this argument one more time establishes the Jacobi identity for $\mathfrak{g}_{\mathbb{C}}$ in general.

Proposition 3.38. Suppose that $\mathfrak{g} \subset M_{n}(\mathbb{C})$ is a real Lie algebra and that for all nonzero $X$ in $\mathfrak{g}$, the element iX is not in $\mathfrak{g}$. Then the "abstract" complexification $\mathfrak{g}_{\mathbb{C}}$ of $\mathfrak{g}$ in Definition 3.36 is isomorphic to the set of matrices in $M_{n}(\mathbb{C})$ that can be expressed in the form $X+i Y$ with $X$ and $Y$ in $\mathfrak{g}$.

Proof. Consider the map from $\mathfrak{g}_{\mathbb{C}}$ into $M_{n}(\mathbb{C})$ sending the formal linear combination $X+i Y$ to the linear combination $X+i Y$ of matrices. This map is easily seen to be a complex Lie algebra homomorphism. If $\mathfrak{g}$ satisfies the assumption in the statement of the proposition, this map is also injective and thus an isomorphism of $\mathfrak{g}_{\mathbb{C}}$ with $\mathfrak{g}+i \mathfrak{g} \subset M_{n}(\mathbb{C})$.

Using the proposition, we easily obtain the following list of isomorphisms:

$$
\begin{aligned}
\mathrm{gl}(n ; \mathbb{R})_{\mathbb{C}} & \cong \mathrm{gl}(n ; \mathbb{C}), \\
\mathrm{u}(n)_{\mathbb{C}} & \cong \mathrm{gl}(n ; \mathbb{C}), \\
\mathrm{su}(n)_{\mathbb{C}} & \cong \mathrm{sl}(n ; \mathbb{C}), \\
\mathrm{sl}(n ; \mathbb{R})_{\mathbb{C}} & \cong \mathrm{sl}(n ; \mathbb{C}), \\
\mathrm{so}(n)_{\mathbb{C}} & \cong \mathrm{so}(n ; \mathbb{C}), \\
\mathrm{sp}(n ; \mathbb{R})_{\mathbb{C}} & \cong \mathrm{sp}(n ; \mathbb{C}), \\
\mathrm{sp}(n)_{\mathbb{C}} & \cong \mathrm{sp}(n ; \mathbb{C}) .
\end{aligned}
$$

Let us verify just one example, that of $\mathrm{u}(n)$. If $X^{*}=-X$, then $(i X)^{*}=i X$. Thus, $X$ and $X^{*}$ cannot both be in $\mathrm{u}(n)$ unless $X$ is zero. Furthermore, every $X$ in $M_{n}(\mathbb{C})$ can be expressed as $X=X_{1}+i X_{2}$, where $X_{1}=\left(X-X^{*}\right) / 2$ and $X_{2}=\left(X+X^{*}\right) /(2 i)$ are both in $\mathrm{u}(n)$. This shows that $\mathrm{u}(n)_{\mathbb{C}} \cong \mathrm{gl}(n ; \mathbb{C})$.

Although both $\mathrm{su}(2)_{\mathbb{C}}$ and $\mathrm{sI}(2 ; \mathbb{R})_{\mathbb{C}}$ are isomorphic to $\mathrm{sI}(2 ; \mathbb{C})$, the Lie algebra su(2) is not isomorphic to $\mathrm{SI}(2 ; \mathbb{R})$. See Exercise 11.

Proposition 3.39. Let $\mathfrak{g}$ be a real Lie algebra, $\mathfrak{g}_{\mathbb{C}}$ its complexification, and $\mathfrak{h}$ an arbitrary complex Lie algebra. Then every real Lie algebra homomorphism of $\mathfrak{g}$ into $\mathfrak{h}$ extends uniquely to a complex Lie algebra homomorphism of $\mathfrak{g}_{\mathbb{C}}$ into $\mathfrak{h}$.

This result is the universal property of the complexification of a real Lie algebra.

Proof. The unique extension is given by $\pi(X+i Y)=\pi(X)+i \pi(Y)$ for all $X, Y \in \mathfrak{g}$. It is easy to check that this map is, indeed, a homomorphism of complex Lie algebras.

\subsection*{The Exponential Map}
Definition 3.40. If $G$ is a matrix Lie group with Lie algebra $\mathfrak{g}$, then the exponential map for $G$ is the map

$$
\exp : \mathfrak{g} \rightarrow G .
$$

That is to say, the exponential map for $G$ is the matrix exponential restricted to the Lie algebra $\mathfrak{g}$ of $G$. We have shown (Theorem 2.10) that every matrix in $\mathrm{GL}(n ; \mathbb{C})$ is the exponential of some $n \times n$ matrix. Nevertheless, if $G \subset \operatorname{GL}(n ; \mathbb{C})$ is a closed subgroup, there may exist $A$ in $G$ such that there is no $X$ in the Lie algebra $\mathfrak{g}$ of $G$ with $\exp X=A$.

Example 3.41. There does not exist a matrix $X \in \mathrm{sI}(2 ; \mathbb{C})$ with

$$
e^{X}=\left(\begin{array}{rr}
-1 & 1 \\
0 & -1
\end{array}\right),
$$

even though the matrix on the right-hand side of (3.18) is in $\operatorname{SL}(2 ; \mathbb{C})$.\\
Proof. If $X \in \operatorname{sl}(2 ; \mathbb{C})$ has distinct eigenvalues, then $X$ is diagonalizable and $e^{X}$ will also be diagonalizable, unlike the matrix on the right-hand side of (3.18). If $X \in \operatorname{sl}(2 ; \mathbb{C})$ has a repeated eigenvalue, this eigenvalue must be 0 or the trace of $X$ would not be zero. Thus, there is a nonzero vector $v$ with $X v=0$, from which it follows that $e^{X} v=e^{0} v=v$. We conclude that $e^{X}$ has 1 as an eigenvalue, unlike the matrix on the right-hand side of (3.18).

We see, then, that the exponential map for a matrix Lie group $G$ does not necessarily map $\mathfrak{g}$ onto $G$. Furthermore, the exponential map may not be one-toone on $\mathfrak{g}$, as may be seen, for example, from the case $\mathfrak{g}=\operatorname{su}(2)$. Nevertheless, it provides a crucial mechanism for passing information between the group and the Lie algebra. Indeed, we will see (Corollary 3.44) that the exponential map is locally one-to-one and onto, a result that will be essential later.

Theorem 3.42. For $0<\varepsilon<\log 2$, let $U_{\varepsilon}=\left\{X \in M_{n}(\mathbb{C}) \mid\|X\|<\varepsilon\right\}$ and let $V_{\varepsilon}=\exp \left(U_{\varepsilon}\right)$. Suppose $G \subset \mathrm{GL}(n ; \mathbb{C})$ is a matrix Lie group with Lie algebra $\mathfrak{g}$. Then there exists $\varepsilon \in(0, \log 2)$ such that for all $A \in V_{\varepsilon}, A$ is in $G$ if and only if $\log A$ is in $\mathfrak{g}$.

The condition $\varepsilon<\log 2$ guarantees (Theorem 2.8) that for all $X \in U_{\varepsilon}, \log \left(e^{X}\right)$ is defined and equal to $X$. Note that if $X=\log A$ is in $\mathfrak{g}$, then $A=e^{X}$ is in $G$. Thus, the content of the theorem is that for some $\varepsilon$, having $A$ in $V_{\varepsilon} \cap G$ implies that $\log A$ must be in $\mathfrak{g}$. See Figure 3.1.

We begin with a lemma.\\
Lemma 3.43. Suppose $B_{m}$ are elements of $G$ and that $B_{m} \rightarrow I$. Let $Y_{m}=\log B_{m}$, which is defined for all sufficiently large $m$. Suppose that $Y_{m}$ is nonzero for all $m$ and that $Y_{m} /\left\|Y_{m}\right\| \rightarrow Y \in M_{n}(\mathbb{C})$. Then $Y$ is in $\mathfrak{g}$.

Proof. For any $t \in \mathbb{R}$, we have $\left(t /\left\|Y_{m}\right\|\right) Y_{m} \rightarrow t Y$. Note that since $B_{m} \rightarrow I$, we have $\left\|Y_{m}\right\| \rightarrow 0$. Thus, we can find integers $k_{m}$ such that $k_{m}\left\|Y_{m}\right\| \rightarrow t$. We have, then,

$$
e^{k_{m} Y_{m}}=\exp \left[\left(k_{m}\left\|Y_{m}\right\|\right) \frac{Y_{m}}{\left\|Y_{m}\right\|}\right] \rightarrow e^{t Y} .
$$

\includegraphics[scale=0.2, center]{2025_08_28_fc86fd4ade69f41b5842g-081}

Fig. 3.1 If $A \in V_{\varepsilon}$ belongs to $G$, then $\log A$ belongs to $\mathfrak{g}$


\includegraphics[scale=0.2, center]{2025_08_28_fc86fd4ade69f41b5842g-082}

Fig. 3.2 The points $k_{m} Y_{m}$ are converging to $t Y$


However,

$$
e^{k_{m} Y_{m}}=\left(e^{Y_{m}}\right)^{k_{m}}=\left(B_{m}\right)^{k_{m}} \in G
$$

and $G$ is closed, and we conclude that $e^{t Y} \in G$. This shows that $Y \in \mathfrak{g}$. (See Figure 3.2.)

Proof of Theorem 3.42. Let us think of $M_{n}(\mathbb{C})$ as $\mathbb{C}^{n^{2}} \cong \mathbb{R}^{2 n^{2}}$ and let $D$ denote the orthogonal complement of $\mathfrak{g}$ with respect to the usual inner product on $\mathbb{R}^{2 n^{2}}$. Consider the map $\Phi: M_{n}(\mathbb{C}) \rightarrow M_{n}(\mathbb{C})$ given by

$$
\Phi(Z)=e^{X} e^{Y},
$$

where $Z=X+Y$ with $X \in \mathfrak{g}$ and $Y \in D$. Since (Proposition 2.16) the exponential is continuously differentiable, the map $\Phi$ is also continuously differentiable, and we may compute that

$$
\begin{aligned}
& \left.\frac{d}{d t} \Phi(t X, 0)\right|_{t=0}=X, \\
& \left.\frac{d}{d t} \Phi(0, t Y)\right|_{t=0}=Y .
\end{aligned}
$$

This calculation shows that the derivative of $\Phi$ at the point $0 \in \mathbb{R}^{2 n^{2}}$ is the identity. (Recall that the derivative at a point of a function from $\mathbb{R}^{2 n^{2}}$ to itself is a linear map of $\mathbb{R}^{2 n^{2}}$ to itself.) Since the derivative of $\Phi$ at the origin is invertible, the inverse function theorem says that $\Phi$ has a continuous local inverse, defined in a neighborhood of $I$.

We need to prove that for some $\varepsilon$, if $A \in V_{\varepsilon} \cap G$, then $\log A \in \mathfrak{g}$. If this were not the case, we could find a sequence $A_{m}$ in $G$ such that $A_{m} \rightarrow I$ as $m \rightarrow \infty$ and such that for all $m, \log A_{m} \notin \mathfrak{g}$. Using the local inverse of the map $\Phi$, we can write $A_{m}$ (for all sufficiently large $m$ ) as

$$
A_{m}=e^{X_{m}} e^{Y_{m}}, \quad X_{m} \in \mathfrak{g}, Y_{m} \in D
$$

with $X_{m}$ and $Y_{m}$ tending to zero as $m$ tends to infinity. We must have $Y_{m} \neq 0$, since otherwise we would have $\log A_{m}=X_{m} \in \mathfrak{g}$. Since $e^{X_{m}}$ and $A_{m}$ are in $G$, we see that

$$
B_{m}:=e^{-X_{m}} A_{m}=e^{Y_{m}}
$$

is in $G$.\\
Since the unit sphere in $D$ is compact, we can choose a subsequence of the $Y_{m}$ 's (still called $Y_{m}$ ) so that $Y_{m} /\left\|Y_{m}\right\|$ converges to some $Y \in D$, with $\|Y\|=1$. Then, by the lemma, $Y \in \mathfrak{g}$. This is a contradiction, because $D$ is the orthogonal complement of $\mathfrak{g}$. Thus, there must be some $\varepsilon$ such that $\log A \in \mathfrak{g}$ for all $A$ in $V_{\varepsilon} \cap G$.

\subsection*{Consequences of Theorem 3.42}
In this section, we derive several consequences of the main result of the last section, Theorem 3.42.

Corollary 3.44. If $G$ is a matrix Lie group with Lie algebra $\mathfrak{g}$, there exists a neighborhood $U$ of 0 in $\mathfrak{g}$ and a neighborhood $V$ of $I$ in $G$ such that the exponential map takes $U$ homeomorphically onto $V$.

Proof. Let $\varepsilon$ be such that Theorem 3.42 holds and set $U=U_{\varepsilon} \cap \mathfrak{g}$ and $V= V_{\varepsilon} \cap G$. The theorem implies that $\exp$ takes $U$ onto $V$. Furthermore, $\exp$ is a homeomorphism of $U$ onto $V$, since there is a continuous inverse map, namely, the restriction of the matrix logarithm to $V$.

Corollary 3.45. Let $G$ be a matrix Lie group with Lie algebra $\mathfrak{g}$ and let $k$ be the dimension of $\mathfrak{g}$ as a real vector space. Then $G$ is a smooth embedded submanifold of $M_{n}(\mathbb{C})$ of dimension $k$ and hence a Lie group.

It follows from the corollary that $G$ is locally path connected: every point in $G$ has a neighborhood $U$ that is homeomorphic to a ball in $\mathbb{R}^{k}$ and hence path connected. It then follows that $G$ is connected (in the usual topological sense) if and only if it is path connected. (See, for example, Proposition 3.4.25 of [Run].)

Proof. Let $\varepsilon \in(0, \log 2)$ be such that Theorem 3.42 holds. Then for any $A_{0} \in G$, consider the neighborhood $A_{0} V_{\varepsilon}$ of $A_{0}$ in $M_{n}(\mathbb{C})$. Note that $A \in A_{0} V_{\varepsilon}$ if and only if $A_{0}^{-1} A \in V_{\varepsilon}$. Define a local coordinate system on $A_{0} V_{\varepsilon}$ by writing each $A \in A_{0} V_{\varepsilon}$ as $A=A_{0} e^{X}$, for $X \in U_{\varepsilon} \subset M_{n}(\mathbb{C})$. It follows from Theorem 3.42 that (for $\left.A \in A_{0} V_{\varepsilon}\right) A \in G$ if and only if $X \in \mathfrak{g}$. Thus, in this local coordinate system defined near $A_{0}$, the group $G$ looks like the subspace $\mathfrak{g}$ of $M_{n}(\mathbb{C})$. Since we can find such local coordinates near any point $A_{0}$ in $G$, we conclude that $G$ is an embedded submanifold of $M_{n}(\mathbb{C})$.

Now, the operation of matrix multiplication is clearly smooth. Furthermore, by the formula for the inverse of a matrix in terms of cofactors, the map $A \mapsto A^{-1}$ is also smooth on $\operatorname{GL}(n ; \mathbb{C})$. The restrictions of these maps to $G$ are then also smooth, showing that $G$ is a Lie group.

Corollary 3.46. Suppose $G \subset \mathrm{GL}(n ; \mathbb{C})$ is a matrix Lie group with Lie algebra $\mathfrak{g}$. Then a matrix $X$ is in $\mathfrak{g}$ if and only if there exists a smooth curve $\gamma$ in $M_{n}(\mathbb{C})$ with $\gamma(t) \in G$ for all $t$ and such that $\gamma(0)=I$ and $d \gamma /\left.d t\right|_{t=0}=X$. Thus, $\mathfrak{g}$ is the tangent space at the identity to $G$.

This result is illustrated in Figure 3.1.\\
Proof. If $X$ is in $\mathfrak{g}$, then we may take $\gamma(t)=e^{t X}$ and then $\gamma(0)=I$ and $d \gamma /\left.d t\right|_{t=0}=X$. In the other direction, suppose that $\gamma(t)$ is a smooth curve in $G$ with $\gamma(0)=I$. For all sufficiently small $t$, we can write $\gamma(t)=e^{\delta(t)}$, where $\delta$ is a smooth curve in $\mathfrak{g}$. Now, the derivative of $\delta(t)$ at $t=0$ is the same as the derivative of $t \mapsto t \delta^{\prime}(0)$ at $t=0$. Thus, by the chain rule, we have

$$
\gamma^{\prime}(0)=\left.\frac{d}{d t} e^{\delta(t)}\right|_{t=0}=\left.\frac{d}{d t} e^{t \delta^{\prime}(0)}\right|_{t=0}=\delta^{\prime}(0) .
$$

Since $\delta(t)$ belongs to $\mathfrak{g}$ for all sufficiently small $t$, we conclude (as in the proof of Theorem 3.20) that $\delta^{\prime}(0)=\gamma^{\prime}(0)$ belongs to $\mathfrak{g}$.

Corollary 3.47. If $G$ is a connected matrix Lie group, every element $A$ of $G$ can be written in the form

$$
A=e^{X_{1}} e^{X_{2}} \cdots e^{X_{m}}
$$

for some $X_{1}, X_{2}, \ldots, X_{m}$ in $\mathfrak{g}$.\\
Even if $G$ is connected, it is in general not possible to write every $A \in G$ as single exponential, $A=\exp X$, with $X \in \mathfrak{g}$. (See Example 3.41.) We begin with a simple analytic lemma.

Lemma 3.48. Suppose $A:[a, b] \rightarrow \mathrm{GI}(n ; \mathbb{C})$ is a continuous map. Then for all $\varepsilon>0$ there exists $\delta>0$ such that if $s, t \in[a, b]$ satisfy $|s-t|<\delta$, then

$$
\left\|A(s) A(t)^{-1}-I\right\|<\varepsilon .
$$

Proof. We note that

$$
\begin{aligned}
\left\|A(s) A(t)^{-1}-I\right\| & =\left\|(A(s)-A(t)) A(t)^{-1}\right\| \\
& \leq\|A(s)-A(t)\|\left\|A(t)^{-1}\right\|
\end{aligned}
$$

Since $[a, b]$ is compact and the map $t \mapsto\left\|A(t)^{-1}\right\|$ is continuous, there is a constant $C$ such that $\left\|A(t)^{-1}\right\| \leq C$ for all $t \in[a, b]$. Furthermore, since $[a, b]$ is compact, Theorem 4.19 in [Rud1] tells us that the map $A$ is actually uniformly continuous on $[a, b]$. Thus, for any $\varepsilon>0$, there exists $\delta>0$ such that when $|s-t|<\delta$, we have $\|A(s)-A(t)\|<\varepsilon / C$. Thus, in light of (3.20), we have the desired $\delta$.

Proof of Corollary 3.47. Let $V_{\varepsilon}$ be as in Theorem 3.42. For any $A \in G$, choose a continuous path $A:[0,1] \rightarrow G$ with $A(0)=I$ and $A(1)=A$. By Lemma 3.48, we can find some $\delta>0$ such that if $|s-t|<\delta$, then $A(s) A(t)^{-1} \in V_{\varepsilon}$. Now divide $[0,1]$ into $m$ pieces, where $1 / m<\delta$. Then for $j=1, \ldots, m$, we see that $A((j-1) / m)^{-1} A(j / m)$ belongs to $V_{\varepsilon}$, so that

$$
A((j-1) / m)^{-1} A(j / m)=e^{X_{j}}
$$

for some elements $X_{1}, \ldots, X_{m}$ of $\mathfrak{g}$. Thus,

$$
\begin{aligned}
A & =A(0)^{-1} A(1) \\
& =A(0)^{-1} A(1 / m) A(1 / m)^{-1} A(2 / m) \cdots A((m-1) / m)^{-1} A(1) \\
& =e^{X_{1}} e^{X_{2}} \cdots e^{X_{m}}
\end{aligned}
$$

as claimed.\\
Corollary 3.49. Let $G$ and $H$ be matrix Lie groups with Lie algebras $\mathfrak{g}$ and $\mathfrak{h}$, respectively, and assume $G$ is connected. Suppose $\Phi_{1}$ and $\Phi_{2}$ are Lie group homomorphisms of $G$ into $H$ and that $\phi_{1}$ and $\phi_{2}$ be the associated Lie algebra homomorphisms of $\mathfrak{g}$ into $\mathfrak{h}$. Then if $\phi_{1}=\phi_{2}$, we have $\Phi_{1}=\Phi_{2}$.

Proof. Let $g$ be any element of $G$. Since $G$ is connected, Corollary 3.47 tells us that every element of $G$ can be written as $e^{X_{1}} e^{X_{2}} \cdots e^{X_{m}}$, with $X_{j} \in \mathfrak{g}$. Thus,

$$
\begin{aligned}
\Phi_{1}\left(e^{X_{1}} \cdots e^{X_{m}}\right) & =e^{\phi_{1}\left(X_{1}\right)} \cdots e^{\phi_{1}\left(X_{m}\right)} \\
& =e^{\phi_{2}\left(X_{1}\right)} \cdots e^{\phi_{2}\left(X_{m}\right)} \\
& =\Phi_{2}\left(e^{X_{1}} \cdots e^{X_{m}}\right)
\end{aligned}
$$

as claimed.\\
Corollary 3.50. Every continuous homomorphism between two matrix Lie groups is smooth.

Proof. For all $g \in G$, we write nearby elements $h \in G$ as $h=g e^{X}$, with $X \in \mathfrak{g}$, so that

$$
\Phi(h)=\Phi(g) \Phi\left(e^{X}\right)=\Phi(g) e^{\phi(X)} .
$$

This relation says that in the exponential coordinates near $g$, the map $\Phi$ is a composition of a linear map, the exponential map, and multiplication on the left by $\Phi(g)$, all of which are smooth. This shows that $\Phi$ is smooth near $g$.

Corollary 3.51. If $G$ is a connected matrix Lie group and the Lie algebra $\mathfrak{g}$ of $G$ is commutative, then $G$ is commutative.

This result is a partial converse to Proposition 3.22.\\
Proof. Since $\mathfrak{g}$ is commutative, any two elements of $G$, when written as in Corollary 3.47, will commute.

Corollary 3.52. If $G \subset M_{n}(\mathbb{C})$ is a matrix Lie group, the identity component $G_{0}$ of $G$ is a closed subgroup of $\mathrm{GI}(n ; \mathbb{C})$ and thus a matrix Lie group. Furthermore, the Lie algebra of $G_{0}$ is the same as the Lie algebra of $G$.

Proof. Suppose that $\left\langle A_{m}\right\rangle$ is a sequence in $G_{0}$ converging to some $A \in \mathrm{GL}(n ; \mathbb{C})$. Then certainly $A \in G$, since $G$ is closed. Furthermore, $A_{m} A^{-1}$ lies in $G$ for all $m$ and $A_{m} A^{-1} \rightarrow I$ as $m \rightarrow \infty$. If $m$ is large enough, Theorem 3.42 tells us that $A_{m} A^{-1}=e^{X}$ for some $X \in \mathfrak{g}$, so that $A=e^{-X} A_{m}$. Since $A_{m} \in G_{0}$, there is a path joining $I$ to $A_{m}$ in $G$. But we can then join $A_{m}$ to $e^{-X} A_{m}=A$ by the path $e^{-t X} A_{m}$, $0 \leq t \leq 1$. By combining these two paths, we can join $I$ to $A$ in $G$, showing that $A$ belongs to $G_{0}$.

Now, since $G_{0} \subset G$, the Lie algebra of $G_{0}$ is contained in the Lie algebra of $G$. In the other direction, if $e^{t X}$ lies in $G$ for all $t$, then it actually lies in $G_{0}$, since any point $e^{t_{0} X}$ on the curve $e^{t X}$ can be connected to $I$ in $G$, using the curve $e^{t X}$ itself, for $0 \leq t \leq t_{0}$.

\subsection*{Exercises}
\begin{enumerate}
  \item (a) If $\mathfrak{g}$ is a Lie algebra, show that the center of $\mathfrak{g}$ is an ideal in $\mathfrak{g}$.\\
(b) If $\mathfrak{g}$ and $\mathfrak{h}$ are Lie algebras and $\phi: \mathfrak{g} \rightarrow \mathfrak{h}$ is a Lie algebra homomorphism, show that the kernel of $\phi$ is an ideal in $\mathfrak{g}$.
  \item Classify up to isomorphism all one-dimensional and two-dimensional real Lie algebras. There is one isomorphism class of one-dimensional algebras and two isomorphism classes of two-dimensional algebras.
  \item Let $\mathfrak{g}$ denote the space of $n \times n$ upper triangular matrices with zeros on the diagonal. Show that $\mathfrak{g}$ is a nilpotent Lie algebra under the bracket given by $[X, Y]=X Y-Y X$.
  \item Give an example of a matrix Lie group $G$ and a matrix $X$ such that $e^{X} \in G$, but $X \notin \mathfrak{g}$.
  \item If $G_{1} \subset \mathrm{GL}\left(n_{1} ; \mathbb{C}\right)$ and $G_{2} \subset \mathrm{GL}\left(n_{2} ; \mathbb{C}\right)$ are matrix Lie groups and $G_{1} \times G_{2}$ is their direct product (regarded as a subgroup of $\mathrm{GL}\left(n_{1}+n_{2} ; \mathbb{C}\right)$ in the obvious way), show that the Lie algebra of $G_{1} \times G_{2}$ is isomorphic to $\mathfrak{g}_{1} \oplus \mathfrak{g}_{2}$.
  \item Let $G$ and $H$ be matrix Lie groups with $H \subset G$, so that the Lie algebra $\mathfrak{h}$ of $H$ is a subalgebra of the Lie algebra $\mathfrak{g}$ of $G$.\\
(a) Show that if $H$ is a normal subgroup of $G$, then $\mathfrak{h}$ is an ideal in $\mathfrak{g}$.\\
(b) Show that if $G$ and $H$ are connected and $\mathfrak{h}$ is an ideal in $\mathfrak{g}$, then $H$ is a normal subgroup of $G$.
  \item Suppose $G \subset \mathrm{GL}(n ; \mathbb{C})$ is a matrix Lie group with Lie algebra $\mathfrak{g}$. Suppose that $A$ is in $G$ and that $\|A-I\|<1$, so that the power series for $\log A$ is convergent. Is it necessarily the case that $\log A$ is in $\mathfrak{g}$ ? Prove or give a counterexample.
  \item Show that two isomorphic matrix Lie groups have isomorphic Lie algebras.
  \item Write out explicitly the general form of a $4 \times 4$ real matrix in $\operatorname{so}(3 ; 1)$ (see Proposition 3.25).
  \item Show that there is an invertible linear map $\phi: \operatorname{su}(2) \rightarrow \mathbb{R}^{3}$ such that $\phi([X, Y])=\phi(X) \times \phi(Y)$ for all $X, Y \in \operatorname{su}(2)$, where $\times$ denotes the cross product (vector product) on $\mathbb{R}^{3}$.
  \item Show that $\operatorname{su}(2)$ and $\operatorname{sl}(2 ; \mathbb{R})$ are not isomorphic Lie algebras, even though $\mathrm{su}(2)_{\mathbb{C}} \cong \mathrm{sl}(2 ; \mathbb{R})_{\mathbb{C}} \cong \mathrm{sl}(2 ; \mathbb{C})$.
  \item Show the groups $\operatorname{SU}(2)$ and $\operatorname{SO}(3)$ are not isomorphic, even though the associated Lie algebras $\operatorname{su}(2)$ and $\operatorname{so}(3)$ are isomorphic.
  \item Let $G$ be a matrix Lie group and let $\mathfrak{g}$ be its Lie algebra. For each $A \in G$, show that $\mathrm{Ad}_{A}$ is a Lie algebra automorphism of $\mathfrak{g}$.
  \item Let $X$ and $Y$ be $n \times n$ matrices. Show by induction that
\end{enumerate}

$$
\left(\operatorname{ad}_{X}\right)^{m}(Y)=\sum_{k=0}^{m}\binom{m}{k} X^{k} Y(-X)^{m-k},
$$

where

$$
\left(\operatorname{ad}_{X}\right)^{m}(Y)=\underbrace{[X, \cdots[X,[X, Y]] \cdots]}_{m} .
$$

Now, show by direct computation that

$$
e^{\mathrm{ad}_{X}}(Y)=\operatorname{Ad}_{e^{X}}(Y)=e^{X} Y e^{-X}
$$

Assume it is legal to multiply power series term by term. (This result was obtained indirectly in Proposition 3.35.)\\
Hint: Use Pascal's triangle.\\
15. If $G$ is a matrix Lie group, a component of $G$ is the collection of all points that can be connected to a fixed $A \in G$ by a continuous path in $G$. Show that if $G$ is compact, $G$ has only finitely many components.

Hint: Suppose $G$ had infinitely many components and consider a sequence $A_{j}$ with each element of the sequence in a different component. Extract a convergent subsequence and $B_{k}=A_{j_{k}}$ and consider $B_{k}^{-1} B_{k+1}$.\\
16. Suppose that $G$ is a connected, commutative matrix Lie group with Lie algebra $\mathfrak{g}$. Show that the exponential map for $G$ maps $\mathfrak{g}$ onto $G$.\\
17. Suppose $G$ is a connected matrix Lie group with Lie algebra $\mathfrak{g}$ and that $A$ is an element of $G$. Show that $A$ belongs to the center of $G$ if and only if $\operatorname{Ad}_{A}(X)= X$ for all $X \in \mathfrak{g}$.\\
18. Show that the exponential map from the Lie algebra of the Heisenberg group to the Heisenberg group is one-to-one and onto.\\
19. Show that the exponential map from $\mathrm{u}(n)$ to $\mathrm{U}(n)$ is onto, but not one-to-one. Hint: Every unitary matrix has an orthonormal basis of eigenvectors.\\
20. Suppose $X$ is a $2 \times 2$ real matrix with trace zero, and assume $X$ has a nonreal eigenvalue.\\
(a) Show that the eigenvalues of $X$ must be of the form $i a,-i a$ with $a$ a nonzero real number.\\
(b) Show that the corresponding eigenvectors of $X$ can be chosen to be complex conjugates of each other, say, $v$ and $\bar{v}$.\\
(c) Show that there exists an invertible real matrix $C$ such that

$$
X=C\left(\begin{array}{rr}
0 & a \\
-a & 0
\end{array}\right) C^{-1} .
$$

Hint: Use $v$ and $\bar{v}$ to construct a real basis for $\mathbb{R}^{2}$, and determine the matrix $X$ in this basis.\\
21. Suppose $A$ is a $2 \times 2$ real matrix with determinant one, and assume $A$ has a nonreal eigenvalue. Show that there exists a real number $\theta$ that is not an integer multiple of $\pi$ and an invertible real matrix $C$ such that

$$
A=C\left(\begin{array}{rr}
\cos \theta & \sin \theta \\
-\sin \theta & \cos \theta
\end{array}\right) C^{-1} .
$$

\begin{enumerate}
  \setcounter{enumi}{21}
  \item Show that the image of the exponential map for $\operatorname{SL}(2 ; \mathbb{R})$ consists of precisely those matrices $A \in \mathrm{SL}(2 ; \mathbb{R})$ such that trace $(A)>-2$, together with the matrix $-I$ (which has trace -2 ). To do this, consider the possibilities for the eigenvalues of a matrix in the Lie algebra $\operatorname{sI}(2 ; \mathbb{R})$ and in the group $\operatorname{SL}(2 ; \mathbb{R})$. In the Lie algebra, show that the eigenvalues are of the form $(\lambda,-\lambda)$ or $(i \lambda,-i \lambda)$, with $\lambda$ real. In the group, show that the eigenvalues are of the form $(a, 1 / a)$ or $(-a,-1 / a)$, with $a$ real and positive, or of the form $\left(e^{i \theta}, e^{-i \theta}\right)$, with $\theta$ real. The case of a repeated eigenvalue ( $(0,0)$ in the Lie algebra and $(1,1)$ or $(-1,-1)$ in the group) will have to be treated separately using the Jordan canonical form (Sect. A.4).
\end{enumerate}

Hint: You may assume that if a real matrix $X$ has real eigenvalues, then $X$ is similar over the reals to its Jordan canonical form. Then use the two previous exercises.

\section*{Chapter 4 Basic Representation Theory}
\subsection*{Representations}
If $V$ is a finite-dimensional real or complex vector space, let $\mathrm{GL}(V)$ denote the group of invertible linear transformations of $V$. If we choose a basis for $V$, we can identify $\mathrm{GL}(V)$ with $\mathrm{GL}(n ; \mathbb{R})$ or $\mathrm{GL}(n ; \mathbb{C})$. Any such identification gives rise to a topology on $\mathrm{GL}(V)$, which is easily seen to be independent of the choice of basis. With this discussion in mind, we think of $\mathrm{GL}(V)$ as a matrix Lie group. Similarly, we let $\operatorname{gl}(V)=\operatorname{End}(V)$ denote the space of all linear operators from $V$ to itself, which forms a Lie algebra under the bracket $[X, Y]=X Y-Y X$.

Definition 4.1. Let $G$ be a matrix Lie group. A finite-dimensional complex representation of $G$ is a Lie group homomorphism

$$
\Pi: G \rightarrow \mathrm{GL}(V),
$$

where $V$ is a finite-dimensional complex vector space (with $\operatorname{dim}(V) \geq 1$ ). A finitedimensional real representation of $G$ is a Lie group homomorphism $\Pi$ of $G$ into $\mathrm{GL}(V)$, where $V$ is a finite-dimensional real vector space.

If $\mathfrak{g}$ is a real or complex Lie algebra, then a finite-dimensional complex representation of $\mathfrak{g}$ is a Lie algebra homomorphism $\pi$ of $\mathfrak{g}$ into $\mathrm{gl}(V)$, where $V$ is a finite-dimensional complex vector space. If $\mathfrak{g}$ is a real Lie algebra, then a finite-dimensional real representation of $\mathfrak{g}$ is a Lie algebra homomorphism $\pi$ of $\mathfrak{g}$ into $\mathrm{gl}(V)$.

If $\Pi$ or $\pi$ is a one-to-one homomorphism, the representation is called faithful.\\
One should think of a representation as a linear action of a group or Lie algebra on a vector space (since, say, to every $g \in G$, there is associated an operator $\Pi(g)$, which acts on the vector space $V)$. If the homomorphism $\Pi: G \rightarrow \mathrm{GL}(V)$ is fixed, we will occasionally use the alternative notation

$$
g \cdot V
$$

in place $\Pi(g) v$. We will often use terminology such as "Let $\Pi$ be a representation of $G$ acting on the space $V$."

If a representation $\Pi$ is a faithful representation of a matrix Lie group $G$, then $\{\Pi(A) \mid A \in G\}$ is a group of matrices that is isomorphic to the original group $G$. Thus, $\Pi$ allows us to represent $G$ as a group of matrices. This is the motivation for the term "representation." (Of course, we still call $\Pi$ a representation even if it is not faithful.) Despite the origin of the term, the goal of representation theory is not simply to represent a group as a group of matrices. After all, the groups we study in this book are already matrix groups! Rather, the goal is to determine (up to isomorphism) all the ways a fixed group can act as a group of matrices.

Linear actions of groups on vector spaces arise naturally in many branches of both mathematics and physics. A typical example would be a linear differential equation in three-dimensional space which has rotational symmetry, such as the equations that describe the energy states of a hydrogen atom in quantum mechanics. Since the equation is rotationally invariant, the space of solutions is invariant under rotations and thus constitutes a representation the rotation group $\mathrm{SO}(3)$. The representation theory of SO (3) (or of its Lie algebra) is very helpful in narrowing down what the space of solutions can be. See, for example, Chapter 18 in [Hall].

Definition 4.2. Let $\Pi$ be a finite-dimensional real or complex representation of a matrix Lie group $G$, acting on a space $V$. A subspace $W$ of $V$ is called invariant if $\Pi(A) w \in W$ for all $w \in W$ and all $A \in G$. An invariant subspace $W$ is called nontrivial if $W \neq\{0\}$ and $W \neq V$. A representation with no nontrivial invariant subspaces is called irreducible.

The terms invariant, nontrivial, and irreducible are defined analogously for representations of Lie algebras.

Even if $\mathfrak{g}$ is a real Lie algebra, we will consider mainly complex representations of $\mathfrak{g}$. It should be emphasized that if we are speaking about complex representations of a real Lie algebra $\mathfrak{g}$ acting on a complex subspace $V$, an invariant subspace $W$ is required to be a complex subspace of $V$.

Definition 4.3. Let $G$ be a matrix Lie group, let $\Pi$ be a representation of $G$ acting on the space $V$, and let $\Sigma$ be a representation of $G$ acting on the space $W$. A linear map $\phi: V \rightarrow W$ is called an intertwining map of representations if

$$
\phi(\Pi(A) v)=\Sigma(A) \phi(v)
$$

for all $A \in G$ and all $v \in V$. The analogous property defines intertwining maps of representations of a Lie algebra.

If $\phi$ is an intertwining map of representations and, in addition, $\phi$ is invertible, then $\phi$ is said to be an isomorphism of representations. If there exists an isomorphism between $V$ and $W$, then the representations are said to be isomorphic.

If we use the "action" notation of (4.1), the defining property of an intertwining map may be written as

$$
\phi(A \cdot v)=A \cdot \phi(v)
$$

for all $A \in G$ and $v \in V$. That is to say, $\phi$ should commute with the action of $G$. A typical problem in representation theory is to determine, up to isomorphism, all of the irreducible representations of a particular group or Lie algebra. In Sect. 4.6, we will determine all the finite-dimensional complex irreducible representations of the Lie algebra $\mathrm{SI}(2 ; \mathbb{C})$.

After identifying $\mathrm{GL}(V)$ with $\mathrm{GL}(n ; \mathbb{R})$ or $\mathrm{GL}(n ; \mathbb{C})$, Theorem 3.28 has the following consequence.

Proposition 4.4. Let $G$ be a matrix Lie group with Lie algebra $\mathfrak{g}$ and let $\Pi$ be a (finite-dimensional real or complex) representation of $G$, acting on the space $V$. Then there is a unique representation $\pi$ of $\mathfrak{g}$ acting on the same space such that

$$
\Pi\left(e^{X}\right)=e^{\pi(X)}
$$

for all $X \in \mathfrak{g}$. The representation $\pi$ can be computed as

$$
\pi(X)=\left.\frac{d}{d t} \Pi\left(e^{t X}\right)\right|_{t=0}
$$

and satisfies

$$
\pi\left(A X A^{-1}\right)=\Pi(A) \pi(X) \Pi(A)^{-1}
$$

for all $X \in \mathfrak{g}$ and all $A \in G$.\\
Given a matrix Lie group $G$ with Lie algebra $\mathfrak{g}$, we may ask whether every representation $\pi$ of $\mathfrak{g}$ comes from a representation $\Pi$ of $G$. As it turns out, this is not true in general, but is true if $G$ is simply connected. See Sect. 4.7 for examples of this phenomenon and Sect. 5.7 for the general result.

\section*{Proposition 4.5.}
\begin{enumerate}
  \item Let $G$ be a connected matrix Lie group with Lie algebra $\mathfrak{g}$. Let $\Pi$ be a representation of $G$ and $\pi$ the associated representation of $\mathfrak{g}$. Then $\Pi$ is irreducible if and only if $\pi$ is irreducible.
  \item Let $G$ be a connected matrix Lie group, let $\Pi_{1}$ and $\Pi_{2}$ be representations of $G$, and let $\pi_{1}$ and $\pi_{2}$ be the associated Lie algebra representations. Then $\pi_{1}$ and $\pi_{2}$ are isomorphic if and only if $\Pi_{1}$ and $\Pi_{2}$ are isomorphic.
\end{enumerate}

Proof. For Point 1, suppose first that $\Pi$ is irreducible. We then want to show that $\pi$ is irreducible. So, let $W$ be a subspace of $V$ that is invariant under $\pi(X)$ for all $X \in \mathfrak{g}$. We want to show that $W$ is either $\{0\}$ or $V$. Now, suppose $A$ is an element\\
of $G$. Since $G$ is assumed connected, Corollary 3.47 tells us that $A$ can be written as $A=e^{X_{1}} \cdots e^{X_{m}}$ for some $X_{1}, \ldots, X_{m}$ in $\mathfrak{g}$. Since $W$ is invariant under $\pi\left(X_{j}\right)$ it will also be invariant under $\exp \left(\pi\left(X_{j}\right)\right)=I+\pi\left(X_{j}\right)+\pi\left(X_{j}\right)^{2} / 2+\cdots$ and, hence, under

$$
\begin{aligned}
\Pi(A) & =\Pi\left(e^{X_{1}} \cdots e^{X_{m}}\right)=\Pi\left(e^{X_{1}}\right) \cdots \Pi\left(e^{X_{m}}\right) \\
& =e^{\pi\left(X_{1}\right)} \cdots e^{\pi\left(X_{m}\right)} .
\end{aligned}
$$

Since $\Pi$ is irreducible and $W$ is invariant under each $\Pi(A), W$ must be either $\{0\}$ or $V$. This shows that $\pi$ is irreducible.

In the other direction, assume that $\pi$ is irreducible and that $W$ is an invariant subspace for $\Pi$. Then $W$ is invariant under $\Pi(\exp t X)$ for all $X \in \mathfrak{g}$ and, hence, under

$$
\pi(X)=\left.\frac{d}{d t} \Pi\left(e^{t X}\right)\right|_{t=0}
$$

Thus, since $\pi$ is irreducible, $W$ is $\{0\}$ or $V$, and we conclude that $\Pi$ is irreducible. This establishes Point 1 of the proposition.

Point 2 of the proposition is similar and is left as an exercise to the reader (Exercise 1).

Proposition 4.6. Let $\mathfrak{g}$ be a real Lie algebra and $\mathfrak{g}_{\mathbb{C}}$ its complexification. Then every finite-dimensional complex representation $\pi$ of $\mathfrak{g}$ has a unique extension to a complex-linear representation of $\mathfrak{g}_{\mathbb{C}}$, also denoted $\pi$. Furthermore, $\pi$ is irreducible as a representation of $\mathfrak{g}_{\mathbb{C}}$ if and only if it is irreducible as a representation of $\mathfrak{g}$.

Of course, the extension of $\pi$ to $\mathfrak{g}_{\mathbb{C}}$ is given by $\pi(X+i Y)=\pi(X)+i \pi(Y)$ for all $X, Y \in \mathfrak{g}$.

Proof. The existence and uniqueness of the extension follow from Proposition 3.39. The claim about irreducibility holds because a complex subspace $W$ of $V$ is invariant under $\pi(X+i Y)$, with $X$ and $Y$ in $\mathfrak{g}$, if and only if it is invariant under the operators $\pi(X)$ and $\pi(Y)$. Thus, the representation of $\mathfrak{g}$ and its extension to $\mathfrak{g}_{\mathbb{C}}$ have precisely the same invariant subspaces.

Definition 4.7. If $V$ is a finite-dimensional inner product space and $G$ is a matrix Lie group, a representation $\Pi: G \rightarrow \mathrm{GL}(V)$ is unitary if $\Pi(A)$ is a unitary operator on $V$ for every $A \in G$.

Proposition 4.8. Suppose $G$ is a matrix Lie group with Lie algebra $\mathfrak{g}$. Suppose $V$ is a finite-dimensional inner product space, $\Pi$ is a representation of $G$ acting on $V$, and $\pi$ is the associated representation of $\mathfrak{g}$. If $\Pi$ is unitary, then $\pi(X)$ is skew selfadjoint for all $X \in \mathfrak{g}$. Conversely, if $G$ is connected and $\pi(X)$ is skew self-adjoint for all $X \in \mathfrak{g}$, then $\Pi$ is unitary.

In a slight abuse of notation, we will say that a representation $\pi$ of a real Lie algebra $\mathfrak{g}$ acting on a finite-dimensional inner product space is unitary if $\pi(X)$ is skew self-adjoint for all $X \in \mathfrak{g}$.

Proof. The proof is similar to the computation of the Lie algebra of the unitary group $\mathrm{U}(n)$. If $\Pi$ is unitary, then for all $X \in \mathfrak{g}$ we have

$$
\left(e^{t \pi(X)}\right)^{*}=\Pi\left(e^{t X}\right)^{*}=\Pi\left(e^{t X}\right)^{-1}=e^{-t \pi(X)}, \quad t \in \mathbb{R}
$$

so that $e^{t \pi(X)^{*}}=e^{-t \pi(X)}$. Differentiating this relation with respect to $t$ at $t=0$ gives $\pi(X)^{*}=-\pi(X)$. In the other direction, if $\pi(X)^{*}=-\pi(X)$, then the above calculation shows that $\Pi\left(e^{t X}\right)=e^{t \pi(X)}$ is unitary. If $G$ is connected, then by Corollary 3.47, every $A \in G$ is a product of exponentials, showing that $\Pi(A)$ is unitary.

\subsection*{Examples of Representations}
A matrix Lie group $G$ is, by definition, a subset of some $\operatorname{GL}(n ; \mathbb{C})$. The inclusion map of $G$ into $\operatorname{GL}(n ; \mathbb{C})$ (i.e., the map $\Pi(A)=A$ ) is a representation of $G$, called the standard representation of $G$. If $G$ happens to be contained in $\operatorname{GL}(n ; \mathbb{R}) \subset \mathrm{GL}(n ; \mathbb{C})$, then can also we can think of the standard representation as a real representation. Thus, for example, the standard representation of $\mathrm{SO}(3)$ is the one in which $\mathrm{SO}(3)$ acts in the usual way on $\mathbb{R}^{3}$ and the standard representation of $\mathrm{SU}(2)$ is the one in which $\operatorname{SU}(2)$ acts on $\mathbb{C}^{2}$ in the usual way. Similarly, if $\mathfrak{g} \subset M_{n}(\mathbb{C})$ is a Lie algebra of matrices, the map $\pi(X)=X$ is called the standard representation of $\mathfrak{g}$.

Consider the one-dimensional complex vector space $\mathbb{C}$. For any matrix Lie group $G$, we can define the trivial representation, $\Pi: G \rightarrow \mathrm{GL}(1 ; \mathbb{C})$, by the formula

$$
\Pi(A)=I
$$

for all $A \in G$. Of course, this is an irreducible representation, since $\mathbb{C}$ has no nontrivial subspaces, let alone nontrivial invariant subspaces. If $\mathfrak{g}$ is a Lie algebra, we can also define the trivial representation of $\mathfrak{g}, \pi: \mathfrak{g} \rightarrow \mathrm{gl}(1 ; \mathbb{C})$, by

$$
\pi(X)=0
$$

for all $X \in \mathfrak{g}$. This is an irreducible representation.\\
Recall the adjoint map of a group or Lie algebra, described in Definitions 3.32 and 3.7.

Definition 4.9. If $G$ is a matrix Lie group with Lie algebra $\mathfrak{g}$, the adjoint representation of $G$ is the map $\mathrm{Ad}: G \rightarrow \mathrm{GL}(\mathfrak{g})$ given by $A \mapsto \mathrm{Ad}_{A}$. Similarly, the adjoint\\
representation of a finite-dimensional Lie algebra $\mathfrak{g}$ is the map ad: $\mathfrak{g} \rightarrow \mathfrak{g l}(\mathfrak{g})$ given by $X \mapsto \operatorname{ad}_{X}$.

If $G$ is a matrix Lie group with Lie algebra $\mathfrak{g}$, then by Proposition 3.34, the Lie algebra representation associated to the adjoint representation of $G$ is the adjoint representation of $\mathfrak{g}$. Note that in the case of $\mathrm{SO}(3)$, the standard representation and the adjoint representation are both three-dimensional real representations. In fact, these two representations are isomorphic (Exercise 2).

Example 4.10. Let $V_{m}$ denote the space of homogeneous polynomials of degree $m$ in two complex variables. For each $U \in \mathrm{SU}(2)$, define a linear transformation $\Pi_{m}(U)$ on the space $V_{m}$ by the formula

$$
\left[\Pi_{m}(U) f\right](z)=f\left(U^{-1} z\right), \quad z \in \mathbb{C}^{2}
$$

Then $\Pi_{m}$ is a representation of $\operatorname{SU}(2)$.\\
Elements of $V_{m}$ have the form

$$
f\left(z_{1}, z_{2}\right)=a_{0} z_{1}^{m}+a_{1} z_{1}^{m-1} z_{2}+a_{2} z_{1}^{m-2} z_{2}^{2}+\cdots+a_{m} z_{2}^{m}
$$

with $z_{1}, z_{2} \in \mathbb{C}$ and the $a_{j}$ 's arbitrary complex constants, from which we see that $\operatorname{dim}\left(V_{m}\right)=m+1$. Explicitly, if $f$ is as in (4.3), then

$$
\left[\Pi_{m}(U) f\right]\left(z_{1}, z_{2}\right)=\sum_{k=0}^{m} a_{k}\left(U_{11}^{-1} z_{1}+U_{12}^{-1} z_{2}\right)^{m-k}\left(U_{21}^{-1} z_{1}+U_{22}^{-1} z_{2}\right)^{k}
$$

By expanding out the right-hand side of this formula, we see that $\Pi_{m}(U) f$ is again a homogeneous polynomial of degree $m$. Thus, $\Pi_{m}(U)$ actually maps $V_{m}$ into $V_{m}$.

To see that $\Pi_{m}$ is actually a representation, compute that

$$
\begin{aligned}
\Pi_{m}\left(U_{1}\right)\left[\Pi_{m}\left(U_{2}\right) f\right](z) & =\left[\Pi_{m}\left(U_{2}\right) f\right]\left(U_{1}^{-1} z\right)=f\left(U_{2}^{-1} U_{1}^{-1} z\right) \\
& =\Pi_{m}\left(U_{1} U_{2}\right) f(z)
\end{aligned}
$$

The inverse on the right-hand side of (4.2) is necessary in order to make $\Pi_{m}$ a representation. We will see in Proposition 4.11 that each $\Pi_{m}$ is irreducible and we will see in Sect. 4.6 that every finite-dimensional irreducible representation of $\operatorname{SU}(2)$ is isomorphic to one (and only one) of the $\Pi_{m}$ 's. (Of course, no two of the $\Pi_{m}$ 's are isomorphic, since they do not even have the same dimension.)

The associated representation $\pi_{m}$ of $\operatorname{su}(2)$ can be computed as

$$
\left(\pi_{m}(X) f\right)(z)=\left.\frac{d}{d t} f\left(e^{-t X} z\right)\right|_{t=0}
$$

Now, let $z(t)=\left(z_{1}(t), z_{2}(t)\right)$ be the curve in $\mathbb{C}^{2}$ defined as $z(t)=e^{-t X} z$. By the chain rule, we have

$$
\pi_{m}(X) f=\left.\frac{\partial f}{\partial z_{1}} \frac{d z_{1}}{d t}\right|_{t=0}+\left.\frac{\partial f}{\partial z_{2}} \frac{d z_{2}}{d t}\right|_{t=0} .
$$

Since $d z /\left.d t\right|_{t=0}=-X z$, so we obtain

$$
\pi_{m}(X) f=-\frac{\partial f}{\partial z_{1}}\left(X_{11} z_{1}+X_{12} z_{2}\right)-\frac{\partial f}{\partial z_{2}}\left(X_{21} z_{1}+X_{22} z_{2}\right) .
$$

We may then take the unique complex-linear extension of $\pi$ to $\operatorname{sI}(2 ; \mathbb{C}) \cong \operatorname{su}(2)_{\mathbb{C}}$, as in Proposition 3.39. This extension is given by the same formula, but with $X \in$ sl(2; C).

If $X, Y$, and $H$ are the following basis elements for $\mathrm{sI}(2 ; \mathbb{C})$ :

$$
H=\left(\begin{array}{rr}
1 & 0 \\
0 & -1
\end{array}\right) ; \quad X=\left(\begin{array}{ll}
0 & 1 \\
0 & 0
\end{array}\right) ; \quad Y=\left(\begin{array}{ll}
0 & 0 \\
1 & 0
\end{array}\right)
$$

then applying formula (4.4) gives

$$
\begin{aligned}
\pi_{m}(H) & =-z_{1} \frac{\partial}{\partial z_{1}}+z_{2} \frac{\partial}{\partial z_{2}} \\
\pi_{m}(X) & =-z_{2} \frac{\partial}{\partial z_{1}} \\
\pi_{m}(Y) & =-z_{1} \frac{\partial}{\partial z_{2}}
\end{aligned}
$$

Applying these operators to a basis element $z_{1}^{m-k} z_{2}^{k}$ for $V_{m}$ gives

$$
\begin{aligned}
\pi_{m}(H)\left(z_{1}^{m-k} z_{2}^{k}\right) & =(-m+2 k) z_{1}^{m-k} z_{2}^{k} \\
\pi_{m}(X)\left(z_{1}^{m-k} z_{2}^{k}\right) & =-(m-k) z_{1}^{m-k-1} z_{2}^{k+1} \\
\pi_{m}(Y)\left(z_{1}^{m-k} z_{2}^{k}\right) & =-k z_{1}^{m-k+1} z_{2}^{k-1}
\end{aligned}
$$

Thus, $z_{1}^{m-k} z_{2}^{k}$ is an eigenvector for $\pi_{m}(H)$ with eigenvalue $-m+2 k$, while $\pi_{m}(X)$ and $\pi_{m}(Y)$ have the effect of shifting the exponent $k$ of $z_{2}$ up or down by one. Note that since $\pi_{m}(X)$ increases the value of $k$, this operator increases the eigenvalue of $\pi_{m}(H)$ by 2 , whereas $\pi_{m}(Y)$ decreases the eigenvalue of $\pi_{m}(H)$ by 2 .


Fig. 4.1 The black dots indicate the nonzero terms for a vector $w$ in the space $V_{6}$. Applying $\pi_{6}(X)^{4}$ to $w$ gives a nonzero multiple of $z_{2}^{6}$

\includegraphics[scale=0.2, center]{2025_08_28_fc86fd4ade69f41b5842g-096}


Proposition 4.11. For each $m \geq 0$, the representation $\pi_{m}$ is irreducible.\\
Proof. It suffices to show that every nonzero invariant subspace of $V_{m}$ is equal to $V_{m}$. So, let $W$ be such a space and let $w$ be a nonzero element of $W$. Then $w$ can be written in the form

$$
w=a_{0} z_{1}^{m}+a_{1} z_{1}^{m-1} z_{2}+a_{2} z_{1}^{m-2} z_{2}^{2}+\cdots+a_{m} z_{2}^{m}
$$

with at least one of the $a_{k}$ 's being nonzero. Let $k_{0}$ be the smallest value of $k$ for which $a_{k} \neq 0$ and consider

$$
\pi_{m}(X)^{m-k_{0}} w .
$$

\section*{(See Figure 4.1.)}
Since each application of $\pi_{m}(X)$ raises the power of $z_{2}$ by $1, \pi_{m}(X)^{m-k_{0}}$ will kill all the terms in $w$ except $a_{k_{0}} z_{1}^{m-k_{0}} z_{2}^{k_{0}}$ term. On the other hand, since $\pi_{m}(X)\left(z_{1}^{m-k} z_{2}^{k}\right)$ is zero only if $k=m$, we see that $\pi_{m}(X)^{m-k_{0}} w$ is a nonzero multiple of $z_{2}^{m}$. Since $W$ is assumed invariant, $W$ must contain this multiple of $z_{2}^{m}$ and so also $z_{2}^{m}$ itself. Now, for $0 \leq k \leq m$, it follows from (4.5) that $\pi_{m}(Y)^{k} z_{2}^{m}$ is a nonzero multiple of $z_{1}^{k} z_{2}^{m-k}$. Therefore, $W$ must also contain $z_{1}^{k} z_{2}^{m-k}$ for all $0 \leq k \leq m$. Since these elements form a basis for $V_{m}$, we see that $W=V_{m}$, as desired.

\subsection*{New Representations from Old}
One way of generating representations is to take some representations one knows and combine them in some fashion. In this section, we will consider three standard methods of obtaining new representations from old, namely direct sums of representations, tensor products of representations, and dual representations.

\subsubsection*{Direct Sums}
Definition 4.12. Let $G$ be a matrix Lie group and let $\Pi_{1}, \Pi_{2}, \ldots, \Pi_{m}$ be representations of $G$ acting on vector spaces $V_{1}, V_{2}, \ldots, V_{m}$. Then the direct sum of $\Pi_{1}, \Pi_{2}, \ldots, \Pi_{m}$ is a representation $\Pi_{1} \oplus \cdots \oplus \Pi_{m}$ of $G$ acting on the space $V_{1} \oplus \cdots \oplus V_{m}$, defined by

$$
\left[\Pi_{1} \oplus \cdots \oplus \Pi_{m}(A)\right]\left(v_{1}, \ldots, v_{m}\right)=\left(\Pi_{1}(A) v_{1}, \ldots, \Pi_{m}(A) v_{m}\right)
$$

for all $A \in G$.

Similarly, if $\mathfrak{g}$ is a Lie algebra, and $\pi_{1}, \pi_{2}, \ldots, \pi_{m}$ are representations of $\mathfrak{g}$ acting on $V_{1}, V_{2}, \ldots, V_{m}$, then we define the direct sum of $\pi_{1}, \pi_{2}, \ldots, \pi_{m}$, acting on $V_{1} \oplus \cdots \oplus V_{m}$ by

$$
\left[\pi_{1} \oplus \cdots \oplus \pi_{m}(X)\right]\left(v_{1}, \ldots, v_{m}\right)=\left(\pi_{1}(X) v_{1}, \ldots, \pi_{m}(X) v_{m}\right)
$$

for all $X \in \mathfrak{g}$.\\
It is straightforward to check that, say, $\Pi_{1} \oplus \cdots \oplus \Pi_{m}$ is really a representation of $G$.

\subsubsection*{Tensor Products}
Let $U$ and $V$ be finite-dimensional real or complex vector spaces. We wish to define the tensor product of $U$ and $V$, which will be a new vector space $U \otimes V$ "built" out of $U$ and $V$. We will discuss the idea of this first and then give the precise definition.

We wish to consider a formal "product" of an element $u$ of $U$ with an element $v$ of $V$, denoted $u \otimes v$. The space $U \otimes V$ is then the space of linear combinations of such products, that is, the space of elements of the form

$$
a_{1} u_{1} \otimes v_{1}+a_{2} u_{2} \otimes v_{2}+\cdots+a_{n} u_{n} \otimes v_{n} .
$$

Of course, if " $\otimes$ " is to be interpreted as a product, then it should be bilinear:

$$
\begin{aligned}
& \left(u_{1}+a u_{2}\right) \otimes v=u_{1} \otimes v+a u_{2} \otimes v, \\
& u \otimes\left(v_{1}+a v_{2}\right)=u \otimes v_{1}+a u \otimes v_{2} .
\end{aligned}
$$

We do not assume that the product is commutative. That is to say, even if $U=V$ so that $U \otimes V$ and $V \otimes U$ are the same space, $u \otimes v$ will not, in general, equal $v \otimes u$.

Now, if $e_{1}, e_{2}, \ldots, e_{n}$ is a basis for $U$ and $f_{1}, f_{2}, \ldots, f_{m}$ is a basis for $V$, then, using bilinearity, it is easy to see that any element of the form (4.6) can be written as a linear combination of the elements $e_{j} \otimes f_{k}$. In fact, it seems reasonable to expect that $\left\{e_{j} \otimes f_{k} \mid 1 \leq j \leq n, 1 \leq k \leq m\right\}$ should be a basis for the space $U \otimes V$. This, in fact, turns out to be the case.\\
Definition 4.13. If $U$ and $V$ are finite-dimensional real or complex vector spaces, then a tensor product of $U$ with $V$ is a vector space $W$, together with a bilinear map $\phi: U \times V \rightarrow W$ with the following property: If $\psi$ is any bilinear map of $U \times V$ into a vector space $X$, there exists a unique linear map $\tilde{\psi}$ of $W$ into $X$ such that the following diagram commutes:\\
\includegraphics[scale=0.2, center]{2025_08_28_fc86fd4ade69f41b5842g-099(1)}

Note that the bilinear map $\psi$ from $U \times V$ into $X$ turns into the linear map $\tilde{\psi}$ of $W$ into $X$. This is one of the points of tensor products: Bilinear maps on $U \times V$ turn into linear maps on $W$.

Theorem 4.14. If $U$ and $V$ are any finite-dimensional real or complex vector spaces, then a tensor product ( $W, \phi$ ) exists. Furthermore, ( $W, \phi$ ) is unique up to canonical isomorphism. That is, if ( $W_{1}, \phi_{1}$ ) and ( $W_{2}, \phi_{2}$ ) are two tensor products, then there exists a unique vector space isomorphism $\Phi: W_{1} \rightarrow W_{2}$ such that the following diagram commutes:\\
\includegraphics[scale=0.2, center]{2025_08_28_fc86fd4ade69f41b5842g-099}

Suppose that ( $W, \phi$ ) is a tensor product and that $e_{1}, e_{2}, \ldots, e_{n}$ is a basis for $U$ and $f_{1}, f_{2}, \ldots, f_{m}$ is a basis for $V$. Then

$$
\left\{\phi\left(e_{j}, f_{k}\right) \mid 1 \leq j \leq n, 1 \leq k \leq m\right\}
$$

is a basis for $W$.\\
In particular, $\operatorname{dim}(U \otimes V)=(\operatorname{dim} U)(\operatorname{dim} V)$.\\
Proof. Exercise 7.\\
Notation 4.15 Since the tensor product of $U$ and $V$ is essentially unique, we will let $U \otimes V$ denote an arbitrary tensor product space and we will write $u \otimes v$ instead of $\phi(u, v)$. In this notation, Theorem 4.14 says that

$$
\left\{e_{j} \otimes f_{k} \mid 1 \leq j \leq n, 1 \leq k \leq m\right\}
$$

is a basis for $U \otimes V$, as expected.\\
The defining property of $U \otimes V$ is called the universal property of tensor products. To understand the significance of this property, suppose we want to define a linear map $T$ from $U \otimes V$ into some other space. We could try to define $T$ using bases for $U$ and $V$, but then we would have to worry about whether $T$ depends on the choice of basis. Instead, we could try to define $T$ on elements of the form $u \otimes v$, with $u \in U$ and $v \in V$. While it follows from Theorem 4.14 that elements of this form span $U \otimes V$, the decomposition of an element of $U \otimes V$ as a linear combination of elements of the form $u \otimes v$ is far from unique. Thus, if we wish to define $T$ on such elements, we have to worry whether $T$ is well defined.

This is where the universal property comes in. If $\psi(u, v)$ is any bilinear expression in $(u, v)$, the universal property says precisely that there is a unique linear map $T(=\tilde{\psi})$ such that

$$
T(u \otimes v)=\psi(u, v) .
$$

Thus, we can construct a well-defined linear map $T$ on $U \otimes V$ simply by defining it on elements of the form $u \otimes v$, provided that our definition of $T(u \otimes v)$ is bilinear in $u$ and $v$. The following result is an application of this line of reasoning.

Proposition 4.16. Let $U$ and $V$ be finite-dimensional real or complex vector spaces. Let $A: U \rightarrow U$ and $B: V \rightarrow V$ be linear operators. Then there exists a unique linear operator from $U \otimes V$ to $U \otimes V$, denoted $A \otimes B$, such that

$$
(A \otimes B)(u \otimes v)=(A u) \otimes(B v)
$$

for all $u \in U$ and $v \in V$. If $A_{1}$ and $A_{2}$ are linear operators on $U$ and $B_{1}$ and $B_{2}$ are linear operators on $V$, then

$$
\left(A_{1} \otimes B_{1}\right)\left(A_{2} \otimes B_{2}\right)=\left(A_{1} A_{2}\right) \otimes\left(B_{1} B_{2}\right) .
$$

Proof. Define a map $\psi$ from $U \times V$ into $U \otimes V$ by

$$
\psi(u, v)=(A u) \otimes(B v) .
$$

Since $A$ and $B$ are linear and since $\otimes$ is bilinear, $\psi$ is a bilinear map of $U \times V$ into $U \otimes V$. By the universal property, there is an associated linear map $\tilde{\psi}: U \otimes V \rightarrow U \otimes V$ such that

$$
\tilde{\psi}(u \otimes v)=\psi(u, v)=(A u) \otimes(B v) .
$$

Thus, $\tilde{\psi}$ is the desired map $A \otimes B$. An elementary calculation shows that the identity (4.7) holds on elements of the form $u \otimes v$. Since, by Theorem 4.14, such elements span $U \otimes V$, the identity holds in general.

We are now ready to define tensor products of representations. There are two different approaches to this, both of which are important. The first approach starts with a representation of a group $G$ acting on a space $U$ and a representation of another group $H$ acting on a space $V$ and produces a representation of the product group $G \times H$ acting on the space $U \otimes V$. The second approach starts with two different representations of the same group $G$, acting on spaces $U$ and $V$, and produces a representation of $G$ acting on $U \otimes V$. Both of these approaches can be adapted to apply to Lie algebras.

Definition 4.17. Let $G$ and $H$ be matrix Lie groups. Let $\Pi_{1}$ be a representation of $G$ acting on a space $U$ and let $\Pi_{2}$ be a representation of $H$ acting on a space $V$.

Then the tensor product of $\Pi_{1}$ and $\Pi_{2}$ is a representation $\Pi_{1} \otimes \Pi_{2}$ of $G \times H$ acting on $U \otimes V$ defined by

$$
\left(\Pi_{1} \otimes \Pi_{2}\right)(A, B)=\Pi_{1}(A) \otimes \Pi_{2}(B)
$$

for all $A \in G$ and $B \in H$.\\
Using Proposition 4.16, it is easy to check that $\Pi_{1} \otimes \Pi_{2}$ is, in fact, a representation of $G \times H$.

Now, if $G$ and $H$ are matrix Lie groups (i.e., $G$ is a closed subgroup of $\mathrm{GL}(n ; \mathbb{C})$ and $H$ is a closed subgroup of $\mathrm{GL}(m ; \mathbb{C})$ ), then $G \times H$ can be regarded in an obvious way as a closed subgroup of $\mathrm{GL}(n+m ; \mathbb{C})$. Thus, the direct product of matrix Lie groups can be regarded as a matrix Lie group. It is easy to check that the Lie algebra of $G \times H$ is isomorphic to the direct sum of the Lie algebra of $G$ and the Lie algebra of $H$.

Proposition 4.18. Let $G$ and $H$ be matrix Lie groups with Lie algebras $\mathfrak{g}$ and $\mathfrak{h}$, respectively. Let $\Pi_{1}$ and $\Pi_{2}$ be representations of $G$ and $H$, respectively, and consider the representation $\Pi_{1} \otimes \Pi_{2}$ of $G \times H$. If $\pi_{1} \otimes \pi_{2}$ denotes the associated representation of $\mathfrak{g} \oplus \mathfrak{h}$, then

$$
\left(\pi_{1} \otimes \pi_{2}\right)(X, Y)=\pi_{1}(X) \otimes I+I \otimes \pi_{2}(Y)
$$

for all $X \in \mathfrak{g}$ and $Y \in \mathfrak{h}$.\\
Proof. Suppose that $u(t)$ is a smooth curve in $U$ and $v(t)$ is a smooth curve in $V$. Then, by repeating the proof of the product rule for scalar-valued functions (or by calculating everything in a basis), we have

$$
\frac{d}{d t}(u(t) \otimes v(t))=\frac{d u}{d t} \otimes v(t)+u(t) \otimes \frac{d v}{d t}
$$

This being the case, we compute as follows:

$$
\begin{aligned}
& \left(\pi_{1} \otimes \pi_{2}\right)(X, Y)(u \otimes v) \\
& =\left.\frac{d}{d t} \Pi_{1}\left(e^{t X}\right) u \otimes \Pi_{2}\left(e^{t Y}\right) v\right|_{t=0} \\
& =\left(\left.\frac{d}{d t} \Pi_{1}\left(e^{t X}\right) u\right|_{t=0}\right) \otimes v+u \otimes\left(\left.\frac{d}{d t} \Pi_{2}\left(e^{t Y}\right) v\right|_{t=0}\right)
\end{aligned}
$$

This establishes the claimed form of $\left(\pi_{1} \otimes \pi_{2}\right)(X, Y)$ on elements of the form $u \otimes v$, which span $U \otimes V$.

The proposition motivates the following definition.

Definition 4.19. Let $\mathfrak{g}$ and $\mathfrak{h}$ be Lie algebras and let $\pi_{1}$ and $\pi_{2}$ be representations of $\mathfrak{g}$ and $\mathfrak{h}$, acting on spaces $U$ and $V$. Then the tensor product of $\pi_{1}$ and $\pi_{2}$, denoted $\pi_{1} \otimes \pi_{2}$, is a representation of $\mathfrak{g} \oplus \mathfrak{h}$ acting on $U \otimes V$, given by

$$
\left(\pi_{1} \otimes \pi_{2}\right)(X, Y)=\pi_{1}(X) \otimes I+I \otimes \pi_{2}(Y)
$$

for all $X \in \mathfrak{g}$ and $Y \in \mathfrak{h}$.\\
It is easy to check that this indeed defines a representation of $\mathfrak{g} \oplus \mathfrak{h}$. Note that if we defined $\left(\pi_{1} \otimes \pi_{2}\right)(X, Y)=\pi_{1}(X) \otimes \pi_{2}(Y)$, this would not be a representation of $\mathfrak{g} \oplus \mathfrak{h}$, since this is expression is not linear in ( $X, Y$ ).

We now define a variant of the above definitions in which we take the tensor product of two representations of the same group $G$ and regard the result as a representation of $G$ rather than of $G \times G$.

Definition 4.20. Let $G$ be a matrix Lie group and let $\Pi_{1}$ and $\Pi_{2}$ be representations of $G$, acting on spaces $V_{1}$ and $V_{2}$. Then the tensor product representation of $G$, acting on $V_{1} \otimes V_{2}$, is defined by

$$
\left(\Pi_{1} \otimes \Pi_{2}\right)(A)=\Pi_{1}(A) \otimes \Pi_{2}(A)
$$

for all $A \in G$. Similarly, if $\pi_{1}$ and $\pi_{2}$ are representations of a Lie algebra $\mathfrak{g}$, we define a tensor product representation of $\mathfrak{g}$ on $V_{1} \otimes V_{2}$ by

$$
\left(\pi_{1} \otimes \pi_{2}\right)(X)=\pi_{1}(X) \otimes I+I \otimes \pi_{2}(X) .
$$

It is easy to check that $\Pi_{1} \otimes \Pi_{2}$ and $\pi_{1} \otimes \pi_{2}$ are actually representations of $G$ and $\mathfrak{g}$, respectively. The notation is, unfortunately, ambiguous, since if $\Pi_{1}$ and $\Pi_{2}$ are representations of the same group $G$, we can regard $\Pi_{1} \otimes \Pi_{2}$ either as a representation of $G$ or as a representation of $G \times G$. We must, therefore, be careful to specify which way we are thinking about $\Pi_{1} \otimes \Pi_{2}$.

If $\Pi_{1}$ and $\Pi_{2}$ are irreducible representations of a group $G$, then $\Pi_{1} \otimes \Pi_{2}$ will typically not be irreducible when viewed as a representation of $G$. One can, then, attempt to decompose $\Pi_{1} \otimes \Pi_{2}$ as a direct sum of irreducible representations. This process is called the Clebsch-Gordan theory or, in the physics literature, "addition of angular momentum." See Exercise 12 and Appendix C for more information about this topic.

\subsubsection*{Dual Representations}
Suppose that $\pi$ is a representation of a Lie algebra $\mathfrak{g}$ acting on a finite-dimensional vector space $V$. Let $V^{*}$ denote the dual space of $V$, that is, the space of linear\\
functionals on $V$ (Sect. A.7). If $A$ is a linear operator on $V$, let $A^{t r}$ denote the dual or transpose operator on $V^{*}$, given by

$$
\left(A^{t r} \phi\right)(v)=\phi(A v)
$$

for $\phi \in V^{*}, v \in V$. If $v_{1}, \ldots, v_{n}$ is a basis for $V$, then there is a naturally associated "dual basis" $\phi_{1}, \ldots, \phi_{n}$ with the property that $\phi_{j}\left(v_{k}\right)=\delta_{j k}$. The matrix for $A^{t r}$ in the dual basis is then simply the transpose (not the conjugate transpose!) of the matrix of $A$ in the original basis. If $A$ and $B$ are linear operators on $V$, it is easily verified that

$$
(A B)^{t r}=B^{t r} A^{t r}
$$

Definition 4.21. Suppose $G$ is a matrix Lie group and $\Pi$ is a representation of $G$ acting on a finite-dimensional vector space $V$. Then the dual representation $\Pi$ * to $\Pi$ is the representation of $G$ acting on $V^{*}$ and given by

$$
\Pi^{*}(g)=\left[\Pi\left(g^{-1}\right)\right]^{t r} .
$$

If $\pi$ is a representation of a Lie algebra $\mathfrak{g}$ acting on a finite-dimensional vector space $V$, then $\pi^{*}$ is the representation of $\mathfrak{g}$ acting on $V^{*}$ and given by

$$
\pi^{*}(X)=-\pi(X)^{t r} .
$$

Using (4.8), it is easy to check that both $\Pi^{*}$ and $\pi^{*}$ are actually representations. (Here the inverse on the right-hand side of (4.9) and the minus sign on the right-hand side of (4.10) are essential.) The dual representation is also called contragredient representation.

Proposition 4.22. If $\Pi$ is a representation of a matrix Lie group $G$, then (1) $\Pi^{*}$ is irreducible if and only if $\Pi$ is irreducible and (2) $\left(\Pi^{*}\right)^{*}$ is isomorphic to $\Pi$. Similar statements apply to Lie algebra representations.

Proof. See Exercise 6.

\subsection*{Complete Reducibility}
Much of representation theory is concerned with studying irreducible representations of a group or Lie algebra. In favorable cases, knowing the irreducible representations leads to a description of all representations.

Definition 4.23. A finite-dimensional representation of a group or Lie algebra is said to be completely reducible if it is isomorphic to a direct sum of a finite number of irreducible representations.

Definition 4.24. A group or Lie algebra is said to have the complete reducibility property if every finite-dimensional representation of it is completely reducible.

As it turns out, most groups and Lie algebras do not have the complete reducibility property. Nevertheless, many interesting example groups and Lie algebra do have this property, as we will see in this section and Sect. 10.3.

Example 4.25. Let $\Pi: \mathbb{R} \rightarrow \mathrm{GL}(2 ; \mathbb{C})$ be given by

$$
\Pi(x)=\left(\begin{array}{ll}
1 & x \\
0 & 1
\end{array}\right) .
$$

Then $\Pi$ is not completely reducible.\\
Proof. Direct calculation shows that $\Pi$ is, in fact, a representation of $\mathbb{R}$. If $\left\{e_{1}, e_{2}\right\}$ is the standard basis for $\mathbb{C}^{2}$, then clearly the span of $e_{1}$ is an invariant subspace. We now claim that $\left\langle e_{1}\right\rangle$ is the only nontrivial invariant subspace for $\Pi$. To see this, suppose $V$ is a nonzero invariant subspace and suppose $V$ contains a vector not in the span of $e_{1}$, say, $v=a e_{1}+b e_{2}$ with $b \neq 0$. Then

$$
\Pi(1) v-v=b e_{1}
$$

also belongs to $V$. Thus, $e_{1}$ and $e_{2}=\left(v-a e_{1}\right) / b$ belong to $V$, showing that $V=\mathbb{C}^{2}$. We conclude, then, that $\mathbb{C}^{2}$ does not decompose as a direct sum of irreducible invariant subspaces.

Proposition 4.26. If $V$ is a completely reducible representation of a group or Lie algebra, then the following properties hold.

\begin{enumerate}
  \item For every invariant subspace $U$ of $V$, there is another invariant subspace $W$ such that $V$ is the direct sum of $U$ and $W$.
  \item Every invariant subspace of $V$ is completely reducible.
\end{enumerate}

Proof. For Point 1, suppose that $V$ decomposes as

$$
V=U_{1} \oplus U_{2} \oplus \cdots \oplus U_{k},
$$

where the $U_{j}$ 's are irreducible invariant subspaces, and that $U$ is any invariant subspace of $V$. If $U$ is all of $V$, then we can take $W=\{0\}$ and we are done. If $W \neq V$, there must be some $j_{1}$ such that $U_{j_{1}}$ is not contained in $U$. Since $U_{j_{1}}$ is irreducible, it follows that the invariant subspace $U_{j_{1}} \cap U$ must be $\{0\}$. Suppose now that $U+U_{j_{1}}=V$. If so, the sum is direct (since $U_{j_{1}} \cap U=\{0\}$ ) and we are done.

If $U+U_{j_{1}} \neq V$, there is some $j_{2}$ such that $U+U_{j_{1}}$ does not contain $U_{j_{2}}$, in which case, $\left(U+U_{j_{1}}\right) \cap U_{j_{2}}=\{0\}$. Proceeding on in the same way, we must eventually obtain some family $j_{1}, j_{2}, \ldots, j_{l}$ such that $U+U_{j_{1}}+\cdots+U_{j_{l}}=V$ and the sum is direct. Then $W:=U_{j_{1}}+\cdots+U_{j_{l}}$ is the desired complement to $U$.

For Point 2, suppose $U$ is an invariant subspace of $V$. We first establish that $U$ has the "invariant complement property" in Point 1 . Suppose, then, that $X$ is another invariant subspace of $V$ with $X \subset U$. By Point 1, we can find invariant subspace $Y$\\
such that $V=X \oplus Y$. Let $Z=Y \cap U$, which is then an invariant subspace. We want to show that $U=X \oplus Z$. For all $u \in U$, we can write $u=x+y$ with $x \in X$ and $y \in Y$. But since $X \subset U$, we have $x \in U$ and therefore $y=u-x \in U$. Thus, $y \in Z=Y \cap U$. We have shown, then, that every $u \in U$ can be written as the sum of an element of $X$ and an element of $Z$. Furthermore, $X \cap Z \subset X \cap Y=\{0\}$, so actually $U$ is the direct sum of $X$ and $Z$.

We may now easily show that $U$ is completely reducible. If $U$ is irreducible, we are done. If not, $U$ has a nontrivial invariant subspace $X$ and thus $U$ decomposes as $U=X \oplus Z$ for some invariant subspace $Z$. If $X$ and $Z$ are irreducible, we are done, and if not, we proceed on in the same way. Since $U$ is finite dimensional, this process must eventually terminate with $U$ being decomposed as a direct sum of irreducibles.

Proposition 4.27. If $G$ is a matrix Lie group and $\Pi$ is a finite-dimensional unitary representation of $G$, then $\Pi$ is completely reducible. Similarly, if $\mathfrak{g}$ is a real Lie algebra and $\pi$ is a finite-dimensional "unitary"representation of $\mathfrak{g}$ (meaning that $\pi(X)^{*}=-\pi(X)$ for all $\left.X \in \mathfrak{g}\right)$, then $\pi$ is completely reducible.

Proof. Let $V$ denote the Hilbert space on which $\Pi$ acts and let $\langle\cdot, \cdot\rangle$ denote the inner product on $V$. If $W \subset V$ be an invariant subspace, let $W^{\perp}$ be the orthogonal complement of $W$, so that $V$ is the direct sum of $W$ and $W^{\perp}$. We claim that $W^{\perp}$ is also an invariant subspace for $\Pi$ or $\pi$.

To see this, note that since $\Pi$ is unitary, $\Pi(A)^{*}=\Pi(A)^{-1}=\Pi\left(A^{-1}\right)$ for all $A \in G$. Then, for any $w \in W$ and any $v \in W^{\perp}$, we have

$$
\begin{aligned}
\langle\Pi(A) v, w\rangle & =\left\langle v, \Pi(A)^{*} w\right\rangle=\left\langle v, \Pi\left(A^{-1}\right) w\right\rangle \\
& =\left\langle v, w^{\prime}\right\rangle=0 .
\end{aligned}
$$

In the last step, we have used that $w^{\prime}=\Pi\left(A^{-1}\right) w$ is in $W$, since $W$ is invariant. This shows that $\Pi(A) v$ is orthogonal to every element of $W$, as claimed. A similar argument, with $\Pi\left(A^{-1}\right)$ replaced by $-\pi(X)$, shows that the orthogonal complement of an invariant subspace for $\pi$ is also invariant.

We have established, then, that for a unitary representation, the orthogonal complement of an invariant subspace is again invariant. Suppose now that $V$ is not irreducible. Then we can find an invariant subspace $W$ that is neither $\{0\}$ nor $V$, and we decompose $V$ as $W \oplus W^{\perp}$. Then $W$ and $W^{\perp}$ are both invariant subspaces and thus unitary representations of $G$ in their own right. Then $W$ is either irreducible or it splits as an orthogonal direct sum of invariant subspaces, and similarly for $W^{\perp}$. We continue this process, and since $V$ is finite dimensional, it cannot go on forever, and we eventually arrive at a decomposition of $V$ as a direct sum of irreducible invariant subspaces.

Theorem 4.28. If $G$ is a compact matrix Lie group, every finite dimensional representation of $G$ is completely reducible.

See also Sect. 10.3 for a similar result for semisimple Lie algebras. The argument below is sometimes called "Weyl's unitarian trick" for the role of unitarity in the\\
proof. We require a notion of integration over matrix Lie groups that is invariant under the right action of the group. One way to construct such a right-invariant integral is to construct a right-invariant measure on $G$, known as a Haar measure. It is, however, simpler to introduce the integral by means of a right-invariant differential form on $G$. (See Appendix B for a quick introduction to the notion of differential forms.)

If $G \subset M_{n}(\mathbb{C})$ is a matrix Lie group, then the tangent space to $G$ at the identity is the Lie algebra $\mathfrak{g}$ of $G$ (Corollary 3.46). It is then easy to see that the tangent space $T_{A} G$ at any point $A \in G$ is the space of vectors of the form $X A$ with $X \in \mathfrak{g}$. If the dimension of $\mathfrak{g}$ as a real vector space is $k$, choose a nonzero $k$-linear, alternating form $\alpha: \mathfrak{g}^{k} \rightarrow \mathbb{R}$. (Such form exists and is unique up to multiplication by a constant.) Then we may define a $k$-linear, alternating form $\alpha_{A}:\left(T_{A} G\right)^{k} \rightarrow \mathbb{R}$ by setting

$$
\alpha_{A}\left(Y_{1}, \ldots, Y_{k}\right)=\alpha_{I}\left(Y_{1} A^{-1}, \ldots, Y_{k} A^{-1}\right)
$$

for all $Y_{1}, \ldots, Y_{k} \in T_{A} G$. That is to say, we define $\alpha$ in an arbitrary nonzero fashion at the identity, and we then use the right action of $G$ to "transport" $\alpha$ to every other point in $G$. The resulting family of functionals is a $k$-form on $G$.

Once such an $\alpha$ has been constructed, we can use it to construct an orientation on $G$, by decreeing that an ordered basis $Y_{1}, \ldots, Y_{k}$ for $T_{A} G$ is positively oriented if $\alpha_{A}\left(Y_{1}, \ldots, Y_{k}\right)>0$. If $f: G \rightarrow \mathbb{R}$ is a smooth function, we can integrate the $k$-form $f \alpha$ over nice domains in $G$. If $G$ is compact, we may $f \alpha$ integrate over all of $G$, leading to a notion of integration, which we denote as

$$
\int_{G} f(A) \alpha(A)
$$

Since the orientation on $G$ was defined in terms of the $k$-form $\alpha$ itself, it is not hard to see that if $f(A)>0$ for all $A$, then $\int_{G} f \alpha>0$. Furthermore, since the form $\alpha$ was constructed using the right action of $G$, it is easily seen to be invariant under that action. As a result, the notion of integration of a function over a compact group $G$ is invariant under the right action of $A$ : For all $B \in G$, we have

$$
\int_{G} f(A B) \alpha(A)=\int_{G} f(A) \alpha(A) .
$$

Proof of Theorem 4.28. Choose an arbitrary inner product $\langle\cdot, \cdot\rangle$ on $V$, and then define a map $\langle\cdot, \cdot\rangle_{G}: V \times V \rightarrow \mathbb{C}$ by the formula

$$
\langle v, w\rangle_{G}=\int_{G}\langle\Pi(A) v, \Pi(A) w\rangle \alpha(A)
$$

It is easy to check that $\langle\cdot, \cdot\rangle_{G}$ is an inner product; in particular, the positivity of $\langle\cdot, \cdot\rangle_{G}$ holds because $\langle\Pi(A) v, \Pi(A) v\rangle>0$ for all $A$ if $v \neq 0$. We now compute that for each $B \in G$, we have

$$
\begin{aligned}
\langle\Pi(B) v, \Pi(B) w\rangle_{G} & =\int_{G}\langle\Pi(A) \Pi(B) v, \Pi(A) \Pi(B) w\rangle \alpha(A) \\
& =\int_{G}\langle\Pi(A B) v, \Pi(A B) w\rangle \alpha(A) \\
& =\int_{G}\langle\Pi(A) v, \Pi(A) w\rangle \alpha(A) \\
& =\langle v, w\rangle_{G}
\end{aligned}
$$

where we have used the right invariance of the integral in the third equality. This computation shows that for each $B \in G$, the operator $\Pi(B)$ is unitary with respect to $\langle\cdot, \cdot\rangle_{G}$. Thus, by Proposition 4.27, $\Pi$ is completely reducible.

Note that compactness of the group $G$ is needed to ensure that the integral defining $\langle\cdot, \cdot\rangle_{G}$ is convergent.

\subsection*{Schur's Lemma}
It is desirable to be able to state Schur's lemma simultaneously for groups and Lie algebras. In order to do so, we need to indulge in a common abuse of notation. If, say, $\Pi$ is a representation of $G$ acting on a space $V$, we will refer to $V$ as the representation, without explicit reference to $\Pi$.

\section*{Theorem 4.29 (Schur's Lemma).}
\begin{enumerate}
  \item Let $V$ and $W$ be irreducible real or complex representations of a group or Lie algebra and let $\phi: V \rightarrow W$ be an intertwining map. Then either $\phi=0$ or $\phi$ is an isomorphism.
  \item Let $V$ be an irreducible complex representation of a group or Lie algebra and let $\phi: V \rightarrow V$ be an intertwining map of $V$ with itself. Then $\phi=\lambda I$, for some $\lambda \in \mathbb{C}$.
  \item Let $V$ and $W$ be irreducible complex representations of a group or Lie algebra and let $\phi_{1}, \phi_{2}: V \rightarrow W$ be nonzero intertwining maps. Then $\phi_{1}=\lambda \phi_{2}$, for some $\lambda \in \mathbb{C}$.
\end{enumerate}

It is important to note that the last two points in the theorem hold only over $\mathbb{C}$ (or some other algebraically closed field) and not over $\mathbb{R}$. See Exercise 8.

Before proving Schur's lemma, we obtain two corollaries of it.\\
Corollary 4.30. Let $\Pi$ be an irreducible complex representation of a matrix Lie group $G$. If $A$ is in the center of $G$, then $\Pi(A)=\lambda I$, for some $\lambda \in \mathbb{C}$. Similarly, if $\pi$ is an irreducible complex representation of a Lie algebra $\mathfrak{g}$ and if $X$ is in the center of $\mathfrak{g}$, then $\pi(X)=\lambda I$.

Proof. We prove the group case; the proof of the Lie algebra case is similar. If $A$ is in the center of $G$, then for all $B \in G$,

$$
\Pi(A) \Pi(B)=\Pi(A B)=\Pi(B A)=\Pi(B) \Pi(A) .
$$

However, this says exactly that $\Pi(A)$ is an intertwining map of the space with itself. Thus, by Point 2 of Schur's lemma, $\Pi(A)$ is a multiple of the identity.

Corollary 4.31. An irreducible complex representation of a commutative group or Lie algebra is one dimensional.

Proof. Again, we prove only the group case. If $G$ is commutative, the center of $G$ is all of $G$, so by the previous corollary $\Pi(A)$ is a multiple of the identity for each $A \in G$. However, this means that every subspace of $V$ is invariant! Thus, the only way that $V$ can fail to have a nontrivial invariant subspace is if it is one dimensional.

We now provide the proof of Schur's lemma.\\
Proof of Theorem 4.29. As usual, we will prove just the group case; the proof of the Lie algebra case requires only the obvious notational changes. For Point 1, if $v \in \operatorname{ker}(\phi)$, then

$$
\phi(\Pi(A) v)=\Sigma(A) \phi(v)=0 .
$$

This shows that $\operatorname{ker} \phi$ is an invariant subspace of $V$. Since $V$ is irreducible, we must have $\operatorname{ker} \phi=0$ or $\operatorname{ker} \phi=V$. Thus, $\phi$ is either one-to-one or zero.

Suppose $\phi$ is one-to-one. Then the image of $\phi$ is a nonzero subspace of $W$. On the other hand, the image of $\phi$ is invariant, for if $w \in W$ is of the form $\phi(v)$ for some $v \in V$, then

$$
\Sigma(A) w=\Sigma(A) \phi(v)=\phi(\Pi(A) v)
$$

Since $W$ is irreducible and image $(V)$ is nonzero and invariant, we must have $\operatorname{image}(V)=W$. Thus, $\phi$ is either zero or one-to-one and onto.

For Point 2, suppose $V$ is an irreducible complex representation and that $\phi$ : $V \rightarrow V$ is an intertwining map of $V$ to itself, that is that $\phi \Pi(A)=\Pi(A) \phi$ for all $A \in G$. Since we are working over an algebraically closed field, $\phi$ must have at least one eigenvalue $\lambda \in \mathbb{C}$. If $U$ denotes the corresponding eigenspace for $\phi$, then Proposition A. 2 tells us that each $\Pi(A)$ maps $U$ to itself, meaning that $U$ is an invariant subspace. Since $\lambda$ is an eigenvalue, $U \neq 0$, and so we must have $U=V$, which means that $\phi=\lambda I$ on all of $V$.

For Point 3, if $\phi_{2} \neq 0$, then by Point $1, \phi_{2}$ is an isomorphism. Then $\phi_{1} \circ \phi_{2}^{-1}$ is an intertwining map of $W$ with itself. Thus, by Point $2, \phi_{1} \circ \phi_{2}^{-1}=\lambda I$, whence $\phi_{1}=\lambda \phi_{2}$.

\subsection*{Representations of $\mathbf{s I}(\mathbf{2} ; \mathbb{C})$}
In this section, we will compute (up to isomorphism) all of the finite-dimensional irreducible complex representations of the Lie algebra $\mathrm{SI}(2 ; \mathbb{C})$. This computation is important for several reasons. First, $\mathrm{sI}(2 ; \mathbb{C})$ is the complexification of $\mathrm{su}(2)$, which in turn is isomorphic to $s o(3)$, and the representations of $s o(3)$ are of physical significance. Indeed, the computation we will perform in the proof of Theorem 4.32 is found in every standard textbook on quantum mechanics, under the heading "angular momentum." Second, the representation theory of su(2) is an illuminating example of how one uses commutation relations to determine the representations of a Lie algebra. Third, in determining the representations of a semisimple Lie algebra $\mathfrak{g}$ (Chapters 6 and 7), we will make frequent use of the representation theory of $\mathrm{sl}(2 ; \mathbb{C})$, applying it to various subalgebras of $\mathfrak{g}$ that are isomorphic to $\mathrm{sI}(2 ; \mathbb{C})$.

We use the following basis for $\mathrm{sI}(2 ; \mathbb{C})$ :

$$
X=\left(\begin{array}{ll}
0 & 1 \\
0 & 0
\end{array}\right) ; \quad Y=\left(\begin{array}{ll}
0 & 0 \\
1 & 0
\end{array}\right) ; \quad H=\left(\begin{array}{rr}
1 & 0 \\
0 & -1
\end{array}\right)
$$

which have the commutation relations

$$
\begin{aligned}
& {[H, X]=2 X,} \\
& {[H, Y]=-2 Y,} \\
& {[X, Y]=H .}
\end{aligned}
$$

If $V$ is a finite-dimensional complex vector space and $A, B$, and $C$ are operators on $V$ satisfying $[A, B]=2 B,[A, C]=-2 C$, and $[B, C]=A$, then because of the skew symmetry and bilinearity of brackets, the unique linear map $\pi: \operatorname{sl}(2 ; \mathbb{C}) \rightarrow \mathrm{gl}(V)$ satisfying

$$
\pi(H)=A, \quad \pi(X)=B, \quad \pi(Y)=C
$$

will be a representation of $\operatorname{sI}(2 ; \mathbb{C})$.\\
Theorem 4.32. For each integer $m \geq 0$, there is an irreducible complex representation of $\mathrm{sl}(2 ; \mathbb{C})$ with dimension $m+1$. Any two irreducible complex representations of $\operatorname{sl}(2 ; \mathbb{C})$ with the same dimension are isomorphic. If $\pi$ is an irreducible complex representation of $\operatorname{sl}(2 ; \mathbb{C})$ with dimension $m+1$, then $\pi$ is isomorphic to the representation $\pi_{m}$ described in Sect. 4.2.

Our goal is to show that any finite-dimensional irreducible representation of sl(2; $\mathbb{C}$ ) "looks like" one of the representations $\pi_{m}$ coming from Example 4.10. In that example, the space $V_{m}$ is spanned by eigenvectors for $\pi_{m}(H)$ and the operators $\pi_{m}(X)$ and $\pi_{m}(Y)$ act by shifting the eigenvalues up or down in increments of 2 . We now introduce a simple but crucial lemma that allows us develop a similar structure in an arbitrary irreducible representation of $\mathrm{sI}(2 ; \mathbb{C})$.

Lemma 4.33. Let $u$ be an eigenvector of $\pi(H)$ with eigenvalue $\alpha \in \mathbb{C}$. Then we have

$$
\pi(H) \pi(X) u=(\alpha+2) \pi(X) u .
$$

Thus, either $\pi(X) u=0$ or $\pi(X) u$ is an eigenvector for $\pi(H)$ with eigenvalue $\alpha+2$. Similarly,

$$
\pi(H) \pi(Y) u=(\alpha-2) \pi(Y) u,
$$

so that either $\pi(Y) u=0$ or $\pi(Y) u$ is an eigenvector for $\pi(H)$ with eigenvalue $\alpha-2$.

Proof. We know that $[\pi(H), \pi(X)]=\pi([H, X])=2 \pi(X)$. Thus,

$$
\begin{aligned}
\pi(H) \pi(X) u & =\pi(X) \pi(H) u+2 \pi(X) u \\
& =\pi(X)(\alpha u)+2 \pi(X) u \\
& =(\alpha+2) \pi(X) u
\end{aligned}
$$

The argument with $\pi(X)$ replaced by $\pi(Y)$ is similar.\\
Proof of Theorem 4.32. Let $\pi$ be an irreducible representation of $\operatorname{sl}(2 ; \mathbb{C})$ acting on a finite-dimensional complex vector space $V$. Our strategy is to diagonalize the operator $\pi(H)$. Since we are working over $\mathbb{C}$, the operator $\pi(H)$ must have at least one eigenvector. Let $u$ be an eigenvector for $\pi(H)$ with eigenvalue $\alpha$. Applying Lemma 4.33 repeatedly, we see that

$$
\pi(H) \pi(X)^{k} u=(\alpha+2 k) \pi(X)^{k} u .
$$

Since operator on a finite-dimensional space can have only finitely many eigenvalues, the $\pi(X)^{k} u$ 's cannot all be nonzero. Thus, there is some $N \geq 0$ such that

$$
\pi(X)^{N} u \neq 0
$$

but

$$
\pi(X)^{N+1} u=0
$$

If we set $u_{0}=\pi(X)^{N} u$ and $\lambda=\alpha+2 N$, then,

$$
\begin{aligned}
\pi(H) u_{0} & =\lambda u_{0}, \\
\pi(X) u_{0} & =0 .
\end{aligned}
$$

Let us then define

$$
u_{k}=\pi(Y)^{k} u_{0}
$$

for $k \geq 0$. By Lemma 4.33, we have

$$
\pi(H) u_{k}=(\lambda-2 k) u_{k} .
$$

Now, it is easily verified by induction (Exercise 3) that

$$
\pi(X) u_{k}=k[\lambda-(k-1)] u_{k-1} \quad(k \geq 1) .
$$

Furthermore, since $\pi(H)$ can have only finitely many eigenvalues, the $u_{k}$ 's cannot all be nonzero. There must, therefore, be a non-negative integer $m$ such that

$$
u_{k}=\pi(Y)^{k} u_{0} \neq 0
$$

for all $k \leq m$, but

$$
u_{m+1}=\pi(Y)^{m+1} u_{0}=0 .
$$

If $u_{m+1}=0$, then $\pi(X) u_{m+1}=0$ and so, by (4.15),

$$
0=\pi(X) u_{m+1}=(m+1)(\lambda-m) u_{m} .
$$

Since $u_{m}$ and $m+1$ are nonzero, we must have $\lambda-m=0$. Thus, $\lambda$ must coincide with the non-negative integer $m$.

Thus, for every irreducible representation ( $\pi, V$ ), there exists an integer $m \geq 0$ and nonzero vectors $u_{0}, \ldots, u_{m}$ such that

$$
\begin{aligned}
\pi(H) u_{k} & =(m-2 k) u_{k} \\
\pi(Y) u_{k} & =\left\{\begin{array}{cll}
u_{k+1} & \text { if } k<m \\
0 & \text { if } k=m
\end{array}\right. \\
\pi(X) u_{k} & =\left\{\begin{array}{cc}
k(m-(k-1)) u_{k-1} & \text { if } k>0 \\
0 & \text { if } k=0
\end{array} .\right.
\end{aligned}
$$

The vectors $u_{0}, \ldots, u_{m}$ must be linearly independent, since they are eigenvectors of $\pi(H)$ with distinct eigenvalues (Proposition A.1). Moreover, the ( $m+1$ )dimensional span of $u_{0}, \ldots, u_{m}$ is explicitly invariant under $\pi(H), \pi(X)$, and $\pi(Y)$ and, hence, under $\pi(Z)$ for all $Z \in \operatorname{sl}(2 ; \mathbb{C})$. Since $\pi$ is irreducible, this space must be all of $V$.

We have shown that every irreducible representation of $\operatorname{sI}(2 ; \mathbb{C})$ is of the form (4.16). Conversely, if we define $\pi(H), \pi(X)$, and $\pi(Y)$ by (4.16) (where the $u_{k}$ 's are basis elements for some ( $m+1$ )-dimensional vector space), it is\\
not hard to check that operators defined as in (4.16) really do satisfy the $\operatorname{sl}(2 ; \mathbb{C})$ commutation relations (Exercise 4). Furthermore, the we may prove irreducibility of this representation in the same way as in the proof of Proposition 4.11.

The preceding analysis shows that every irreducible representation of dimension $m+1$ must have the form in (4.16), which shows that any two such representations are isomorphic. In particular, the ( $m+1$ )-dimensional representation $\pi_{m}$ described in Sect. 4.2 must be isomorphic to (4.16).

This completes the proof of Theorem 4.32.\\
As mentioned earlier in this section, the representation theory of $\operatorname{sl}(2 ; \mathbb{C})$ plays a key role in the representation theory of other Lie algebras, such as $\mathrm{SI}(3 ; \mathbb{C})$, because these Lie algebras contain subalgebras isomorphic to $\mathrm{SI}(2 ; \mathbb{C})$. For such applications, we need a few results about finite-dimensional representations of $\operatorname{sI}(2 ; \mathbb{C})$ that are not necessarily irreducible.

Theorem 4.34. If ( $\pi, V$ ) is a finite-dimensional representation of $\mathrm{sl}(2 ; \mathbb{C})$, not necessarily irreducible, the following results hold.

\begin{enumerate}
  \item Every eigenvalue of $\pi(H)$ is an integer. Furthermore, if $v$ is an eigenvector for $\pi(H)$ with eigenvalue $\lambda$ and $\pi(X) v=0$, then $\lambda$ is a non-negative integer.
  \item The operators $\pi(X)$ and $\pi(Y)$ are nilpotent.
  \item If we define $S: V \rightarrow V$ by
\end{enumerate}

$$
S=e^{\pi(X)} e^{-\pi(Y)} e^{\pi(X)}
$$

then $S$ satisfies

$$
S \pi(H) S^{-1}=-\pi(H) .
$$

\begin{enumerate}
  \setcounter{enumi}{3}
  \item If an integer $k$ is an eigenvalue for $\pi(H)$, so is each of the numbers
\end{enumerate}

$$
-|k|,-|k|+2, \ldots,|k|-2,|k| .
$$

Since $\mathrm{SU}(2)$ is simply connected, Theorem 5.6 will tell us that the representations of $\mathrm{SI}(2 ; \mathbb{C}) \cong \mathrm{SU}(2)_{\mathbb{C}}$ are in one-to-correspondence with the representations of $\mathrm{SU}(2)$. Since $\mathrm{SU}(2)$ is compact, Theorem 4.28 then tells us that every representation of $\operatorname{sI}(2 ; \mathbb{C})$ is completely reducible. If we decompose $V$ as a direct sum of irreducibles, it is easy to prove the theorem for each summand separately. It is, however, preferable to give a proof of the theorem that does not rely on Theorem 5.6, which in turn relies on the Baker-Campbell-Hausdorff formula.

See also Exercise 13 for a different approach to the first part of Point 1, and Exercise 14 for a different approach to Point 3.

Proof. For Point 1, suppose $v$ is an eigenvalue of $\pi(H)$ with eigenvalue $\lambda$. Then there is some $N \geq 0$ such that $\pi(X)^{N} v \neq 0$ but $\pi(X)^{N+1} v=0$, where $\pi(X)^{N} v$ is an eigenvector of $\pi(H)$ with eigenvalue $\lambda+2 N$. The proof of Theorem 4.32 shows that $m:=\lambda+2 N$ must be a non-negative integer, so that $\lambda$ is an integer. If $\pi(X) v=0$ then we take $N=0$ and $\lambda=m$ is non-negative.

For Point 2, it follows from the $S N$ decomposition (Sect. A.3) that $\pi(H)$ has a basis of generalized eigenvectors, that is, vectors $v$ for which $(\pi(H)-\lambda I)^{k} v=0$ for some positive integer $k$. But, using the commutation relation $[H, X]=2 X$ and induction on $k$, we can see that

$$
[\pi(H)-(\lambda+2) I]^{k} \pi(X)=\pi(X)[\pi(H)-\lambda I]^{k} .
$$

Thus, if $v$ is a generalized eigenvector for $\pi(H)$ with eigenvalue $\lambda$, then $\pi(X) v$ is either zero or a generalized eigenvector with eigenvalue $\lambda+2$. Applying $\pi(X)$ repeatedly to a generalized eigenvector for $\pi(H)$ must eventually give zero, since $\pi(H)$ can have only finitely many generalized eigenvalues. Thus, $\pi(X)$ is nilpotent. A similar argument applies to $\pi(Y)$.

For Point 3, we note that

$$
S \pi(H) S^{-1}=e^{\pi(X)} e^{-\pi(Y)} e^{\pi(X)} \pi(H) e^{-\pi(X)} e^{\pi(Y)} e^{-\pi(X)} .
$$

By Proposition 3.35, we have

$$
e^{\pi(X)} \pi(H) e^{-\pi(X)}=\operatorname{Ad}_{e^{\pi(X)}}(\pi(H))=e^{\operatorname{ad}_{\pi(X)}}(\pi(H))
$$

and similarly for the remaining products in (4.17).\\
Now, $\operatorname{ad}_{X}(X)=0, \operatorname{ad}_{X}(H)=-2 X$ and $\operatorname{ad}_{X}(Y)=H$, and similarly with $\pi$ applied to each Lie algebra element. Thus,

$$
e^{\mathrm{ad}_{\pi(X)}}(\pi(H))=\pi(H)-2 \pi(X) .
$$

Meanwhile, $\operatorname{ad}_{Y}(X)=-H, \operatorname{ad}_{Y}(H)=2 Y$, and $\operatorname{ad}_{Y}(Y)=0$. Thus,

$$
\begin{aligned}
& e^{-\mathrm{ad}_{\pi(Y)}}(\pi(H)-2 \pi(X)) \\
& =\pi(H)-2 \pi(X)-2 \pi(Y)-2 \pi(H)+\frac{1}{2} 4 \pi(Y) \\
& =-\pi(H)-2 \pi(X)
\end{aligned}
$$

Finally,

$$
\begin{aligned}
e^{\operatorname{ad}_{\pi(X)}(-\pi(H)-2 \pi(X))} & =-\pi(H)-2 \pi(X)+2 \pi(X) \\
& =-\pi(H),
\end{aligned}
$$

which establishes Point 3.\\
For Point 4, assume first that $k$ is non-negative and let $v$ be an eigenvector for $\pi(H)$ with eigenvalue $k$. Then as in Point 1 , there is then another eigenvector $v_{0}$ for $\pi(H)$ with eigenvalue $m:=k+2 N \geq k$ and such that $\pi(X) v_{0}=0$. Then by the proof of Theorem 4.32, we obtain a chain of eigenvectors $v_{0}, v_{1}, \ldots, v_{m}$ for\\
$\pi(H)$ with eigenvalues ranging from $m$ to $-m$ in increments of 2 . These eigenvalues include all of the numbers $k, k-2, \ldots,-k$. If $k$ is negative and $v$ is an eigenvector for $\pi(H)$ with eigenvalue $k$, then $S v$ is an eigenvector for $\pi(H)$ with eigenvalue $|k|$. Hence, by the preceding argument, each of the numbers from $|k|$ to $-|k|$ in increments of 2 is an eigenvalue.

\subsection*{Group Versus Lie Algebra Representations}
We know from Chapter 3 (Theorem 3.28) that every Lie group homomorphism gives rise to a Lie algebra homomorphism. In particular, every representation of a matrix Lie group gives rise to a representation of the associated Lie algebra. In Chapter 5, we will prove a partial converse to this result: If $G$ is a simply connected matrix Lie group with Lie algebra $\mathfrak{g}$, every representation of $\mathfrak{g}$ comes from a representation of $G$. (See Theorem 5.6). Thus, for a simply connected matrix Lie group $G$, there is a natural one-to-one correspondence between the representations of $G$ and the representations of $\mathfrak{g}$.

It is instructive to see how this general theory works out in the case of $\mathrm{SU}(2)$ (which is simply connected) and $\mathrm{SO}(3)$ (which is not). For every irreducible representation $\pi$ of $\operatorname{su}(2)$, the complex-linear extension of $\pi$ to $\operatorname{sl}(2 ; \mathbb{C})$ must be isomorphic (Theorem 4.32) to one of the representations $\pi_{m}$ described in Sect. 4.2. Since those representations are constructed from representations of the group $\mathrm{SU}(2)$, we can see directly (without appealing to Theorem 5.6) that every irreducible representation of $\operatorname{su}(2)$ comes from a representation of $\operatorname{SU}(2)$.

Now, by Example 3.27, there is a Lie algebra isomorphism $\phi: \operatorname{su}(2) \rightarrow \operatorname{so}(3)$ such that $\phi\left(E_{j}\right)=F_{j}, j=1,2,3$, where $\left\{E_{1}, E_{2}, E_{3}\right\}$ and $\left\{F_{1}, F_{2}, F_{3}\right\}$ are the bases listed in the example. Thus, the irreducible representations of so(3) are precisely of the form $\sigma_{m}=\pi_{m} \circ \phi^{-1}$. We wish to determine, for a particular $m$, whether or not there is a representation $\Sigma_{m}$ of the group SO (3) such that $\Sigma_{m}(\exp X)=\exp \left(\sigma_{m}(X)\right)$ for all $X$ in so(3).

Proposition 4.35. Let $\sigma_{m}=\pi_{m} \circ \phi^{-1}$ be an irreducible complex representations of the Lie algebra $\mathrm{SO}(3)(m \geq 0)$. If $m$ is even, there is a representation $\Sigma_{m}$ of the group SO (3) such that $\Sigma_{m}\left(e^{X}\right)=e^{\sigma_{m}(X)}$ for all $X$ in so (3). If $m$ is odd, there is no such representation of $\mathrm{SO}(3)$.

Note that the condition that $m$ be even is equivalent to the condition that $\operatorname{dim} V_{m}=m+1$ be odd. Thus, it is the odd-dimensional representations of the Lie algebra $\mathbf{S O}(\mathbf{3})$ which come from group representations. In the physics literature, the representations of $\operatorname{su}(2) \cong \operatorname{so}(3)$ are labeled by the parameter $l=m / 2$. In terms of this notation, a representation of so(3) comes from a representation of SO(3) if and only if $l$ is an integer. The representations with $l$ an integer are called "integer spin"; the others are called "half-integer spin."

For any $m$, one could attempt to construct $\Sigma_{m}$ by the construction in the proof of Theorem 5.6. The construction is based on defining $\Sigma_{m}(A)$ along a path joining\\
$I$ to $A$ and then proving that the value of $\Sigma_{m}(A)$ is independent of the choice of path. The construction of $\Sigma_{m}$ along a path goes through without change. Since SO (3) is not simply connected, however, two paths in $\mathrm{SO}(3)$ with the same endpoint are not necessarily homotopic with endpoints fixed and the proof of independence of the path breaks down. One can join the identity to itself, for example, either by the constant path or by the path consisting of rotations by angle $2 \pi t$ in the $(y, z)$ plane, $0 \leq t \leq 1$. If one defines $\Sigma_{m}$ along the constant path, one gets the value $\Sigma_{m}(I)=I$. If $m$ is odd, however, and one defines $\Sigma_{m}$ along the path of rotations in the ( $y, z$ )-plane, then one gets the value $\Sigma_{m}(I)=-I$, as the calculations in the proof of Proposition 4.35 will show. This strongly suggests (and Proposition 4.35 confirms) that when $m$ is odd, there is no way to define $\Sigma_{m}$ as a "single-valued" representation of $\mathrm{SO}(3)$.

An electron, for example, is a "spin one-half" particle, which means that it is described in quantum mechanics in a way that involves the two-dimensional representation $\sigma_{1}$ of $\mathbf{s o}(3)$. In the physics literature, one finds statements to the effect that applying a $360^{\circ}$ rotation to the wave function of the electron gives back the negative of the original wave function. This statement reflects that if one attempts to construct the nonexistent representation $\Sigma_{1}$ of $\mathrm{SO}(3)$, then when defining $\Sigma_{1}$ along a path of rotations in the $(x, y)$-plane, one gets that $\Sigma_{1}(I)=-I$.

Proof. Suppose, first, that $m$ is odd and suppose that there a $\Sigma_{m}$ existed. Computing as in Sect. 2.2, we see that

$$
e^{2 \pi F_{1}}=\left(\begin{array}{rrr}
1 & 0 & 0 \\
0 & \cos 2 \pi & -\sin 2 \pi \\
0 & \sin 2 \pi & \cos 2 \pi
\end{array}\right)=I .
$$

Meanwhile, $\sigma_{m}\left(F_{1}\right)=\pi_{m}\left(\phi^{-1}\left(F_{1}\right)\right)=\pi_{m}\left(E_{1}\right)$, with $E_{1}=i H / 2$, where, as usual, $H$ is the diagonal matrix with diagonal entries $(1,-1)$. We know that there is a basis $u_{0}, u_{1}, \ldots, u_{m}$ for $V_{m}$ such that $u_{k}$ is an eigenvector for $\pi_{m}(H)$ with eigenvalue $m-2 j$. This means that $u_{j}$ is also an eigenvector for $\sigma_{m}\left(F_{1}\right)=i \pi_{m}(H) / 2$, with eigenvalue $i(m-2 j) / 2$. Thus, in the basis $\left\{u_{j}\right\}$, we have

$$
\sigma_{m}\left(F_{1}\right)=\left(\begin{array}{llll}
\frac{i}{2} m & & & \\
& \frac{i}{2}(m-2) & & \\
& & \ddots & \\
& & & \frac{i}{2}(-m)
\end{array}\right)
$$

Since we are assuming $m$ is odd, $m-2 j$ is an odd integer. Thus, $e^{2 \pi \sigma_{m}\left(F_{1}\right)}$ has eigenvalues $e^{2 \pi i(m-2 j) / 2}=-1$ in the basis $\left\{u_{j}\right\}$, showing that $e^{2 \pi \sigma_{m}\left(F_{1}\right)}=-I$. Thus, on the one hand,

$$
\Sigma_{m}\left(e^{2 \pi F_{1}}\right)=\Sigma_{m}(I)=I
$$

whereas, on the other hand,

$$
\Sigma_{m}\left(e^{2 \pi F_{1}}\right)=e^{2 \pi \sigma_{m}\left(F_{1}\right)}=-I
$$

Thus, there can be no such group representation $\Sigma_{m}$.\\
Suppose now that $m$ is even. Recall that the Lie algebra isomorphism $\phi$ comes from the surjective group homomorphism $\Phi$ in Proposition 1.19, where $\operatorname{ker}(\Phi)= \{I,-I\}$. Let $\Pi_{m}$ be the representation of $\mathrm{SU}(2)$ in Example 4.10. Now, $e^{2 \pi E_{1}}= -I$, and, thus,

$$
\Pi_{m}(-I)=\Pi_{m}\left(e^{2 \pi E_{1}}\right)=e^{\pi_{m}\left(2 \pi E_{1}\right)}
$$

If, however $m$ is even, then $e^{\pi_{m}\left(2 \pi E_{1}\right)}$ is diagonal in the basis $\left\{u_{j}\right\}$ with eigenvalues $e^{2 \pi i(m-2 j) / 2}=1$, showing that $\Pi(-I)=e^{\pi_{m}\left(2 \pi E_{1}\right)}=I$.

Now, for each $R \in \mathrm{SO}(3)$, there is a unique pair of elements $\{U,-U\}$ such that $\Phi(U)=\Phi(-U)=R$. Since $\Pi_{m}(-I)=I$, we see that $\Pi_{m}(U)=\Pi_{m}(-U)$. It thus makes sense to define

$$
\Sigma_{m}(R)=\Pi_{m}(U) .
$$

It is easy to see that $\Sigma_{m}$ is a Lie group homomorphism, and, by construction, we have $\Pi_{m}=\Sigma_{m} \circ \Phi$. Thus, the Lie algebra representation $\sigma_{m}$ associated to $\Sigma_{m}$ satisfies $\pi_{m}=\sigma_{m} \circ \phi$ or $\sigma_{m}=\pi \circ \phi^{-1}$, showing that $\Sigma_{m}$ is the desired representation of SO(3).

\subsection*{A Nonmatrix Lie Group}
In this section, we will show that the Lie group introduced in Sect. 1.5 is not isomorphic to a matrix Lie group. (The universal cover of $\operatorname{SL}(2 ; \mathbb{R})$ is another example of a Lie group that is not a matrix Lie group; see Sect. 5.8.) The group in question is $G=\mathbb{R} \times \mathbb{R} \times S^{1}$, with the group product defined by

$$
\left(x_{1}, y_{1}, u_{1}\right) \cdot\left(x_{2}, y_{2}, u_{1}\right)=\left(x_{1}+x_{2}, y_{1}+y_{2}, e^{i x_{1} y_{2}} u_{1} u_{2}\right) .
$$

Meanwhile, let $H$ be the Heisenberg group and consider the map $\Phi: H \rightarrow G$ given by

$$
\Phi\left(\begin{array}{lll}
1 & a & b \\
0 & 1 & c \\
0 & 0 & 1
\end{array}\right)=\left(a, c, e^{2 \pi i b}\right)
$$

Direct computation shows that $\Phi$ is a homomorphism. The kernel of $\Phi$ is the discrete normal subgroup $N$ of $H$ given by

$$
N=\left\{\left.\left(\begin{array}{lll}
1 & 0 & n \\
0 & 1 & 0 \\
0 & 0 & 1
\end{array}\right) \right\rvert\, n \in \mathbb{Z}\right\}
$$

Now, suppose that $\Sigma$ is any finite-dimensional representation of $G$. Then we can define an associated representation $\Pi$ of $H$ by $\Pi=\Sigma \circ \Phi$. Clearly, the kernel of any such representation of $H$ must include the kernel $N$ of $\Phi$. Now, let $Z(H)$ denote the center of $H$, which is easily shown to be

$$
Z(H)=\left\{\left.\left(\begin{array}{lll}
1 & 0 & b \\
0 & 1 & 0 \\
0 & 0 & 1
\end{array}\right) \right\rvert\, b \in \mathbb{R}\right\}
$$

Theorem 4.36. Let $\Pi$ be any finite-dimensional representation of $H$. If $\operatorname{ker} \Pi \supset N$, then $\operatorname{ker} \Pi \supset Z(H)$.

Once this is established, we will be able to conclude that there are no faithful finite-dimensional representations of $G$. After all, if $\Sigma$ is any finite-dimensional representation of $G$, then the kernel of $\Pi=\Sigma \circ \Phi$ will contain $N$ and, thus, $Z(H)$, by the theorem. Thus, for all $b \in \mathbb{R}$,

$$
\Pi\left(\begin{array}{lll}
1 & 0 & b \\
0 & 1 & 0 \\
0 & 0 & 1
\end{array}\right)=\Sigma\left(0,0, e^{2 \pi i b}\right)=I
$$

This means that the kernel of $\Sigma$ contains all elements of the form $(0,0, u)$ and $\Sigma$ is not faithful. Thus, we obtain the following result.

Corollary 4.37. The Lie group $G$ has no faithful finite-dimensional representations. In particular, $G$ is not isomorphic to any matrix Lie group.

We now begin the proof of Theorem 4.36.\\
Lemma 4.38. If $X$ is a nilpotent matrix and $e^{t X}=I$ for some nonzero $t$, then $X=0$.

Proof. Since $X$ is nilpotent, the power series for $e^{t X}$ terminates after a finite number of terms. Thus, each entry of $e^{t X}$ depends polynomially on $t$; that is, there exist polynomials $p_{j k}(t)$ such that $\left(e^{t X}\right)_{j k}=p_{j k}(t)$. If $e^{t X}=I$ for some nonzero $t$, then $e^{n t X}=I$ for all $n \in \mathbb{Z}$, showing that $p_{j k}(n t)=\delta_{j k}$ for all $n$. However, a polynomial that takes on a certain value infinitely many times must be constant. Thus, actually, $e^{t X}=I$ for all $t$, from which it follows that $X=0$.

Proof of Theorem 4.36. Let $\pi$ be the associated representation of the Lie algebra $\mathfrak{h}$ of $H$. Let $\{X, Y, Z\}$ be the following basis for $\mathfrak{h}$ :

$$
X=\left(\begin{array}{lll}
0 & 1 & 0 \\
0 & 0 & 0 \\
0 & 0 & 0
\end{array}\right), Y=\left(\begin{array}{lll}
0 & 0 & 0 \\
0 & 0 & 1 \\
0 & 0 & 0
\end{array}\right), Z=\left(\begin{array}{lll}
0 & 0 & 1 \\
0 & 0 & 0 \\
0 & 0 & 0
\end{array}\right) .
$$

These satisfy the commutation relations $[X, Y]=Z$ and $[X, Z]=[Y, Z]=0$.\\
We now claim that $\pi(Z)$ must be nilpotent, or, equivalently, that all of the eigenvalues of $\pi(Z)$ are zero. Let $\lambda$ be an eigenvalue for $\pi(Z)$ and let $V_{\lambda}$ be the associated eigenspace. Then $V_{\lambda}$ is certainly invariant under $\pi(Z)$. Furthermore, since $\pi(X)$ and $\pi(Y)$ commute with $\pi(Z)$, Proposition A. 2 tells us that $V_{\lambda}$ is invariant under $\pi(X)$ and $\pi(Y)$. Thus, the restriction of $\pi(Z)$ to $V_{\lambda}$-namely, $\lambda I$ is the commutator of the restrictions to $V_{\lambda}$ of $\pi(X)$ and $\pi(Y)$. Since the trace of a commutator is zero, we have $0=\lambda \operatorname{dim}\left(V_{\lambda}\right)$ and $\lambda$ must be zero.

Now, direct calculation shows that $e^{n Z}$ belongs to $N$ for all integers $n$. Thus, if $\Pi$ is a representation of $H$ and $\operatorname{ker} \Pi \supset N$, we have $\Pi\left(e^{n Z}\right)=I$ for all $n$. Since $\pi(Z)$ is nilpotent, Lemma 4.38 tells us that $\pi(Z)$ is zero and thus that $\Pi\left(e^{t Z}\right)= e^{t \pi(Z)}=I$ for all $t \in \mathbb{R}$. Since every element of $Z(H)$ is of the form $e^{t Z}$ for some $t$, we have the desired conclusion.

\subsection*{Exercises}
\begin{enumerate}
  \item Prove Point 2 of Proposition 4.5.
  \item Show that the adjoint representation and the standard representation are isomorphic representations of the Lie algebra so(3). Show that the adjoint and standard representations of the group $\mathrm{SO}(3)$ are isomorphic.
  \item Using the commutation relation $[\pi(X), \pi(Y)]=\pi(H)$ and induction on $k$, verify the relation (4.15).
  \item Define a vector space with basis $u_{0}, u_{1}, \ldots, u_{m}$. Now, define operators $\pi(H)$, $\pi(X)$, and $\pi(Y)$ by formula (4.16). Verify by direct computation that the operators defined by (4.16) satisfy the commutation relations (4.11) for $\mathrm{SI}(2 ; \mathbb{C})$.\\
Hint: When dealing with $\pi(Y)$, treat the case of $u_{k}, k<m$, separately from the case of $u_{m}$, and similarly for $\pi(X)$.
  \item Consider the standard representation of the Heisenberg group, acting on $\mathbb{C}^{3}$. Determine all subspaces of $\mathbb{C}^{3}$ which are invariant under the action of the Heisenberg group. Is this representation completely reducible?
  \item Prove Proposition 4.22.
\end{enumerate}

Hint: There is a one-to-one correspondence between subspaces of $V$ and subspaces of $V^{*}$ as follows: For any subspace $W$ of $V$, the annihilator of $W$ is the subspace of all $\phi$ in $V^{*}$ such that $\phi$ is zero on $W$. See Sect. A.7.\\
7. Prove Theorem 4.14.

Hints: For existence, choose bases $\left\{e_{j}\right\}$ and $\left\{f_{k}\right\}$ for $U$ and $V$. Then define a space $W$ which has as a basis $\left\{w_{j k} \mid 1 \leq j \leq n, 1 \leq k \leq m\right\}$. Define $\phi\left(e_{j}, f_{k}\right)=w_{j k}$ and extend by bilinearity. For uniqueness, use the universal property.\\
8. Let $\mathrm{SO}(2)$ act on $\mathbb{R}^{2}$ in the obvious way. Show that $\mathbb{R}^{2}$ is an irreducible real representation under this action, but that Point 2 of Schur's lemma (Theorem 4.29) fails.\\
9. Suppose $V$ is a finite-dimensional representation of a group or Lie algebra and that $W$ is a nonzero invariant subspace of $V$. Show that there exists a nonzero irreducible invariant subspace for $V$ that is contained in $W$.\\
10. Suppose that $V_{1}$ and $V_{2}$ are nonisomorphic irreducible representations of a group or Lie algebra, and consider the associated representation $V_{1} \oplus V_{2}$. Regard $V_{1}$ and $V_{2}$ as subspaces of $V_{1} \oplus V_{2}$ in the obvious way. Following the outline below, show that $V_{1}$ and $V_{2}$ are the only nontrivial invariant subspaces of $V_{1} \oplus V_{2}$.\\
(a) First assume that $U$ is a nontrivial irreducible invariant subspace. Let $P_{1}$ : $V_{1} \oplus V_{2} \rightarrow V_{1}$ be the projection onto the first factor and let $P_{2}$ be the projection onto the second factor. Show that $P_{1}$ and $P_{2}$ are intertwining maps. Show that $U=V_{1}$ or $U=V_{2}$.\\
(b) Using Exercise 9, show that $V_{1}$ and $V_{2}$ are the only nontrivial invariant subspaces of $V_{1} \oplus V_{2}$.\\
11. Suppose that $V$ is an irreducible finite-dimensional representation of a group or Lie algebra over $\mathbb{C}$, and consider the associated representation $V \oplus V$. Show that every nontrivial invariant subspace $U$ of $V \oplus V$ is isomorphic to $V$ and is of the form

$$
U=\left\{\left(\lambda_{1} v, \lambda_{2} v\right) \mid v \in V\right\}
$$

for some constants $\lambda_{1}$ and $\lambda_{2}$, not both zero.\\
12. Recall the spaces $V_{m}$ introduced in Sect. 4.2, viewed as representations of the Lie algebra $\operatorname{sI}(2 ; \mathbb{C})$. In particular, consider the space $V_{1}$ (which has dimension 2).\\
(a) Regard $V_{1} \otimes V_{1}$ as a representation of $\mathrm{sI}(2 ; \mathbb{C})$, as in Definition 4.20. Show that this representation is not irreducible.\\
(b) Now, view $V_{1} \otimes V_{1}$ as a representation of $\operatorname{sl}(2 ; \mathbb{C}) \oplus \operatorname{sl}(2 ; \mathbb{C})$, as in Definition 4.19. Show that this representation is irreducible.\\
13. Let ( $\Pi, V$ ) be a finite-dimensional representation of $\mathrm{SU}(2)$ with associated representation $\pi$ of $\mathrm{su}(2)$, which extends by complex linearity to $\mathrm{sI}(2 ; \mathbb{C})$. (Since SU(2) is simply connected, Theorem 5.6 will show that every representation $\pi$ of $\operatorname{sl}(2 ; \mathbb{C})$ arise in this way.) If $H$ is the diagonal matrix with diagonal entries $(1,-1)$, show that $e^{2 \pi i H}=I$ and use this to prove (independently of Theorem 4.34) that every eigenvalue of $\pi(H)$ is an integer.\\
14. Let ( $\Pi, V$ ) be a finite-dimensional representation of $\mathrm{SU}(2)$ with associated representation $\pi$ of $\operatorname{su}(2)$, which extends by complex linearity to $\operatorname{sl}(2 ; \mathbb{C})$. If $X, Y$, and $H$ are the usual basis element for $\mathrm{sI}(2 ; \mathbb{C})$, compute $e^{X} e^{-Y} e^{X}$ and show that

$$
e^{X} e^{-Y} e^{X} H\left(e^{X} e^{-Y} e^{X}\right)^{-1}=-H
$$

Use this result to give a different proof of Point 3 of Theorem 4.34.

\section*{Chapter 5 \\
 The Baker-Campbell-Hausdorff Formula and Its Consequences}
\subsection*{The "Hard" Questions}
Consider three elementary results from the preceding chapters of this book: (1) Every matrix Lie group $G$ has a Lie algebra $\mathfrak{g}$. (2) A continuous homomorphism $\Phi$ between matrix Lie groups $G$ and $H$ gives rise to a Lie algebra homomorphism $\phi: \mathfrak{g} \rightarrow \mathfrak{h}$. (3) If $G$ and $H$ are matrix Lie groups and $H$ is a subgroup of $G$, then the Lie algebra $\mathfrak{h}$ of $H$ is a subalgebra of the Lie algebra $\mathfrak{g}$ of $G$. Each of these results goes in the "easy" direction, from a group notion to an associated Lie algebra notion. In this chapter, we attempt to go in the "hard" direction, from the Lie algebra to the Lie group. We will investigate three questions relating to the preceding three theorems.

\begin{itemize}
  \item Question 1: Is every finite-dimensional, real Lie algebra the Lie algebra of some matrix Lie group?
  \item Question 2: Let $G$ and $H$ be matrix Lie groups with Lie algebras $\mathfrak{g}$ and $\mathfrak{h}$, respectively, and let $\phi: \mathfrak{g} \rightarrow \mathfrak{h}$ be a Lie algebra homomorphism. Does there exists a Lie group homomorphism $\Phi: G \rightarrow H$ such that $\Phi\left(e^{X}\right)=e^{\phi(X)}$ for all $X \in \mathfrak{g}$ ?
  \item Question 3: If $G$ is a matrix Lie group with Lie algebra $\mathfrak{g}$ and $\mathfrak{h}$ is a subalgebra of $\mathfrak{g}$, is there a matrix Lie group $H \subset G$ whose Lie algebra is $\mathfrak{h}$ ?
\end{itemize}

The answer to Question 1 is yes; see Sect. 5.10. The answer to Question 2 is, in general, no, but is yes if $G$ is simply connected; see Sect. 5.7. The answer to Question 3 is no, in general, but is yes if we allow $H$ to be a "connected Lie subgroup" that is not necessarily closed; see Sect. 5.9.

Our tool for investigating these questions is the Baker-Campbell-Hausdorff formula, which expresses $\log \left(e^{X} e^{Y}\right)$, where $X$ and $Y$ are sufficiently small $n \times n$ matrices, in Lie-algebraic terms, that is, in terms of iterated commutators involving $X$ and $Y$. The formula implies that all information about the product operation on a matrix Lie group, at least near the identity, is encoded in the Lie algebra. In the\\
case of Questions 2 and 3 in the preceding paragraph, we will give a complete proof of the theorem that answers the question. In the case of Question 1, we will need to assume Ado's theorem, which asserts that every finite-dimensional real Lie algebra is isomorphic to an algebra of matrices.

\subsection*{An Illustrative Example}
In this section, we prove one of the main theorems alluded to above (the answer to Question 2), in the case of the Heisenberg group. We introduce the problem in a general setting and then specialize to the Heisenberg case. Suppose $G$ and $H$ are matrix Lie groups with Lie algebras $\mathfrak{g}$ and $\mathfrak{h}$, respectively, and suppose $\phi$ : $\mathfrak{g} \rightarrow \mathfrak{h}$ is a Lie algebra homomorphism. We would like to construct a Lie group homomorphism $\Phi: G \rightarrow H$ such that $\Phi\left(e^{X}\right)=e^{\phi(X)}$ for all $X$ in $\mathfrak{g}$. In light of Theorem 3.42, we can define a map $\Phi$ from a neighborhood $U$ of the identity in $G$ into $H$ as follows:

$$
\Phi(A)=e^{\phi(\log A)}
$$

so that

$$
\Phi\left(e^{X}\right)=e^{\phi(X)}
$$

at least for sufficiently small $X$.\\
A key issue is then to show that $\Phi$, as defined near the identity by (5.1) or (5.2), is a "local homomorphism." Suppose $A=e^{X}$ and $B=e^{Y}$, where $X, Y \in \mathfrak{g}$ are small that $e^{X}, e^{Y}$, and $e^{X} e^{Y}$ are all in the domain of $\Phi$. To compute $\Phi(A B)=\Phi\left(e^{X} e^{Y}\right)$, we need to express $e^{X} e^{Y}$ in the form $e^{Z}$, so that $\Phi\left(e^{X} e^{Y}\right)$ will equal $e^{\phi(Z)}$. The Baker-Campbell-Hausdorff formula states that for sufficiently small $X$ and $Y$, we have

$$
\begin{aligned}
Z & =\log \left(e^{X} e^{Y}\right) \\
& =X+Y+\frac{1}{2}[X, Y]+\frac{1}{12}[X,[X, Y]]-\frac{1}{12}[Y,[X, Y]]+\cdots,
\end{aligned}
$$

where the " $\ldots$ " refers to additional terms involving iterated brackets of $X$ and $Y$. (A precise statement of and a proof of the formula will be given in subsequent sections.) If $\phi$ is a Lie algebra homomorphism, then

$$
\begin{aligned}
\phi\left(\log \left(e^{X} e^{Y}\right)\right)= & \phi(X)+\phi(Y)+\frac{1}{2}[\phi(X), \phi(Y)] \\
& +\frac{1}{12}[\phi(X),[\phi(X), \phi(Y)]]-\frac{1}{12}[\phi(Y),[\phi(X), \phi(Y)]]+\cdots \\
= & \log \left(e^{\phi(X)} e^{\phi(Y)}\right)
\end{aligned}
$$

It then follows that $\Phi\left(e^{X} e^{Y}\right)=e^{\phi(X)} e^{\phi(Y)}=\Phi\left(e^{X}\right) \Phi\left(e^{Y}\right)$.

In the general case, it requires considerable effort to prove the Baker-CampbellHausdorff formula and then to prove that, when $G$ is simply connected, $\Phi$ can be extended to all of $G$. In the case of the Heisenberg group (which is simply connected), the argument can be greatly simplified.

Theorem 5.1. Suppose $X$ and $Y$ are $n \times n$ complex matrices, and that $X$ and $Y$ commute with their commutator:

$$
[X,[X, Y]]=[Y,[X, Y]]=0 .
$$

Then we have

$$
e^{X} e^{Y}=e^{X+Y+\frac{1}{2}[X, Y]}
$$

This is the special case of (5.3) in which the series terminates after the [ $X, Y$ ] term. See Exercise 5 for another special case of the Baker-Campbell-Hausdorff formula.

Proof. Consider $X$ and $Y$ in $M_{n}(\mathbb{C})$ satisfying (5.5). We will prove that

$$
e^{t X} e^{t Y}=\exp \left(t X+t Y+\frac{t^{2}}{2}[X, Y]\right)
$$

which reduces to the desired result in the case $t=1$. Since, by assumption, $[X, Y]$ commutes with $X$ and $Y$, the above relation is equivalent to

$$
e^{t X} e^{t Y} e^{-\frac{t^{2}}{2}[X, Y]}=e^{t(X+Y)}
$$

Let us denote by $A(t)$ the left-hand side of (5.6) and by $B(t)$ the right-hand side. Our strategy will be to show that $A(t)$ and $B(t)$ satisfy the same differential equation, with the same initial conditions. We can see immediately that

$$
\frac{d B}{d t}=B(t)(X+Y)
$$

On the other hand, differentiating $A(t)$ by means of the product rule gives

$$
\begin{aligned}
\frac{d A}{d t}= & e^{t X} X e^{t Y} e^{-\frac{t^{2}}{2}[X, Y]}+e^{t X} e^{t Y} Y e^{-\frac{t^{2}}{2}[X, Y]} \\
& +e^{t X} e^{t Y} e^{-\frac{t^{2}}{2}[X, Y]}(-t[X, Y])
\end{aligned}
$$

(The correctness of the last term may be verified by differentiating $e^{-t^{2}[X, Y] / 2}$ term by term.) Now, since $Y$ commutes with $[X, Y]$, it also commute with $e^{-t^{2}[X, Y] / 2}$. Thus, the second term on the right in (5.7) can be rewritten as

$$
e^{t X} e^{t Y} e^{-\frac{t^{2}}{2}[X, Y]} Y .
$$

For the first term on the right-hand side of (5.7), we compute, using Proposition 3.35, that

$$
\begin{aligned}
X e^{t Y} & =e^{t Y} e^{-t Y} X e^{t Y} \\
& =e^{t Y} \operatorname{Ad}_{e^{-t Y}}(X) \\
& =e^{t Y} e^{-t \operatorname{ad}_{Y}}(X)
\end{aligned}
$$

Since $[Y,[Y, X]]=-[Y,[X, Y]]=0$, we have

$$
e^{-t \mathrm{ad}_{Y}}(X)=X-t[Y, X]=X+t[X, Y]
$$

with all higher terms being zero. We may then simplify (5.7) to

$$
\frac{d A}{d t}=e^{t X} e^{t Y} e^{-\frac{t^{2}}{2}[X, Y]}(X+Y)=A(t)(X+Y)
$$

We see, then, that $A(t)$ and $B(t)$ satisfy the same differential equation, with the same initial condition $A(0)=B(0)=I$. Thus, by standard uniqueness results for (linear) ordinary differential equations, $A(t)=B(t)$ for all $t$.

Theorem 5.2. Let $H$ denote the Heisenberg group and $\mathfrak{h}$ its Lie algebra. Let $G$ be a matrix Lie group with Lie algebra $\mathfrak{g}$ and let $\phi: \mathfrak{h} \rightarrow \mathfrak{g}$ be a Lie algebra homomorphism. Then there exists a unique Lie group homomorphism $\Phi: H \rightarrow G$ such that

$$
\Phi\left(e^{X}\right)=e^{\phi(X)}
$$

for all $X \in \mathfrak{h}$.\\
Proof. Recall (Exercise 18 in Chapter 3) that the Heisenberg group has the special property that its exponential map is one-to-one and onto. Let "log" denote the inverse of this map and define $\Phi: H \rightarrow G$ by the formula

$$
\Phi(A)=e^{\phi(\log A)}
$$

We will show that $\Phi$ is a group homomorphism.\\
If $X$ and $Y$ are in the Lie algebra of the Heisenberg group ( $3 \times 3$ strictly upper triangular matrices), direct computation shows that every entry of [ $X, Y$ ] is zero except possibly for the entry in the upper right corner. It is then easily seen that $[X, Y]$ commutes with both $X$ and $Y$. Since $\phi$ is a Lie algebra homomorphism, $\phi(X)$ and $\phi(Y)$ will also commute with their commutator. Thus, by Theorem 5.1, for any $X$ and $Y$ in the Lie algebra of the Heisenberg group, we have

$$
\begin{aligned}
\Phi\left(e^{X} e^{Y}\right) & =\Phi\left(e^{X+Y+\frac{1}{2}[X, Y]}\right) \\
& =e^{\phi(X)+\phi(Y)+\frac{1}{2}[\phi(X), \phi(Y)]} \\
& =e^{\phi(X)} e^{\phi(Y)} \\
& =\Phi\left(e^{X}\right) \Phi\left(e^{Y}\right)
\end{aligned}
$$

Thus, $\Phi$ is a group homomorphism, which is continuous because each of exp, log, and $\phi$ is continuous.

\subsection*{The Baker-Campbell-Hausdorff Formula}
The goal of the Baker-Campbell-Hausdorff formula (BCH formula) is to compute $\log \left(e^{X} e^{Y}\right)$. One may well ask, "Why do we not simply expand both exponentials and the logarithm in power series and multiply everything out?" While it is possible to do this, what is not clear is why the answer is expressible in terms of commutators. Consider, for example, the quadratic term in the expression for $\log \left(e^{X} e^{Y}\right)$, which will be a linear combination of $X^{2}, Y^{2}, X Y$, and $Y X$. For this term to be expressible in terms of commutators, it must be a multiple of $X Y-Y X$. Although direct computation verifies that this is, indeed, the case, it is far from obvious how to prove that a similar result occurs for all the higher terms.

We will actually state and prove an integral form of the BCH formula, rather than the series form (5.3). The integral version of the formula, along with the argument we present in Sect. 5.5, is actually due to Poincaré. (See [Poin1, Poin2] and Section 1.1.2.2 of [BF].)

Consider the function

$$
g(z)=\frac{\log z}{1-\frac{1}{z}}
$$

which is defined and holomorphic in the disk $\{|z-1|<1\}$. Thus, $g(z)$ can be expressed as a series

$$
g(z)=\sum_{m=0}^{\infty} a_{m}(z-1)^{m}
$$

with radius of convergence one. If $V$ is a finite-dimensional vector space, we may identify $V$ with $\mathbb{C}^{n}$ by means of an arbitrary basis, so that the Hilbert-Schmidt norm (Definition 2.2) of a linear operator on $V$ can be defined. For any operator $A$ on $V$ with $\|A-I\|<1$, we can define

$$
g(A)=\sum_{m=0}^{\infty} a_{m}(A-I)^{m}
$$

We are now ready to state the integral form of the BCH formula.\\
Theorem 5.3 (Baker-Campbell-Hausdorff). For all $n \times n$ complex matrices $X$ and $Y$ with $\|X\|$ and $\|Y\|$ sufficiently small, we have

$$
\log \left(e^{X} e^{Y}\right)=X+\int_{0}^{1} g\left(e^{\mathrm{ad}_{X}} e^{t \mathrm{ad}_{Y}}\right)(Y) d t
$$

The proof of this theorem is given in Sect. 5.5 of this chapter. Note that $e^{\operatorname{ad}_{X}} e^{t \operatorname{ad}_{Y}}$ and, hence, also $g\left(e^{\mathrm{ad}_{X}} e^{t^{\mathrm{ad}_{Y}}}\right)$ are linear operators on the space $M_{n}(\mathbb{C})$ of all $n \times n$ complex matrices. In (5.8), this operator is being applied to the matrix $Y$. The fact that $X$ and $Y$ are assumed small guarantees that $e^{\mathrm{ad}_{X}} e^{t \mathrm{ad}_{Y}}$ is close to the identity operator on $M_{n}(\mathbb{C})$ for $0 \leq t \leq 1$, so that $g\left(e^{\mathrm{ad}_{X}} e^{t \mathrm{ad}_{Y}}\right)$ is well defined. Although the right-hand side of (5.8) is rather complicated to compute explicitly, we are not so much interested in the details of the formula but in the fact that it expresses $\log \left(e^{X} e^{Y}\right)$ (and hence $e^{X} e^{Y}$ ) in terms of the Lie-algebraic quantities $\operatorname{ad}_{X}$ and ad $_{Y}$.

\subsection*{The Derivative of the Exponential Map}
In this section we prove a result that is useful in its own right and will play a key role in our proof of the BCH formula. Consider the directional derivative of exp at a point $X$ in the direction of $Y$ :

$$
\left.\frac{d}{d t} e^{X+t Y}\right|_{t=0}
$$

Unless $X$ and $Y$ commute, this derivative may not equal $e^{X+t Y} Y$. Nevertheless, since exp is continuously differentiable (Proposition 2.16), the directional derivatives in (5.9) will depend linearly on $Y$ with $X$ fixed. Now, the function $\left(1-e^{-z}\right) / z$ is an entire function of $z$ (with a removable singularity at the origin) and is given by the power series

$$
\frac{1-e^{-z}}{z}=\sum_{k=0}^{\infty}(-1)^{k} \frac{z^{k}}{(k+1)!}
$$

which has infinite radius of convergence. It thus makes sense to replace $z$ by an arbitrary linear operator $A$ on a finite-dimensional vector space.

Theorem 5.4 (Derivative of Exponential). For all $X, Y \in M_{n}(\mathbb{C})$, we have

$$
\begin{aligned}
\left.\frac{d}{d t} e^{X+t Y}\right|_{t=0} & =e^{X}\left\{\frac{I-e^{-\mathrm{ad}_{X}}}{\operatorname{ad}_{X}}(Y)\right\} \\
& =e^{X}\left\{Y-\frac{[X, Y]}{2!}+\frac{[X,[X, Y]]}{3!}-\cdots\right\}
\end{aligned}
$$

More generally, if $X(t)$ is a smooth matrix-valued function, then

$$
\frac{d}{d t} e^{X(t)}=e^{X(t)}\left\{\frac{I-e^{-\mathrm{ad}_{X(t)}}}{\operatorname{ad}_{X(t)}}\left(\frac{d X}{d t}\right)\right\}
$$

Our proof follows [Tuy].\\
Lemma 5.5. If $Z$ is a linear operator on a finite-dimensional vector space, then

$$
\lim _{m \rightarrow \infty} \frac{1}{m} \sum_{k=0}^{m-1}\left(e^{-Z / m}\right)^{k}=\frac{1-e^{-Z}}{Z}
$$

Proof. If we formally applied the formula for the sum of a finite geometric series to $e^{-Z / m}$, we would get

$$
\frac{1}{m} \sum_{k=0}^{m-1}\left(e^{-Z / m}\right)^{k}=\frac{1}{m} \frac{1-e^{-Z}}{1-e^{-Z / m}} \xrightarrow{m \rightarrow \infty} \frac{1-e^{-Z}}{Z}
$$

To give a rigorous argument, we observe that

$$
\frac{1-e^{-x}}{x}=\int_{0}^{1} e^{-t x} d t
$$

from which it follows that

$$
\frac{1-e^{-Z}}{Z}=\int_{0}^{1} e^{-t Z} d t
$$

(The reader may check, using term-by-term integration of the series expansion of $e^{-t Z}$, that this formula for $\left(1-e^{-Z}\right) / Z$ agrees with our earlier definition.)

Since $\left(e^{-Z / m}\right)^{k}=e^{-k Z / m}$, the left-hand side of (5.12) is a Riemann-sum approximation to the matrix-valued integral on the right-hand side of (5.13). These Riemann sums converge to the integral of $e^{-t Z}$-which is a continuous function of $t$-establishing (5.12).

Proof of Theorem 5.4. The formula (5.11) follows from (5.10) by applying the chain rule to the composition of exp and $X(t)$. Thus, it suffices to prove (5.10). For any $n \times n$ matrices $X$ and $Y$, set

$$
\Delta(X, Y)=\left.\frac{d}{d t} e^{X+t Y}\right|_{t=0}
$$

Since (Proposition 2.16) exp is a continuously differentiable map, $\Delta(X, Y)$ is jointly continuous in $X$ and $Y$ and is linear in $Y$ for each fixed $X$.

Now, for every positive integer $m$, we have

$$
e^{X+t Y}=\left[\exp \left(\frac{X}{m}+t \frac{Y}{m}\right)\right]^{m}
$$

Applying the product rule, we will get $m$ terms, where in each term, $m-1$ of the factors in (5.14) are simply evaluated at $t=0$ and the remaining factor is differentiated at $t=0$. Thus,

$$
\begin{aligned}
\left.\frac{d}{d t} e^{X+t Y}\right|_{t=0} & =\sum_{k=0}^{m-1}\left(e^{X / m}\right)^{m-k-1}\left[\left.\frac{d}{d t} \exp \left(\frac{X}{m}+t \frac{Y}{m}\right)\right|_{t=0}\right]\left(e^{X / m}\right)^{k} \\
& =e^{(m-1) X / m} \sum_{k=0}^{m-1}\left(e^{X / m}\right)^{-k} \Delta\left(\frac{X}{m}, \frac{Y}{m}\right)\left(e^{X / m}\right)^{k} \\
& =e^{(m-1) X / m} \frac{1}{m} \sum_{k=0}^{m-1} \exp \left(-\frac{\mathrm{ad}_{X}}{m}\right)^{k}\left(\Delta\left(\frac{X}{m}, Y\right)\right)
\end{aligned}
$$

In the third equality, we have used the linearity of $\Delta(X, Y)$ in $Y$ and the relationship between Ad and ad (Proposition 3.35).

We now wish to let $m$ tend to infinity in (5.15). The factor in front tends to $\exp (X)$. Since $\Delta(X, Y)$ is jointly continuous in $X$ and $Y$, the expression $\Delta(X / m, Y)$ tends to $\Delta(0, Y)$, where it is easily verified that $\Delta(0, Y)=Y$. Finally, applying Lemma 5.5 with $Z=\operatorname{ad}_{X}$, we see that

$$
\lim _{m \rightarrow \infty} \frac{1}{m} \sum_{k=0}^{m-1} \exp \left(-\frac{\mathrm{ad}_{X}}{m}\right)^{k}=\frac{1-e^{-\mathrm{ad}_{X}}}{\mathrm{ad}_{X}}
$$

Thus, by letting $m$ tend to infinity in (5.15), we obtain the desired result.

\subsection*{Proof of the BCH Formula}
We now turn to the proof of Theorem 5.3. For sufficiently small $X$ and $Y$ in $M_{n}(\mathbb{C})$, let

$$
Z(t)=\log \left(e^{X} e^{t Y}\right)
$$

for $0 \leq t \leq 1$. Our goal is to compute $Z(1)$. Since $e^{Z(t)}=e^{X} e^{t Y}$, we have

$$
e^{-Z(t)} \frac{d}{d t} e^{Z(t)}=\left(e^{X} e^{t Y}\right)^{-1} e^{X} e^{t Y} Y=Y
$$

On the other hand, by Theorem 5.4,

$$
e^{-Z(t)} \frac{d}{d t} e^{Z(t)}=\left\{\frac{I-e^{-\mathrm{ad}_{Z(t)}}}{\operatorname{ad}_{Z(t)}}\right\}\left(\frac{d Z}{d t}\right)
$$

Hence,

$$
\left\{\frac{I-e^{-\mathrm{ad}_{Z(t)}}}{\operatorname{ad}_{Z(t)}}\right\}\left(\frac{d Z}{d t}\right)=Y .
$$

Now, if $X$ and $Y$ are small enough, $Z(t)$ will also be small, so that $[I- \left.e^{-\mathrm{ad}_{Z(t)}}\right] / \mathrm{ad}_{Z(t)}$ will be close to the identity and thus invertible. In that case, we obtain

$$
\frac{d Z}{d t}=\left\{\frac{I-e^{-\mathrm{ad}_{Z(t)}}}{\operatorname{ad}_{Z(t)}}\right\}^{-1}(Y)
$$

Meanwhile, if we apply the homomorphism "Ad" to the equation $e^{Z(t)}=e^{X} e^{t Y}$, use the relationship between "Ad" and "ad," and take a logarithm, we obtain the following relations:

$$
\begin{aligned}
\operatorname{Ad}_{e^{Z(t)}} & =\operatorname{Ad}_{e^{X}} \operatorname{Ad}_{e^{t Y}} \\
e^{\operatorname{ad}_{Z(t)}} & =e^{\operatorname{ad}_{X}} e^{t \operatorname{ad}_{Y}} \\
\operatorname{ad}_{Z(t)} & =\log \left(e^{\operatorname{ad}_{X}} e^{t \operatorname{ad}_{Y}}\right)
\end{aligned}
$$

Plugging the last two of these relations into (5.16) gives

$$
\frac{d Z}{d t}=\left\{\frac{I-\left(e^{\mathrm{ad}_{X}} e^{t \mathrm{ad}_{Y}}\right)^{-1}}{\log \left(e^{\mathrm{ad}_{X}} e^{t \mathrm{ad}_{Y}}\right)}\right\}^{-1}(Y)
$$

Now, observe that

$$
g(z)=\left\{\frac{1-z^{-1}}{\log z}\right\}^{-1}
$$

so that (5.17) is the same as

$$
\frac{d Z}{d t}=g\left(e^{\mathrm{ad}_{X}} e^{t \mathrm{ad}_{Y}}\right)(Y) .
$$

Noting that $Z(0)=X$ and integrating (5.18) gives

$$
\log \left(e^{X} e^{Y}\right)=Z(1)=X+\int_{0}^{1} g\left(e^{\mathrm{ad}_{X}} e^{t \mathrm{ad}_{Y}}\right)(Y) d t
$$

which is the Baker-Campbell-Hausdorff formula.

\subsection*{The Series Form of the BCH Formula}
Let us see how to get the first few terms of the series form of the BCH formula from the integral form in Theorem 5.3. Using the Taylor series (2.7) for the logarithm, we may easily compute that

$$
g(z)=1+\frac{1}{2}(z-1)-\frac{1}{6}(z-1)^{2}+\frac{1}{12}(z-1)^{3}-\cdots
$$

Meanwhile,

$$
\begin{aligned}
& e^{\mathrm{ad}_{X}} e^{t \mathrm{ad}_{Y}}-I \\
& \quad=\left(I+\mathrm{ad}_{X}+\frac{\left(\mathrm{ad}_{X}\right)^{2}}{2}+\cdots\right)\left(I+t \mathrm{ad}_{Y}+\frac{t^{2}\left(\mathrm{ad}_{Y}\right)^{2}}{2}+\cdots\right)-I \\
& \quad=\mathrm{ad}_{X}+t \mathrm{ad}_{Y}+t \mathrm{ad}_{X} \mathrm{ad}_{Y}+\frac{\left(\mathrm{ad}_{X}\right)^{2}}{2}+\frac{t^{2}\left(\mathrm{ad}_{Y}\right)^{2}}{2}+\cdots
\end{aligned}
$$

Since $e^{\operatorname{ad}_{X}} e^{t \operatorname{ad}_{Y}}-I$ has no zeroth-order term, $\left(e^{\operatorname{ad}_{X}} e^{t \operatorname{ad}_{Y}}-I\right)^{m}$ will contribute only terms of degree $m$ or higher in $\mathrm{ad}_{X}$ and/or $\operatorname{ad}_{Y}$. Computing up to degree 2 in $\mathrm{ad}_{X}$ and $\mathrm{ad}_{Y}$ gives

$$
\begin{aligned}
& g\left(e^{\mathrm{ad}_{X}} e^{t \mathrm{ad}_{Y}}\right) \\
& \quad=I+\frac{1}{2}\left(\operatorname{ad}_{X}+t \operatorname{ad}_{Y}+t \operatorname{ad}_{X} \operatorname{ad}_{Y}+\frac{\left(\operatorname{ad}_{X}\right)^{2}}{2}+\frac{t^{2}\left(\operatorname{ad}_{Y}\right)^{2}}{2}\right)
\end{aligned}
$$

$$
\begin{aligned}
& -\frac{1}{6}\left[\left(\operatorname{ad}_{X}\right)^{2}+t^{2}\left(\operatorname{ad}_{Y}\right)^{2}+t \operatorname{ad}_{X} \operatorname{ad}_{Y}+t \operatorname{ad}_{Y} \operatorname{ad}_{X}\right] \\
& + \text { higher-order terms. }
\end{aligned}
$$

We now apply $g\left(e^{\operatorname{ad}_{X}} e^{t \operatorname{ad}_{Y}}\right)$ to $Y$ and integrate. Computing to second order and noting that any term with $\operatorname{ad}_{Y}$ acting first is zero, we obtain:

$$
\begin{aligned}
& \log \left(e^{X} e^{Y}\right) \\
& \approx X+\int_{0}^{1}\left[Y+\frac{1}{2}[X, Y]+\frac{1}{4}[X,[X, Y]]-\frac{1}{6}[X,[X, Y]]-\frac{t}{6}[Y,[X, Y]]\right] d t \\
& \approx X+Y+\frac{1}{2}[X, Y]+\frac{1}{12}[X,[X, Y]]-\frac{1}{12}[Y,[X, Y]]
\end{aligned}
$$

which is the expression in (5.3).

\subsection*{Group Versus Lie Algebra Homomorphisms}
Recall Theorem 3.28, which says that given matrix Lie groups $G$ and $H$ and a Lie group homomorphism $\Phi: G \rightarrow H$, we can find a Lie algebra homomorphism $\phi: \mathfrak{g} \rightarrow \mathfrak{h}$ such that $\Phi\left(e^{X}\right)=e^{\phi(X)}$ for all $X \in \mathfrak{g}$. In this section, we prove a converse to this result in the case that $G$ is simply connected.

Theorem 5.6. Let $G$ and $H$ be matrix Lie groups with Lie algebras $\mathfrak{g}$ and $\mathfrak{h}$, respectively, and let $\phi: \mathfrak{g} \rightarrow \mathfrak{h}$ be a Lie algebra homomorphism. If $G$ is simply connected, there exists a unique Lie group homomorphism $\Phi: G \rightarrow H$ such that $\Phi\left(e^{X}\right)=e^{\phi(X)}$ for all $X \in \mathfrak{g}$.

This result has the following corollary.\\
Corollary 5.7. Suppose $G$ and $H$ are simply connected matrix Lie groups with Lie algebras $\mathfrak{g}$ and $\mathfrak{h}$, respectively. If $\mathfrak{g}$ is isomorphic to $\mathfrak{h}$, then $G$ is isomorphic to $H$.

Proof. Let $\phi: \mathfrak{g} \rightarrow \mathfrak{h}$ be a Lie algebra isomorphism. By Theorem 5.6, there exists an associated Lie group homomorphism $\Phi: G \rightarrow H$. Since $\psi:=\phi^{-1}$ is also a Lie algebra homomorphism, there is a corresponding Lie group homomorphism $\Psi: H \rightarrow G$. We want to show that $\Phi$ and $\Psi$ are inverses of each other.

Now, the Lie algebra homomorphism associated to $\Phi \circ \Psi$ is, by Proposition 3.30, equal to $\phi \circ \psi=I$, and similarly for $\Psi \circ \Phi$. Thus, by Corollary 3.49, both $\Phi \circ \Psi$ and $\Psi \circ \Phi$ are equal to the identity maps on $H$ and $G$, respectively.

We now proceed with the proof of Theorem 5.6. The first step is to construct a "local homomorphism" from $\phi$. This step is the only place in the argument in which we use the BCH formula.

Definition 5.8. If $G$ and $H$ are matrix Lie groups, a local homomorphism of $G$ to $H$ is pair ( $U, f$ ) where $U$ is a path-connected neighborhood of the identity in $G$ and $f: U \rightarrow H$ is a continuous map such that $f(A B)=f(A) f(B)$ whenever $A$, $B$, and $A B$ all belong to $U$.

The definition says that $f$ is as much of a homomorphism as it makes sense to be, given that $U$ is not necessarily a subgroup of $G$.

Proposition 5.9. Let $G$ and $H$ be matrix Lie groups with Lie algebras $\mathfrak{g}$ and $\mathfrak{h}$, respectively, and let $\phi: \mathfrak{g} \rightarrow \mathfrak{h}$ be a Lie algebra homomorphism. Define $U_{\varepsilon} \subset G$ by

$$
U_{\varepsilon}=\{A \in G \mid\|A-I\|<1 \text { and }\|\log A\|<\varepsilon\}
$$

Then there exists some $\varepsilon>0$ such that the map $f: U_{\varepsilon} \rightarrow H$ given by

$$
f(A)=e^{\phi(\log A)}
$$

is a local homomorphism.\\
Note that by Theorem 3.42, if $\varepsilon$ is small enough, $\log A$ will be in $\mathfrak{g}$ for all $A \in U_{\varepsilon}$, so that $\Phi$ makes sense.

Proof. Choose $\varepsilon$ small enough that Theorem 3.42 applies and small enough that for all $A, B \in U_{\varepsilon}$, the BCH formula applies to $X:=\log A$ and $Y:=\log B$ and also to $\phi(X)$ and $\phi(Y)$. Then if $A B$ happens to be in $U_{\varepsilon}$, we have

$$
f(A B)=f\left(e^{X} e^{Y}\right)=e^{\phi\left(\log \left(e^{X} e^{Y}\right)\right)}
$$

We now compute $\log \left(e^{X} e^{Y}\right)$ by the BCH formula and then apply $\phi$. Since $\phi$ is a Lie algebra homomorphism, it will change all the Lie-algebraic quantities involving $X$ and $Y$ in the BCH formula into the analogous quantities involving $\phi(X)$ and $\phi(Y)$. Thus, as in (5.4), we have

$$
\begin{aligned}
\phi\left[\log \left(e^{X} e^{Y}\right)\right] & =\phi(X)+\int_{0}^{1} \sum_{m=0}^{\infty} a_{m}\left(e^{\operatorname{ad}_{\phi(X)}} e^{t \operatorname{ad}_{\phi(Y)}}-I\right)^{m}(\phi(Y)) d t \\
& =\log \left(e^{\phi(X)} e^{\phi(Y)}\right)
\end{aligned}
$$

We obtain, then,

$$
\begin{aligned}
f(A B) & =\exp \left\{\log \left(e^{\phi(X)} e^{\phi(Y)}\right)\right\} \\
& =e^{\phi(X)} e^{\phi(Y)} \\
& =f(A) f(B),
\end{aligned}
$$

as claimed.

Theorem 5.10. Let $G$ and $H$ be matrix Lie groups, with $G$ simply connected. If ( $U, f$ ) is a local homomorphism of $G$ into $H$, there exists a unique (global) Lie group homomorphism $\Phi: G \rightarrow H$ such that $\Phi$ agrees with $f$ on $U$.

Proof. Step 1: Define $\Phi$ along a path. Since $G$ is simply connected and thus connected, for any $A \in G$, there exists a path $A(t) \in G$ with $A(0)=I$ and $A(1)=A$. Let us call a partition $0=t_{0}<t_{1}<t_{2} \cdots<t_{m}=1$ of [ 0,1 ] a good partition if for all $s$ and $t$ belonging to the same subinterval of the partition, we have

$$
A(t) A(s)^{-1} \in U
$$

Lemma 3.48 guarantees that good partitions exist. If a partition is good, then, in particular, since $t_{0}=0$ and $A(0)=I$, we have $A\left(t_{1}\right) \in U$. Choose a good partition and write $A$ as

$$
A=\left[A(1) A\left(t_{m-1}\right)^{-1}\right]\left[A\left(t_{m-1}\right) A\left(t_{m-2}\right)^{-1}\right] \cdots\left[A\left(t_{2}\right) A\left(t_{1}\right)^{-1}\right] A\left(t_{1}\right) .
$$

Since $\Phi$ is supposed to be a homomorphism and is supposed to agree with $f$ near the identity it is reasonable to "define" $\Phi(A)$ by

$$
\Phi(A)=f\left(A(1) A\left(t_{m-1}\right)^{-1}\right) \cdots f\left(A\left(t_{2}\right) A\left(t_{1}\right)^{-1}\right) f\left(A\left(t_{1}\right)\right) .
$$

In the next two steps, we will prove that $\Phi(A)$ is independent of the choice of partition for a fixed path and independent of the choice of path.\\
Step 2: Prove independence of the partition. For any good partition, if we insert an extra partition point $s$ between $t_{j}$ and $t_{j+1}$, the result is easily seen to be another good partition. This change in the partition has the effect of replacing the factor $f\left(A\left(t_{j+1}\right) A\left(t_{j}\right)^{-1}\right)$ in (5.20) by

$$
f\left(A\left(t_{j+1}\right) A(s)^{-1}\right) f\left(A(s) A\left(t_{j}\right)^{-1}\right) .
$$

Since $s$ is between $t_{j}$ and $t_{j+1}$, the condition (5.19) on the original partition guarantees that $A\left(t_{j+1}\right) A(s)^{-1}, A(s) A\left(t_{j}\right)^{-1}$ and $A\left(t_{j+1}\right) A\left(t_{j}\right)^{-1}$ are all in $U$. Thus, since $f$ is a local homomorphism, we have

$$
f\left(A\left(t_{j+1}\right) A\left(t_{j}\right)^{-1}\right)=f\left(A\left(t_{j+1}\right) A(s)^{-1}\right) f\left(A(s) A\left(t_{j}\right)^{-1}\right)
$$

showing that the value of $\Phi(A)$ is unchanged by the addition of the extra partition point.\\
By repeating this argument, we see that the value of $\Phi(A)$ does not change by the addition of any finite number of points to the partition. Now, any two good partitions have a common refinement, namely their union, which is also a good partition. The above argument shows that the value of $\Phi(A)$ computed from the first partition is the same as for the common refinement, which is the same as for the second partition.

Step 3: Prove independence of the path. It is in this step that we use the simple connectedness of $G$. Suppose $A_{0}(t)$ and $A_{1}(t)$ are two paths joining the identity to some $A \in G$. Then, since $G$ is simply connected, a standard topological argument (e.g., Proposition 1.6 in [Hat]) shows that $A_{0}$ and $A_{1}$ are homotopic with endpoints fixed. This means that there exists a continuous map $A:[0,1] \times[0,1] \rightarrow G$ with

$$
A(0, t)=A_{0}(t), \quad A(1, t)=A_{1}(t)
$$

for all $t \in[0,1]$ and also

$$
A(s, 0)=I, \quad A(s, 1)=A
$$

for all $s \in[0,1]$.\\
As in the proof of Lemma 3.48, there exists an integer $N$ such that for all $(s, t)$ and ( $s^{\prime}, t^{\prime}$ ) in $[0,1] \times[0,1]$ with $\left|s-s^{\prime}\right|<2 / N$ and $\left|t-t^{\prime}\right|<2 / N$, we have

$$
A(s, t) A\left(s^{\prime}, t^{\prime}\right)^{-1} \in U .
$$

We now deform $A_{0}$ "a little bit at a time" into $A_{1}$. This means that we define a sequence $B_{k, l}$ of paths, with $k=0, \ldots, N-1$ and $l=0, \ldots, N$. We define these paths so that $B_{k, l}(t)$ coincides with $A((k+1) / N, t)$ for $t$ between 0 and $(l-1) / N$, and $B_{k, l}(t)$ coincides with $A(k / N, t)$ for $t$ between $l / N$ and 1 . For $t$ between $(l-1) / N$ and $l / N$, we define $B_{k, l}(t)$ to coincide with the values of $A(\cdot, \cdot)$ on the path that goes "diagonally" in the ( $s, t$ )-plane, as indicated in Figure 5.1. When $l=0$, there are no $t$-values between 0 and $(l-1) / N$, so $B_{k, 0}(t)=A(k / N, t)$ for all $t \in[0,1]$. In particular, $B_{0,0}(t)=A_{0}(t)$.


Fig. 5.1 The path in the $(s, t)$ plane defining $B_{k, l}(t)$

\includegraphics[scale=0.2, center]{2025_08_28_fc86fd4ade69f41b5842g-134}


We now deform $A_{0}=B_{0,0}$ into $B_{0,1}$ and then into $B_{0,2}, B_{0,3}$, and so on until we reach $B_{0, N}$, which we then deform into $B_{1,0}$ and so on until we reach $B_{N-1, N}$, which we finally deform into $A_{1}$. We claim that the value of $\Phi(A)$ is the same at each stage of this deformation. Note that for $k<l, B_{k, l}(t)$ and $B_{k, l+1}(t)$ are the same except for $t$ 's in the interval

$$
[(l-1) / N,(l+1) / N] .
$$

Furthermore, by Step 2, we are free to choose any good partition we like to compute $\Phi(A)$. For both $B_{k, l}$ and $B_{k, l+1}$, we choose the partition points to be

$$
0, \frac{1}{N}, \ldots, \frac{l-1}{N}, \frac{l+1}{N}, \frac{l+2}{N}, \ldots, 1,
$$

which gives a good partition by the way $N$ was chosen.\\
Now, from (5.20), the value of $\Phi(A)$ depends only on the values of the path at the partition points. Since we have chosen our partition in such a way that the values of $B_{k, l}$ and $B_{k, l+1}$ are identical at all the partition points, the value of $\Phi(A)$ is the same for these two paths. (See Figure 5.2.) A similar argument shows that the value of $\Phi(A)$ computed along $B_{k, N}$ is the same as along $B_{k+1,0}$. Thus, the value of $\Phi(A)$ is the same for each path from $A_{0}=B_{0,0}$ all the way to $B_{N-1, N}$ and then the same as $A_{1}$.


Fig. 5.2 The paths $B_{k, l}$ and $B_{k, l+1}$ agree at each partition point. In the figure, $s$ increases as we move from the top toward the bottom


\includegraphics[scale=0.2, center]{2025_08_28_fc86fd4ade69f41b5842g-135}


Step 4: Prove that $\Phi$ is a homomorphism and agrees with $f$ on $U$. The proof that $\Phi$ is a homomorphism is a straightforward unpacking of the definition of $\Phi$ and is left to the reader; see Exercise 7. To show that $\Phi$ agrees with $f$ on $U$, choose $A \in U$. Since $U$ is path connected, we can find a path $A(t), 0 \leq t \leq 1$, lying in $U$ joining $I$ to $A$. Choose a good partition $\left\{t_{j}\right\}_{j=1}^{m}$ for $A(t)$, and we then claim that for all $j$, we have $\Phi\left(A\left(t_{j}\right)\right)=f\left(A\left(t_{j}\right)\right)$.\\
Note that $A(t), 0 \leq t \leq t_{j}$ is a path joining $I$ to $A\left(t_{j}\right)$ and that $\left\{t_{0}, t_{1}, \ldots, t_{j}\right\}$ is a good partition of this path. (Technically, we should reparameterize this path so that the time interval is $[0,1]$.) Hence,

$$
\Phi\left(A\left(t_{j}\right)\right)=f\left(A\left(t_{j}\right) A\left(t_{j-1}\right)^{-1}\right) \cdots f\left(A\left(t_{2}\right) A\left(t_{1}\right)^{-1}\right) f\left(A\left(t_{1}\right)\right),
$$

for all $j$. In particular,

$$
\Phi\left(A\left(t_{1}\right)\right)=f\left(A\left(t_{1}\right)\right)
$$

Now assume that $\Phi\left(A\left(t_{j}\right)\right)=f\left(A\left(t_{j}\right)\right)$, and compute that

$$
\begin{aligned}
\Phi\left(A\left(t_{j+1}\right)\right) & =f\left(A\left(t_{j+1}\right) A\left(t_{j}\right)^{-1}\right) f\left(A\left(t_{j}\right) A\left(t_{j-1}\right)^{-1}\right) \cdots f\left(A\left(t_{1}\right)\right) \\
& =f\left(A\left(t_{j+1}\right) A\left(t_{j}\right)^{-1}\right) \Phi\left(A\left(t_{j}\right)\right) \\
& =f\left(A\left(t_{j+1}\right) A\left(t_{j}\right)^{-1}\right) f\left(A\left(t_{j}\right)\right) \\
& =f\left(A\left(t_{j}\right)\right)
\end{aligned}
$$

The last equality holds because $f$ is a local homomorphism and because $A\left(t_{j+1}\right) A\left(t_{j}\right)^{-1}, A\left(t_{j}\right)$, and their product all lie in $U$. Thus, by induction, $\Phi\left(A\left(t_{j}\right)\right)=f\left(A\left(t_{j}\right)\right)$ for all $j$; when $j=m$, we obtain $\Phi(A)=f(A)$.

It is important to note that when proving independence of the path (Step 3 of the proof), it is essential to know already that the value of $\Phi(A)$ is independent of the choice of good partition (Step 2 of the proof). Specifically, when we move from $B_{k, l}$ to $B_{k, l+1}$, we use one partition for $B_{k, l+1}$, but when we move from $B_{k, l+1}$ to $B_{k, l+2}$, we use a different partition for $B_{k, l+1}$. Essentially, the proof proceeds by deforming the path between partition points-which clearly does not change the value of $\Phi(A)$-then picking a new partition and doing the same thing again. Note also that it is in proving independence of the partition that we use the assumption that $f$ is a local homomorphism.

Proof of Theorem 5.6. For the existence part of the proof, let $f$ be the local homomorphism in Proposition 5.9 and let $\Phi$ be the global homomorphism in Theorem 5.10. Then for any $X \in \mathfrak{g}$, the element $e^{X / m}$ will be in $U$ for all sufficiently large $m$, showing that

$$
\Phi\left(e^{X / m}\right)=f\left(e^{X / m}\right)=e^{\phi(X) / m}
$$

Since $\Phi$ is a homomorphism, we have

$$
\Phi\left(e^{X}\right)=\Phi\left(e^{X / m}\right)^{m}=e^{\phi(X)},
$$

as required.\\
For the uniqueness part of the proof, suppose $\Phi_{1}$ and $\Phi_{2}$ are two homomorphisms related in the desired way to $\phi$. Then for any $A \in G$, we express $A$ as $e^{X_{1}} \cdots e^{X_{N}}$ with $X_{j} \in \mathfrak{g}$, as in Corollary 3.47, and observe that

$$
\Phi_{1}(A)=\Phi_{2}(A)=e^{\phi\left(X_{1}\right)} \cdots e^{\phi\left(X_{N}\right)}
$$

so that $\Phi_{1}$ agrees with $\Phi_{2}$.\\
We conclude this section with a typical application of Theorem 5.6.\\
Theorem 5.11. Suppose that $G$ is a simply connected matrix Lie group and that the Lie algebra $\mathfrak{g}$ of $G$ decomposes as a Lie algebra direct sum $\mathfrak{g}=\mathfrak{h}_{1} \oplus \mathfrak{h}_{2}$, for two subalgebras $\mathfrak{h}_{1}$ and $\mathfrak{h}_{2}$ of $\mathfrak{g}$. Then there exists closed, simply connected subgroups $H_{1}$ and $H_{2}$ of $G$ whose Lie algebras are $\mathfrak{h}_{1}$ and $\mathfrak{h}_{2}$, respectively. Furthermore, $G$ is isomorphic to the direct product of $H_{1}$ and $H_{2}$.

Proof. Consider the Lie algebra homomorphism $\phi: \mathfrak{g} \rightarrow \mathfrak{g}$ that sends $X+Y$ to $X$, where $X \in \mathfrak{h}_{1}$ and $Y \in \mathfrak{h}_{2}$. Since $G$ is simply connected, there is a Lie group homomorphism $\Phi: G \rightarrow G$ associated to $\phi$, and the Lie algebra of the kernel of $\Phi$ is the kernel of $\phi$ (Proposition 3.31), which is $\mathfrak{h}_{2}$. Let $H_{2}$ be the identity component of $\operatorname{ker} \Phi$. Since $\Phi$ is continuous, $\operatorname{ker} \Phi$ is closed, and so is its identity component $H_{2}$ (Corollary 3.52). Thus, $H_{2}$ is a closed, connected subgroup of $G$ with Lie algebra $\mathfrak{h}_{2}$. By a similar argument, we may construct a closed, connected subgroup $H_{1}$ of $G$ whose Lie algebra is $\mathfrak{h}_{1}$.

Suppose now that $A(t)$ is a loop in $H_{1}$. Since $G$ is simply connected, there is a homotopy $A(s, t)$ shrinking $A(t)$ to a point in $G$. Now, $\phi$ is the identity on $\mathfrak{h}_{1}$, from which it follows that $\Phi$ is the identity on $H_{1}$. Thus, if we define $B(s, t)= \Phi(A(s, t))$, we see that

$$
B(0, t)=\Phi(A(t))=A(t)
$$

Furthermore, since $\phi$ maps $G$ into $\mathfrak{h}_{1}$, we see that $\Phi$ maps $G$ into $H_{1}$. We conclude that $B$ is a homotopy of $A(t)$ to a point in $H_{1}$. Thus, $H_{1}$ is simply connected, and, by a similar argument, so is $\mathrm{H}_{2}$.

Finally, since $\mathfrak{g}$ is the Lie algebra direct sum of $\mathfrak{h}_{1}$ and $\mathfrak{h}_{2}$, elements of $\mathfrak{h}_{1}$ commutes with elements of $\mathfrak{h}_{2}$. It follows that elements of $H_{1}$ (which are all product of exponentials of elements of $\mathfrak{h}_{1}$ ) commute with elements of $H_{2}$. Thus, we have a Lie group homomorphism $\Psi: H_{1} \times H_{2} \rightarrow G$ given by $\Psi(A, B)=A B$. The associated Lie algebra homomorphism $\psi$ is then just the original isomorphism of $\mathfrak{h}_{1} \oplus \mathfrak{h}_{2}$ with $\mathfrak{g}$. Since $G$ is simply connected, there is a homomorphism $\Gamma: G \rightarrow H_{1} \times H_{2}$ for which the associated Lie algebra homomorphism is $\psi^{-1}$. By the proof of Corollary 5.7, $\Gamma$ and $\Psi$ are inverses of each other, showing that $G$ is isomorphic to $H_{1} \times H_{2}$.

\subsection*{Universal Covers}
Theorem 5.6 says that if $G$ is simply connected, every homomorphism of the Lie algebra $\mathfrak{g}$ of $G$ can be exponentiated to a homomorphism of $G$. If $G$ is not simply connected, we may look for another group $\tilde{G}$ that has the same Lie algebra as $G$ but such that $\tilde{G}$ is simply connected.

Definition 5.12. Let $G$ be a connected matrix Lie group. Then a universal cover of $G$ is a simply connected matrix Lie group $H$ together with a Lie group homomorphism $\Phi: H \rightarrow G$ such that the associated Lie algebra homomorphism $\phi: \mathfrak{h} \rightarrow \mathfrak{g}$ is a Lie algebra isomorphism. The homomorphism $\Phi$ is called the covering map.

If a universal cover of $G$ exists, it is unique up to "canonical isomorphism," as follows.

Proposition 5.13. If $G$ is a connected matrix Lie group and ( $H_{1}, \Phi_{1}$ ) and ( $H_{2}, \Phi_{2}$ ) are universal covers of $G$, then there exists a Lie group isomorphism $\Psi: H_{1} \rightarrow H_{2}$ such that $\Phi_{2} \circ \Psi=\Phi_{1}$.

Proof. See Exercise 9.\\
Since a connected matrix Lie group has at most one universal cover (up to canonical isomorphism), it is reasonable to speak of the universal cover ( $\tilde{G}, \Phi$ ) of $G$. Furthermore, if $H$ is a simply connected Lie group and $\phi: \mathfrak{h} \rightarrow \mathfrak{g}$ is a Lie algebra isomorphism, then by Theorem 5.6, we can construct an associated Lie group homomorphism $\Phi: H \rightarrow G$, so that ( $H, \Phi$ ) is a universal cover of $G$. Since $\phi$ is an isomorphism, we can use $\phi$ to identify $\tilde{\mathfrak{g}}$ with $\mathfrak{g}$. Thus, in slightly less formal terms, we may define the notion of universal cover as follows: The universal cover of a matrix Lie group $G$ is a simply connected matrix Lie group $\tilde{G}$ such that the Lie algebra of $\tilde{G}$ is equal to the Lie algebra of $G$. With this perspective, we have the following immediate corollary of Theorem 5.6.\\
Corollary 5.14. Let $G$ be a connected matrix Lie group and let $\tilde{G}$ be the universal cover of $G$, where we think of $G$ and $\tilde{G}$ as having the same Lie algebra $\mathfrak{g}$. If $H$ is a matrix Lie group with Lie algebra $\mathfrak{h}$ and $\phi: \mathfrak{g} \rightarrow \mathfrak{h}$ is a Lie algebra homomorphism, there exists a unique homomorphism $\Phi: \tilde{G} \rightarrow H$ such that $\Phi\left(e^{X}\right)=e^{\phi(X)}$ for all $X \in \mathfrak{g}$.

An example of importance in physics is the universal cover of $\mathrm{SO}(3)$.\\
Example 5.15. The universal cover of $\mathrm{SO}(3)$ is $\mathrm{SU}(2)$.\\
Proof. The group SU(2) is simply connected by Proposition 1.15. Proposition 1.19 and Example 3.29 then provide the desired covering map.

The topic of universal covers is one place where we pay a price for our decision to consider only matrix Lie groups: a matrix Lie group may not have a universal cover that is a matrix Lie group. It is not hard to show that every Lie group $G$\\
has a universal cover in the class of (not necessarily matrix) Lie groups. Indeed, $G$ has a universal cover in the topological sense (Definition 13.1), and this cover can be given a group structure in such a way that the covering map is a Lie group homomorphism. It turns out, however, that the universal cover of a matrix Lie group may not be a matrix group.

We now show that the group $\operatorname{SL}(2 ; \mathbb{R})$ does not have a universal cover in the class of matrix Lie groups. We begin by showing that $\operatorname{SL}(2 ; \mathbb{R})$ is not simply connected. By Theorem 2.17 and Proposition 2.19, $\operatorname{SL}(2 ; \mathbb{R})$, as a manifold, is homeomorphic to $\mathrm{SO}(2) \times V$, where $V$ is the space of $2 \times 2$ real, symmetric matrices with trace zero. Now, $V$, being a vector space, is certainly simply connected, but $\mathrm{SO}(2)$, which is homeomorphic to the unit circle $S^{1}$, is not. Thus, $\operatorname{SL}(2 ; \mathbb{R})$ itself is not simply connected. By contrast, the group $\mathrm{SL}(2 ; \mathbb{C})$ decomposes as $\mathrm{SU}(2) \times W$, where $W$ is the space of $2 \times 2$ self-adjoint, complex matrices with trace zero. Since both $\mathrm{SU}(2)$ and $W$ are simply connected, $\mathrm{SL}(2 ; \mathbb{C})$ is simply connected. (See Appendix 13.3 for more information on this type of calculation.)

Proposition 5.16. Let $G \subset \mathrm{GL}(n ; \mathbb{C})$ be a connected matrix Lie group with Lie algebra $\mathfrak{g}$. Suppose $\Phi: G \rightarrow \mathrm{SL}(2 ; \mathbb{R})$ is a Lie group homomorphism for which the associated Lie algebra $\phi: \mathfrak{g} \rightarrow \mathrm{sl}(2 ; \mathbb{R})$ is a Lie algebra isomorphism. Then $\Phi$ is a Lie group isomorphism and, therefore, $G$ cannot be simply connected. Thus, $\mathrm{SL}(2 ; \mathbb{R})$ has no universal cover in the class of matrix Lie groups.

The result relies essentially on the assumption that $G$ is a matrix Lie group-or, more precisely, on the assumption that $G$ is contained in a group whose Lie algebra is complex.

Lemma 5.17. Suppose $\psi: \operatorname{sl}(2 ; \mathbb{R}) \rightarrow \mathrm{gl}(n ; \mathbb{C})$ is a Lie algebra homomorphism. Then there exists a Lie group homomorphism $\Phi: \mathrm{SL}(2 ; \mathbb{R}) \rightarrow \mathrm{GL}(n ; \mathbb{C})$ such that $\Phi\left(e^{X}\right)=e^{\phi(X)}$ for all $X \in \operatorname{sl}(2 ; \mathbb{R})$.

The significance of the lemma is that the result holds even though $\operatorname{SL}(2 ; \mathbb{R})$ is not simply connected.

Proof. Let $\psi_{\mathbb{C}}: \operatorname{sl}(2 ; \mathbb{C}) \rightarrow \mathrm{gl}(n ; \mathbb{C})$ be the complex-linear extension of $\psi$ to $\mathrm{sI}(2 ; \mathbb{C}) \cong \mathrm{sI}(2 ; \mathbb{R})_{\mathbb{C}}$, which is a Lie algebra homomorphism (Proposition 3.39). Since $\operatorname{SL}(2 ; \mathbb{C})$ is simply connected, there exists a Lie group homomorphism $\Psi_{\mathbb{C}}: \mathrm{SL}(2 ; \mathbb{C}) \rightarrow \mathrm{GL}(n ; \mathbb{C})$ such that $\Psi_{\mathbb{C}}\left(e^{X}\right)=e^{\psi_{\mathbb{C}}(X)}$ for all $X \in \mathrm{SI}(2 ; \mathbb{C})$. If we let $\Psi$ be the restriction of $\Psi_{\mathbb{C}}$ to $\operatorname{SL}(2 ; \mathbb{R})$, then $\Psi$ is a Lie group homomorphism which satisfies $\Psi\left(e^{X}\right)=e^{\psi(X)}$ for $X \in \mathrm{Sl}(2 ; \mathbb{R})$.

Proof of Proposition 5.16. Since $\phi$ is a Lie algebra isomorphism, the inverse map $\psi: \mathrm{sI}(2 ; \mathbb{R}) \rightarrow \mathfrak{g}$ is a Lie algebra homomorphism. Thus, by the lemma, there is a Lie group homomorphism $\Psi: \operatorname{SL}(2 ; \mathbb{R}) \rightarrow G$ corresponding to $\psi$. Since $\phi$ and $\psi$ are inverses of each other, it follows from Proposition 3.30 and Corollary 3.49 that $\Phi$ and $\Psi$ are also inverses of each other.

\subsection*{Subgroups and Subalgebras}
In this section, we address Question 3 from Sect. 5.1: If $G$ is a matrix Lie group with Lie algebra $\mathfrak{g}$ and $\mathfrak{h}$ is a subalgebra of $\mathfrak{g}$, does there exist a matrix Lie group $H \subset G$ whose Lie algebra is $H$ ? If the exponential map for $G$ were a homeomorphism between $\mathfrak{g}$ and $G$ and if the BCH formula worked globally instead of locally, the answer would be yes, since we could simply define $H$ to be the set of elements of the form $e^{X}, X \in \mathfrak{g}$, and the BCH formula would show that $H$ is a subgroup.

In reality, the answer to Question 3, as stated, is no. Suppose, for example, that $G=\mathrm{GL}(2 ; \mathbb{C})$ and

$$
\mathfrak{h}=\left\{\left.\left(\begin{array}{cc}
\text { it } & 0 \\
0 & \text { ita }
\end{array}\right) \right\rvert\, t \in \mathbb{R}\right\},
$$

where $a$ is irrational. If there is going to be a matrix Lie group $H$ with Lie algebra $\mathfrak{h}$, then $H$ would have to contain the closure of the group

$$
H_{0}=\left\{\left.\left(\begin{array}{cc}
e^{i t} & 0 \\
0 & e^{i t a}
\end{array}\right) \right\rvert\, t \in \mathbb{R}\right\},
$$

which is (Exercise 10 in Chapter 1) is the group

$$
H_{1}=\left\{\left.\left(\begin{array}{cc}
e^{i \theta} & 0 \\
0 & e^{i \phi}
\end{array}\right) \right\rvert\, \theta, \phi \in \mathbb{R}\right\}
$$

But then the Lie algebra of $H$ would have to contain the Lie algebra of $H_{1}$, which is two dimensional!

Fortunately, all is not lost. We can still get a subgroup $H$ for each subalgebra $\mathfrak{h}$ if we weaken the condition that $H$ be a matrix Lie group. In the above example, the subgroup we want is $H_{0}$, despite the fact that $H_{0}$ is not closed.

Definition 5.18. If $H$ is any subgroup of $\mathrm{GL}(n ; \mathbb{C})$, the Lie algebra $\mathfrak{h}$ of $H$ is the set of all matrices $X$ such that

$$
e^{t X} \in H
$$

for all real $t$.\\
It is possible to prove that for any subgroup $H$ of $\mathrm{GL}(n ; \mathbb{C})$, the Lie algebra $\mathfrak{h}$ of $H$ is actually a Lie algebra, that is, a real vector space-possibly zero dimensionaland closed under brackets. (See Proposition 1 and Corollary 7 in Chapter 2 of [Ross].) This result is not, however, directly relevant to our goal in this section, which is to construct, for each subalgebra $\mathfrak{h}$ of $\mathrm{gl}(n ; \mathbb{C})$ a subgroup with Lie algebra $\mathfrak{h}$. Note, however, that if $\mathfrak{h}$ is at least a real subspace of $\operatorname{gl}(n ; \mathbb{C})$, then the proof of Point 4 of Theorem 3.20 shows that $\mathfrak{h}$ is also closed under brackets.

Definition 5.19. If $G$ is a matrix Lie group with Lie algebra $\mathfrak{g}$, then $H \subset G$ is a connected Lie subgroup of $G$ if the following conditions are satisfied:

\begin{enumerate}
  \item $H$ is a subgroup of $G$.
  \item The Lie algebra $\mathfrak{h}$ of $H$ is a Lie subalgebra of $\mathfrak{g}$.
  \item Every element of $H$ can be written in the form $e^{X_{1}} e^{X_{2}} \cdots e^{X_{m}}$, with $X_{1}, \ldots, X_{m} \in \mathfrak{h}$.
\end{enumerate}

Connected Lie subgroups are also called analytic subgroups. Note that any group $H$ as in the definition is path connected, since each element of $H$ can be connected to the identity in $H$ by a path of the form

$$
t \mapsto e^{(1-t) X_{1}} e^{(1-t) X_{2}} \cdots e^{(1-t) X_{m}}
$$

The group $H_{0}$ in (5.22) is a connected Lie subgroup of $\mathrm{GL}(2 ; \mathbb{C})$ whose Lie algebra is the algebra $\mathfrak{h}$ in (5.21).

We are now ready to state the main result of this section, which is our second major application of the Baker-Campbell-Hausdorff formula.

Theorem 5.20. Let $G$ be a matrix Lie group with Lie algebra $\mathfrak{g}$ and let $\mathfrak{h}$ be a Lie subalgebra of $\mathfrak{g}$. Then there exists a unique connected Lie subgroup $H$ of $G$ with Lie algebra $\mathfrak{h}$.

If $\mathfrak{h}$ is the subalgebra of $\mathrm{gl}(2 ; \mathbb{C})$ in (5.21), then the connected Lie subgroup $H$ is the group $H_{0}$ in (5.22), which is not closed. In practice, Theorem 5.20 is most useful in those cases where the connected Lie subgroup $H$ turns out to be closed. See Proposition 5.24 and Exercises 10, 13, and 14 for conditions under which this is the case.

We now begin working toward the proof of Theorem 5.20. Since $G$ is assumed to be a matrix Lie group, we may as well assume that $G=\mathrm{GL}(n ; \mathbb{C})$. After all, if $G$ is a closed subgroup of $\mathrm{GL}(n ; \mathbb{C})$ and $H$ is a connected Lie subgroup of $\mathrm{GL}(n ; \mathbb{C})$ whose Lie algebra $\mathfrak{h}$ is contained in $\mathfrak{g}$, then $H$ is also a connected Lie subgroup of $G$. We now let

$$
H=\left\{e^{X_{1}} e^{X_{2}} \cdots e^{X_{N}} \mid X_{1}, \ldots, X_{N} \in \mathfrak{h}\right\}
$$

which is a subgroup of $G$. The key issue is to prove that the Lie algebra of $H$, in the sense of Definition 5.18, is $\mathfrak{h}$. Once we know that $\operatorname{Lie}(H)=\mathfrak{h}$, we will immediately conclude that $H$ is a connected Lie subgroup with Lie algebra $\mathfrak{h}$, the remaining properties in Definition 5.19 being true by definition. Note that for the claim $\operatorname{Lie}(H)=\mathfrak{h}$ to be true, it essential that $\mathfrak{h}$ be a subalgebra of $\operatorname{gl}(n ; \mathbb{C})$, and not merely a subspace; compare Exercise 11.

As in the proof of Theorem 3.42, we think of $\operatorname{gl}(n ; \mathbb{C})$ as $\mathbb{R}^{2 n^{2}}$ and we decompose $\operatorname{gl}(n ; \mathbb{C})$ as the direct sum of $\mathfrak{h}$ and $D$, where $D$ is the orthogonal complement of $\mathfrak{h}$ with respect to the usual inner product on $\mathbb{R}^{2 n^{2}}$. Then, as shown in the proof of Theorem 3.42, there exist neighborhoods $U$ and $V$ of the origin in $\mathfrak{h}$ and $D$, respectively, and a neighborhood $W$ of $I$ in $\mathrm{GL}(n ; \mathbb{C})$ with the following properties:

Each $A \in W$ can be written uniquely as

$$
A=e^{X} e^{Y}, \quad X \in U, Y \in V,
$$

in such a way that $X$ and $Y$ depend continuously on $A$. We think of the decomposition in (5.24) as our local coordinates in a neighborhood of the identity in $\mathrm{GL}(n ; \mathbb{C})$.

If $X$ is a small element of $\mathfrak{h}$, the decomposition of $e^{X}$ is just $e^{X} e^{0}$. If we take the product of two elements of the form $e^{X_{1}} e^{X_{2}}$, with $X_{1}$ and $X_{2}$ small elements of $\mathfrak{h}$, then since $\mathfrak{h}$ is a subalgebra, if we combine the exponentials as $e^{X_{1}} e^{X_{2}}=e^{X_{3}}$ by means of the Baker-Campbell-Hausdorff formula, $X_{3}$ will again be in $\mathfrak{h}$. Thus, if we take a small number of products as in (5.23) with the $X_{j}$ 's being small elements of $\mathfrak{h}$, we will move from the identity in the $X$-direction in the decomposition (5.24). Globally, however, $H$ may wind around and come back to points in $W$ of the form (5.24) with $Y \neq 0$. (See Figure 5.3.) Indeed, as the example of the "irrational line" in (5.22) shows, there may be elements of $H$ in $W$ with arbitrarily small nonzero values of $Y$. Nevertheless, we will see that the set of $Y$ values that occurs is at most countable.

Lemma 5.21. Decompose $\operatorname{gl}(n ; \mathbb{C})$ as $\mathfrak{h} \oplus D$ and let $V$ be a neighborhood of the origin in $D$ as in (5.24). If $E \subset V$ is defined by

$$
E=\left\{Y \in V \mid e^{Y} \in H\right\},
$$

then $E$ is at most countable.\\
Assuming the lemma, we may now prove Theorem 5.20.


\includegraphics[scale=0.2, center]{2025_08_28_fc86fd4ade69f41b5842g-142}

Fig. 5.3 The black lines indicate the portion of $H$ in the set $W$. The group $H$ intersects $e^{V}$ in at most countably many points


Proof of Theorem 5.20. As we have already observed, it suffices to show that the Lie algebra of $H$ is $\mathfrak{h}$. Let $\mathfrak{h}^{\prime}$ be the Lie algebra of $H$, which clearly contains $\mathfrak{h}$. For $Z \in \mathfrak{h}^{\prime}$, we may write, for all sufficiently small $t$,

$$
e^{t Z}=e^{X(t)} e^{Y(t)}
$$

where $X(t) \in U \subset \mathfrak{h}$ and $Y(t) \in V \subset D$ and where $X(t)$ and $Y(t)$ are continuous functions of $t$. Since $Z$ is in the Lie algebra of $H$, we have $e^{t Z} \in H$ for all $t$. Since, also, $e^{X(t)}$ is in the group $H$, we conclude that $e^{Y(t)}$ is in $H$ for all sufficiently small $t$. If $Y(t)$ were not constant, then it would take on uncountably many values, which would mean that $E$ is uncountable, violating Lemma 5.21. So, $Y(t)$ must be constant, and since $Y(0)=0$, this means that $Y(t)$ is identically equal to zero. Thus, for small $t$, we have $e^{t Z}=e^{X(t)}$ and, therefore, $t Z=X(t) \in \mathfrak{h}$. This means $Z \in \mathfrak{h}$ and we conclude that $\mathfrak{h}^{\prime} \subset \mathfrak{h}$.

Before proving Lemma 5.21, we prove another lemma.\\
Lemma 5.22. Pick a basis for $\mathfrak{h}$ and call an element of $\mathfrak{h}$ rational if its coefficients with respect to this basis are rational. Then for every $\delta>0$ and every $A \in H$, there exist rational elements $R_{1}, \ldots, R_{m}$ of $\mathfrak{h}$ such that

$$
A=e^{R_{1}} e^{R_{2}} \cdots e^{R_{m}} e^{X}
$$

where $X$ is in $\mathfrak{h}$ and $\|X\|<\delta$.\\
Suppose we take $\delta$ small enough that the ball of radius $\delta$ in $\mathfrak{h}$ is contained in $U$. Then since there are only countably many $m$-tuples of the form ( $R_{1}, \ldots, R_{m}$ ) with $R_{j}$ rational, the lemma tells us that $H$ can be covered by countably many translates of the set $e^{U}$.

Proof. Choose $\varepsilon>0$ so that for all $X, Y \in \mathfrak{h}$ with $\|X\|<\varepsilon$ and $\|Y\|<\varepsilon$, the Baker-Campbell-Hausdorff holds for $X$ and $Y$. Let $C(\cdot, \cdot)$ denote the right-hand side of the formula, so that

$$
e^{X} e^{Y}=e^{C(X, Y)}
$$

whenever $\|X\|,\|Y\|<\varepsilon$. It is not hard to see that $C(\cdot, \cdot)$ is a continuous. Now, if the lemma holds for some $\delta$, it also holds for any $\delta^{\prime}>\delta$. Thus, it is harmless to assume $\delta$ is less than $\varepsilon$ and small enough that if $\|X\|,\|Y\|<\delta$, we have $\|C(X, Y)\|<\varepsilon$.

Since $e^{X}=\left(e^{X / k}\right)^{k}$, every element $A$ of $H$ can be written as

$$
A=e^{X_{1}} \cdots e^{X_{N}}
$$

with $X_{j} \in \mathfrak{h}$ and $\left\|X_{j}\right\|<\delta$. We now proceed by induction on $N$. If $N=0$, then $A=I=e^{0}$, and there is nothing to prove. Assume the lemma for $A$ 's that can be expressed as in (5.25) for some integer $N$, and consider $A$ of the form

$$
A=e^{X_{1}} \cdots e^{X_{N}} e^{X_{N+1}}
$$

with $X_{j} \in \mathfrak{h}$ and $\left\|X_{j}\right\|<\delta$. Applying our induction hypothesis to $e^{X_{1}} \cdots e^{X_{N}}$, we obtain

$$
\begin{aligned}
A & =e^{R_{1}} \cdots e^{R_{m}} e^{X} e^{X_{N+1}} \\
& =e^{R_{1}} \cdots e^{R_{m}} e^{C\left(X, X_{N+1}\right)} .
\end{aligned}
$$

where the $R_{j}$ 's are rational and $\left\|C\left(X, X_{N+1}\right)\right\|<\varepsilon$. Since $\mathfrak{h}$ is a subalgebra of $\mathrm{gl}(n ; \mathbb{C})$, the element $C\left(X, X_{N+1}\right)$ is again in $\mathfrak{h}$, but may not have norm less than $\delta$.

Now choose a rational element $R_{m+1}$ of $\mathfrak{h}$ that is very close to $C\left(X, X_{N+1}\right)$ and such that $\left\|R_{m+1}\right\|<\varepsilon$. We then have

$$
\begin{aligned}
A & =e^{R_{1}} \cdots e^{R_{m}} e^{R_{m+1}} e^{-R_{m+1}} e^{C\left(X, X_{N+1}\right)} \\
& =e^{R_{1}} \cdots e^{R_{m}} e^{R_{m+1}} e^{X^{\prime}}
\end{aligned}
$$

where

$$
X^{\prime}=C\left(-R_{m+1}, C\left(X, X_{N+1}\right)\right)
$$

Then $X^{\prime}$ will be in $\mathfrak{h}$, and by choosing $R_{m+1}$ sufficiently close to $C\left(X, X_{N+1}\right)$, we can make $\left\|X^{\prime}\right\|<\delta$. After all, since $C(-Z, Z)=\log \left(e^{-Z} e^{Z}\right)=0$ for all small $Z$, if $Z^{\prime}$ is close to $Z$, then $C\left(-Z^{\prime}, Z\right)$ will be small.

We now supply the proof of Lemma 5.21.\\
Proof of Lemma 5.21. Fix $\delta$ so that for all $X$ and $Y$ with $\|X\|,\|Y\|<\delta$, the quantity $C(X, Y)$ is defined and contained in $U$. We then claim that for each sequence $R_{1}, \ldots, R_{m}$ of rational elements in $\mathfrak{h}$, there is at most one $X \in \mathfrak{h}$ with $\|X\|<\delta$ such that the element

$$
e^{R_{1}} e^{R_{2}} \cdots e^{R_{m}} e^{X}
$$

belongs to $e^{V}$. After all, if we have

$$
\begin{aligned}
& e^{R_{1}} e^{R_{2}} \cdots e^{R_{m}} e^{X_{1}}=e^{Y_{1}} \\
& e^{R_{1}} e^{R_{2}} \cdots e^{R_{m}} e^{X_{2}}=e^{Y_{2}}
\end{aligned}
$$

with $Y_{1}, Y_{2} \in V$, then

$$
e^{-X_{1}} e^{X_{2}}=e^{-Y_{1}} e^{Y_{2}}
$$

and so

$$
e^{-Y_{1}}=e^{-X_{1}} e^{X_{2}} e^{-Y_{2}}=e^{C\left(-X_{1}, X_{2}\right)} e^{-Y_{2}}
$$

with $C\left(-X_{1}, X_{2}\right) \in U$. However, each element of $e^{U} e^{V}$ has a unique representation as $e^{Y} e^{X}$ with $X \in U$ and $Y \in V$. Thus, we must have $-Y_{2}=-Y_{1}$ and, by (5.28) and (5.29), $e^{X_{1}}=e^{X_{2}}$ and $X_{1}=X_{2}$.

By Lemma 5.22, every element of $H$ can be expressed in the form (5.27) with $\|X\|<\delta$. Now, there are only countably many rational elements in $\mathfrak{h}$ and thus only countably many expressions of the form $e^{R_{1}} \cdots e^{R_{m}}$, each of which produces at most one element of the form (5.27) that belongs to $e^{V}$. Thus, the set $E$ in Lemma 5.21 is at most countable.

This completes the proof of Theorem 5.20.\\
If a connected Lie subgroup $H$ of $\mathrm{GL}(n ; \mathbb{C})$ is not closed, the topology $H$ inherits from $\mathrm{GL}(n ; \mathbb{C})$ may be pathological, e.g., not locally connected. (Compare Figure 1.1.) Nevertheless, we can give $H$ a new topology that is much nicer.

Theorem 5.23. Let $H$ be a connected Lie subgroup of $\mathrm{GL}(n ; \mathbb{C})$ with Lie algebra $\mathfrak{h}$. Then $H$ can be given the structure of a smooth manifold in such a way that the group operations on $H$ are smooth and the inclusion map of $H$ into $\mathrm{GL}(n ; \mathbb{C})$ is smooth.

Thus, every connected Lie subgroup of $\mathrm{GL}(n ; \mathbb{C})$ can be made into a Lie group. In the case of the group $H_{0}$ in (5.22), the new topology on $H_{0}$ is obtained by identifying $H_{0}$ with $\mathbb{R}$ by means of the parameter $t$ in the definition of $H_{0}$.

Proof. For any $A \in H$ and any $\varepsilon>0$, define

$$
U_{A, \varepsilon}=\left\{A e^{X} \mid X \in \mathfrak{h} \text { and }\|X\|<\varepsilon\right\} .
$$

Now define a topology on $H$ as follows: A set $U \subset H$ is open if for each $A \in U$ there exists $\varepsilon>0$ such that $U_{A, \varepsilon} \subset U$. (See Figure 5.4.) In this topology, two elements $A$ and $B$ of $H$ are "close" if we can express $B$ as $B=A e^{X}$ with $X \in \mathfrak{h}$ and $\|X\|$ small. This topology is finer than the topology $H$ inherits from $G$; that is, if $A$ and $B$ are close in this new topology, they are certainly close in the ordinary sense in $G$, but not vice versa.

It is easy to check that this topology is Hausdorff, and using Lemma 5.22, it is not hard to see that the topology is second countable. Furthermore, in this topology, $H$ is locally homeomorphic to $\mathbb{R}^{N}$, where $N=\operatorname{dimh}$, by identifying each $U_{A, \varepsilon}$ with the ball of radius $\varepsilon$ in $\mathfrak{h}$.

We may define a smooth structure on $H$ by using the $U_{A, \varepsilon}$ 's, with $\varepsilon$ less than some small number $\varepsilon_{0}$, as our "atlas." If two of these sets overlap, then some element $C$ of $H$ can be written as $C=A e^{X}=B e^{Y}$ for some $A, B \in H$ and $X, Y \in \mathfrak{h}$. It follows that $B=A e^{X} e^{-Y}$, which means (since $\|X\|$ and $\|Y\|$ are less than $\varepsilon_{0}$ ) that $A$ and $B$ are close. The change-of-coordinates map is then $Y=\log \left(B^{-1} A e^{X}\right)$. Since $A$ and $B$ are close and $\|X\|$ is small, we will have that $\left\|B^{-1} A e^{X}-I\right\|<1$, so that $B^{-1} A e^{X}$ is in the domain where the matrix logarithm is defined and smooth. Thus, the change-of-coordinates map is smooth as function of $X$. Finally, in any of the coordinate neighborhoods $U_{A, \varepsilon}$, the inclusion of $H$ into $G$ is given by $X \mapsto A e^{X}$, which is smooth as a function of $X$.


Fig. 5.4 The set $U$ in $H$ is open the new topology but not in the topology inherited from $\mathrm{GL}(2 ; \mathbb{C})$. The element $B$ is close to $A$ in $\mathrm{GL}(2 ; \mathbb{C})$ but not in the new topology on $H$

\includegraphics[scale=0.2, center]{2025_08_28_fc86fd4ade69f41b5842g-146}


As we have already noted, Theorem 5.20 is most useful in cases where the connected Lie subgroup $H$ is actually closed. The following result gives one condition under which this is guaranteed to be the case. See also Exercises 10, 13, and 14.

Proposition 5.24. Suppose $G \subset \mathrm{GL}(n ; \mathbb{C})$ is a matrix Lie group with Lie algebra $\mathfrak{g}$ and that $\mathfrak{h}$ is a maximal commutative subalgebra of $\mathfrak{g}$, meaning that $\mathfrak{h}$ is commutative and $\mathfrak{h}$ is not contained in any larger commutative subalgebra of $\mathfrak{g}$. Then the connected Lie subgroup $H$ of $G$ with Lie algebra $\mathfrak{h}$ is closed.

Proof. Since $\mathfrak{h}$ is commutative, $H$ is also commutative, since every element of $H$ is a product of exponentials of elements of $\mathfrak{h}$. It easily follows that the closure $\bar{H}$ of $H$ in $\mathrm{GL}(n ; \mathbb{C})$ is also commutative. We now claim that $\bar{H}$ is connected. To see this, take $A \in \bar{H}$, so that $A$ is in $G$ (since $G$ is closed) and $A$ is the limit of a sequence $A_{m}$ in $H$. Since $\bar{H}$ is closed, Theorem 3.42 applies. Thus, for all sufficiently large $m$, the element $A A_{m}^{-1}$ is expressible as $A A_{m}^{-1}=e^{X}$, for some $X$ in the Lie algebra $\mathfrak{h}^{\prime}$ of $\bar{H}$. Thus, $A=e^{X} A_{m}$, which means that $A$ can be connected to $A_{m}$ by the path $A(t)=e^{(1-t) X} A_{m}, 0 \leq t \leq 1$, in $\bar{H}$. Since $A_{m}$ can be connected to the identity in $H \subset \bar{H}$, we see that $A$ can be connected to the identity in $\bar{H}$.

Now, since $\bar{H}$ is commutative, its Lie algebra $\mathfrak{h}^{\prime}$ is also commutative. But since $\mathfrak{h}$ was maximal commutative, we must have $\mathfrak{h}^{\prime}=\mathfrak{h}$. Since, also, $\bar{H}$ is connected, we conclude that $\bar{H}=H$, showing that $H$ is closed.

\subsection*{Lie's Third Theorem}
Lie's third theorem (in its modern, global form) says that for every finitedimensional, real Lie algebra $\mathfrak{g}$, there exists a Lie group $G$ with Lie algebra $\mathfrak{g}$. We will construct $G$ as a connected Lie subgroup of $\operatorname{GL}(n ; \mathbb{C})$.

Theorem 5.25. If $\mathfrak{g}$ is any finite-dimensional, real Lie algebra, there exists a connected Lie subgroup $G$ of $\mathrm{GL}(n ; \mathbb{C})$ whose Lie algebra is isomorphic to $\mathfrak{g}$.

Our proof assumes Ado's theorem, which asserts that every finite-dimensional real or complex Lie algebra is isomorphic to an algebra of matrices. (See, for example, Theorem 3.17.7 in [Var].)

Proof. By Ado's theorem, we may identify $\mathfrak{g}$ with a real subalgebra of $\mathfrak{g l}(n ; \mathbb{C})$. Then by Theorem 5.20, there is a connected Lie subgroup of $\operatorname{GL}(n ; \mathbb{C})$ with Lie algebra $\mathfrak{g}$.

It is actually possible to choose the subgroup $G$ in Theorem 5.25 to be closed. Indeed, according to Theorem 9 on p. 105 of [Got], if a connected Lie group $G$ can be embedded into some $\mathrm{GL}(n ; \mathbb{C})$ as a connected Lie subgroup, then $G$ can be embedded into some other $\mathrm{GL}\left(n^{\prime} ; \mathbb{C}\right)$ as a closed subgroup. Assuming this result, we may reach the following conclusion.

Conclusion 5.26. Every finite-dimensional, real Lie algebra is isomorphic to the Lie algebra of some matrix Lie group.

This result does not, however, mean that every Lie group is isomorphic to a matrix Lie group, since there can be several nonisomorphic Lie groups with the same Lie algebra. See, for example, Sect. 4.8.

\subsection*{Exercises}
\begin{enumerate}
  \item Let $X$ be a linear transformation on a finite-dimensional real or complex vector space. Show that
\end{enumerate}

$$
\frac{I-e^{-X}}{X}
$$

is invertible if and only if none of the eigenvalues of $X$ (over $\mathbb{C}$ ) is of the form $2 \pi i n$, with $n$ an nonzero integer.

Remark. This exercise, combined with the formula in Theorem 5.4, gives the following result (in the language of differentiable manifolds): The exponential map exp : $\mathfrak{g} \rightarrow G$ is a local diffeomorphism near $X \in \mathfrak{g}$ if and only if $\operatorname{ad}_{X}$ : $\mathfrak{g} \rightarrow \mathfrak{g}$ has no eigenvalue of the form $2 \pi i n$, with $n$ a nonzero integer.\\
2. Show that for any $X$ and $Y$ in $M_{n}(\mathbb{C})$, even if $X$ and $Y$ do not commute,

$$
\left.\frac{d}{d t} \operatorname{trace}\left(e^{X+t Y}\right)\right|_{t=0}=\operatorname{trace}\left(e^{X} Y\right)
$$

\begin{enumerate}
  \setcounter{enumi}{2}
  \item Compute $\log \left(e^{X} e^{Y}\right)$ through third order in $X$ and $Y$ by calculating directly with the power series for the exponential and the logarithm. Show this gives the same answer as the Baker-Campbell-Hausdorff formula.
  \item Suppose that $X$ and $Y$ are upper triangular matrices with zeros on the diagonal. Show that the power series for $\log \left(e^{X} e^{Y}\right)$ is convergent. What happens to the series form of the Baker-Campbell-Hausdorff formula in this case?
  \item Suppose $X$ and $Y$ are $n \times n$ complex matrices satisfying $[X, Y]=\alpha Y$ for some complex number $\alpha$. Suppose further that there is no nonzero integer $n$ such that $\alpha=2 \pi i$ in. Show that
\end{enumerate}

$$
e^{X} e^{Y}=\exp \left\{X+\frac{\alpha}{1-e^{-\alpha}} Y\right\} .
$$

Hint: Let $A(t)=e^{X} e^{t Y}$ and let

$$
B(t)=\exp \left\{X+\frac{\alpha}{1-e^{-\alpha}} t Y\right\}
$$

Using Theorem 5.4, show that $A(t)$ and $B(t)$ satisfy the same differential equation with the same value at $t=0$.\\
6. Give an example of matrices $X$ and $Y$ in $\mathrm{sl}(2 ; \mathbb{C})$ such that $[X, Y]=2 \pi i Y$ but such that there does not exist any $Z$ in $\mathrm{sl}(2 ; \mathbb{C})$ with $e^{X} e^{Y}=e^{Z}$. Use Example 3.41 and compare Exercise 5.\\
7. Complete Step 4 in the proof of Theorem 5.6 by showing that $\Phi$ is a homomorphism. For all $A, B \in G$, choose a path $A(t)$ connecting $I$ to $A$ and a path $B(t)$ connecting $I$ to $B$. Then define a path $C$ connecting $I$ to $A B$ by setting $C(t)=B(2 t)$ for $0 \leq t \leq 1 / 2$ and setting $C(t)=A(2 t-1) B$ for $1 / 2 \leq t \leq 1$. If $t_{0}, \ldots, t_{m}$ is a good partition for $A(t)$ and $s_{0}, \ldots, s_{M}$ is a good partition for $B(t)$, show that

$$
\frac{s_{0}}{2}, \cdots, \frac{s_{M}}{2}, \frac{1}{2}+\frac{t_{0}}{2}, \cdots, \frac{1}{2}+\frac{t_{m}}{2}
$$

is a good partition for $C(t)$. Now, compute $\Phi(A), \Phi(B)$, and $\Phi(A B)$ using these paths and partitions and show that $\Phi(A B)=\Phi(A) \Phi(B)$.\\
8. If $\tilde{G}$ is a universal cover of a connected group $G$ with projection map $\Phi$, show that $\Phi$ maps $\tilde{G}$ onto $G$.\\
9. Prove the uniqueness of the universal cover, as stated in Proposition 5.13.\\
10. Let $\mathfrak{a}$ be a subalgebra of the Lie algebra of the Heisenberg group. Show that $\exp (\mathfrak{a})$ is a connected Lie subgroup of the Heisenberg group and that this subgroup is closed.\\
11. Consider the Lie algebra $\mathfrak{h}$ of the Heisenberg group $H$, as computed in Proposition 3.26. Let $X, Y$, and $Z$ be the basis elements for $\mathfrak{h}$ in (4.18), which satisfy $[X, Y]=Z$ and $[X, Z]=[Y, Z]=0$. Let $V$ be the subspace of $\mathfrak{h}$ spanned by $X$ and $Y$ (which is not a subalgebra of $\mathfrak{h}$ ) and let $K$ denote the subgroup of $H$ consisting of products of exponential of elements of $V$. Show that $K=H$ and, thus, that the Lie algebra of $K$ is not equal to $V$.\\
Hint: Use Theorem 5.1 and the surjectivity of the exponential map for $H$ (Exercise 18 in Chapter 3).\\
12. Show that every connected Lie subgroup of $\mathrm{SU}(2)$ is closed. Show that this is not the case for $\mathrm{SU}(3)$.\\
13. Let $G$ be a matrix Lie group with Lie algebra $\mathfrak{g}$, let $\mathfrak{h}$ be a subalgebra of $\mathfrak{g}$, and let $H$ be the unique connected Lie subgroup of $G$ with Lie algebra $\mathfrak{h}$. Suppose that there exists a simply connected, compact matrix Lie group $K$ such that the Lie algebra of $K$ is isomorphic to $\mathfrak{h}$. Show that $H$ is closed. Is $H$ necessarily isomorphic to $K$ ?\\
14. This exercise asks you to prove, assuming Ado's theorem (Sect. 5.10), the following result: If $G$ is a simply connected matrix Lie group with Lie algebra $\mathfrak{g}$ and $\mathfrak{h}$ is an ideal in $\mathfrak{g}$, then the connected Lie subgroup $H$ with Lie algebra $\mathfrak{h}$ is closed.\\
(a) Show that there exists a Lie algebra homomorphism $\phi: \mathfrak{g} \rightarrow \operatorname{gl}(N ; \mathbb{C})$ with $\operatorname{ker}(\phi)=\mathfrak{h}$.\\
Hint: Since $\mathfrak{h}$ is an ideal in $\mathfrak{g}$, the quotient space $\mathfrak{g} / \mathfrak{h}$ has a natural Lie algebra structure.\\
(b) Since $G$ is simply connected, there exists a Lie group homomorphism $\Phi$ : $G \rightarrow \mathrm{gl}(N ; \mathbb{C})$ for which the associated Lie algebra homomorphism is $\phi$. Show that the identity component of the kernel of $\Phi$ is a closed subgroup of $G$ whose Lie algebra is $\mathfrak{h}$.\\
(c) Show that the result fails if the assumption that $G$ be simply connected is omitted.



\pagebreak
\printbibliography
\end{document}